% Part A.  Included below the root document's lettered section heading.
% Every numbered environment has a status in its title.
\subsection{The typed groupoid construction and a correction on scan 51}

In this part $K$ is a field of characteristic zero, $\overline K$ is a
chosen separable closure, $\Gamma=G_K$, and
$T=\varprojlim_n\mu_n(\overline K)=\widehat{\mathbf Z}(1)$.
For a smooth proper geometrically connected curve $X/K$ and a reduced
finite divisor $S$, put $U=X\setminus S$, $E=\pi_1(U)$ and
$\pi=\pi_1(U_{\overline K})$.  A choice of geometric base point is
understood; its change acts by geometric inner conjugacy.  A loop group
$L_{\widetilde s}$ is the inertia at a lift of a geometric cusp to the
completed universal covering.  Its arithmetic decomposition group is
$E_{\widetilde s}$.  We use $\operatorname{Int}(a)(b)=aba^{-1}$.
The letter $\Delta$ in the CM calculation below denotes
$G_{\mathbf Q(i)^{\rm ab}}$; it never denotes the geometric group in
that calculation.  Grothendieck's scans 43--54 use $\mathcal E,\pi,
\Gamma,\mathcal E_i,L_i,T$ for the respective objects.

\begin{proposition}[Proved: the two functors and the exact boundary diagram]
\label{A:typed-functor}
The literal functor of scan 23 sends a smooth projective anabelian
$K$-scheme of dimension at most one to
$\Pi_1(X)\longrightarrow\Pi_1(K)$.  Its pointed version has morphisms
$(f,\gamma):(X',\xi')\to(X,\xi)$, where $f$ is a $K$-morphism and
$\gamma$ is an \'etale path from $f(\xi')$ to $\xi$.  The path supplies
the actual homomorphism of based groups.  Passing to paths modulo target
geometric conjugacy recovers the unpointed groupoid formulation.

For a pair $(X,S)$ the boundary enhancement of scans 28--35 is the
following diagram of profinite groupoids, with its displayed
commutation isomorphism and its Kummer isomorphism:
\[
 \Pi_{D^*}\xrightarrow{\rho}\Pi_U\xrightarrow{\varphi}\Pi_e,
 \qquad
 \Pi_{D^*}\xrightarrow{\sigma}\Pi_D\xrightarrow{\psi}\Pi_e,
 \qquad \alpha:\varphi\rho\simeq\psi\sigma,
 \qquad \ker(\sigma)\simeq\psi^*T .
\]
Here $D_s=\operatorname{Spec}\widehat{\mathcal O}_{X,s}$,
$D_s^*=D_s\setminus\{s\}$, $D=\coprod_{s\in S}D_s$ and
$D^*=\coprod_{s\in S}D_s^*$.  The map $\sigma$ induces a bijection
on components and surjections on automorphism groups, its kernel is
abelian, $\psi$ is faithful, and $\rho$ is faithful for hyperbolic
curves.  The images in $\Pi_e$ are open.

After choosing one object in each component, the same data are
$p:E\to\Gamma$, its kernel $\pi$, a finite index set $I$, injections
$\rho_i:E_i\hookrightarrow E$, open images
$\Gamma_i=p(E_i)$, and equivariant isomorphisms
$k_i:T\xrightarrow{\sim}L_i=E_i\cap\pi$.
The image of $p$ is permitted to be open rather than all of $\Gamma$.
Disjoint unions and the separate constant fields recover the
non-geometrically-connected version.
\end{proposition}
\begin{proof}
The \'etale fibre functors on a scheme and their isomorphisms form its
profinite fundamental groupoid.  Pullback of finite \'etale covers
gives the stated covariant functor on fundamental groupoids.  In the
pointed construction the specified path identifies the target fibre
functor, which gives the asserted actual homomorphism.  Any two choices
of that path differ by an element of the target geometric group; this
is exactly the displayed conjugacy relation.

The complete local ring at $s$ is a complete discrete valuation ring
with residue field $K(s)$.  In characteristic zero the geometric
punctured local field has tame inertia $T$, by the covers
$t^{1/n}$.  Its conjugation action is the cyclotomic action.  This
proves the local kernel and equivariance statements, including the
nonrational case with quotient $G_{K(s)}$.  The inclusion of local
inertia and decomposition in $\pi_1(U)$ is injective for a hyperbolic
curve: over an algebraic closure this is the injection of a peripheral
profinite cyclic group in the surface group; the arithmetic assertion
then follows by the exact sequences and the injection of
$G_{K(s)}$ in $G_K$.  Choosing objects and paths in a connected
groupoid identifies it with the one-object groupoid of its
automorphism group.  Applying that equivalence simultaneously to the
diagram gives precisely the listed group data.  Reversing this
construction gives the original diagram up to equivalence.
\end{proof}

% Additional case analysis is placed here so that every case of section 8
% has an explicit statement, independently of the subsequent normalizer work.
\subsection{The individual cases of section 8}

\begin{proposition}[Proved using the supplied local theorem: proper source and affine target]
\label{A:proper-affine}
Let $K$ be a field finitely generated over $\mathbf Q$, or a finite
extension of $\mathbf Q_p$. Let $X'/K$ be a smooth proper
geometrically connected curve and $U/K$ an affine hyperbolic curve.
Every continuous homomorphism of extensions
\[
 f:E_{X'}\longrightarrow E_U,\qquad p_Uf=p_{X'},
\]
is trivial on $P_{X'}=\pi_1(X'_{\overline K})$. Consequently it
factors uniquely as $E_{X'}\to G_K\xrightarrow{s}E_U$, where
$s$ is a continuous section. This proves the group-theoretic assertion
of case IV(b), scan 66, manuscript folio 31. The input used to exclude
an infinite-index noncuspidal image is exactly the local section
theorem stipulated in the prompt, with its stated unrefereed status.
\end{proposition}
\begin{proof}
First let $K=k$ be a finite extension of $\mathbf Q_p$. The actual
local construction in Proposition~\ref{F-local-dichotomy} gives
three possibilities for $H=f(P_{X'})$: it is open in $P_U$, it is
trivial, or it lies nontrivially in one cusp inertia $I_c$.
This construction uses the supplied point-or-cusp theorem over
finite extensions of $k$, Carlson--Stix's point-theoretic
Hom theorem, and compactifications of arbitrarily high genus.
In particular it does not require an arithmetic map to preserve
point sections before that local theorem has been applied.

The open possibility is impossible. Mochizuki's Theorem A, also
in its formulation with an arbitrary smooth source, would give a
dominant morphism $X'\to U$. Every regular function on the
geometrically connected proper curve $X'$ is constant, so a
morphism to the affine curve $U$ is constant. This is incompatible
with dominance. Equivalently, for genus zero or one sources one
can first puncture $X'$ to a hyperbolic curve, apply Theorem A,
and use triviality on the newly introduced inertia to extend the
result across those punctures.

Suppose instead that $1\ne H\subset I_c$. Normality of $H$ in
$f(E_{X'})$ and the inertia-normalizer identity put the whole
image in $D_c$. Since it projects onto $G_k$, the cusp is
$k$-rational. The resulting geometric map therefore factors as
the following $G_k$-equivariant continuous map of abelian groups:
\[
 P_{X'}\longrightarrow P_{X'}^{\rm ab}
   =\prod_\ell T_\ell J_{X'}
 \longrightarrow I_c=\widehat{\mathbf Z}(1).
\]
For each prime $\ell$, the final coordinate vanishes. Indeed the
Weil pairing identifies
\[
 \operatorname{Hom}_{G_k}(T_\ell J_{X'},\mathbf Z_\ell(1))
       =(T_\ell J_{X'}^\vee)^{G_k}=0.
\]
The last equality follows from finiteness of the rational torsion
of an abelian variety over a $p$-adic local field: the inverse
limit, under multiplication by $\ell$, of its bounded groups
of rational $\ell$-power torsion is zero. A homomorphism from a
pro-$q$ group to a pro-$\ell$ group is zero for $q\ne\ell$;
hence these coordinate calculations kill the entire displayed
map. This contradicts $H\ne1$. Thus $H=1$, and the quotient
property of $E_{X'}\to G_k$ gives the asserted factorization.

For $K$ finitely generated over $\mathbf Q$, use the embedding
$K\hookrightarrow k$ into a finite $p$-adic extension constructed
in Lemma~\ref{A:local-embedding}. Base change gives the cartesian
maps
\[
 E_{X'_k}=E_{X'}\times_{G_K}G_k,
 \qquad E_{U_k}=E_U\times_{G_K}G_k.
\]
Their geometric groups identify with the original geometric
groups, since geometric fundamental groups of smooth curves in
characteristic zero are invariant under algebraically closed
field extension. The induced $f_k$ thus has exactly the same
geometric map as $f$. The local calculation makes that map zero.
No injectivity or openness of the homomorphism $G_k\to G_K$
is used. This establishes the stated field range by an explicit
base-change morphism, without an additional specialization claim.
\end{proof}

\begin{proposition}[Proved: cases I, II, IV(a), IV(c), IV(d), and the exact point-source remainder]
\label{A:section-eight-cases}
In the category constructed in Proposition~\ref{A:typed-functor},
the cases of scans 55--67 have the following precise reductions.
Cases I and II are finite Galois theory; case III is the point
section map; case IV(b) is Proposition~\ref{A:proper-affine};
case IV(c) has empty Hom set in the literal pair category;
and case IV(d) is the dominant relative Hom bijection of
Proposition~\ref{A:relative-hom}. For case IV(a), over a field
finitely generated over $\mathbf Q$ or a finite $p$-adic
extension, the groupoid Hom set is the disjoint union of the
dominant-morphism classes and the section classes of the target.
The section classes give the remaining constant maps exactly
when those sections are point-induced. Over finite $p$-adic
extensions and proper hyperbolic targets this last assertion
follows from the stipulated local theorem; over number fields it
retains precisely the global point-section surjectivity problem.
\end{proposition}
\begin{proof}
In case I, write $X=\operatorname{Spec}L$ and
$X'=\operatorname{Spec}L'$ with the chosen finite subextensions
of $\overline K/K$. A pair $(f,\alpha)$ in scans 55--56 has
$f=\operatorname{Int}(\alpha)|_{G_{L'}}$ and exists exactly when
$\alpha G_{L'}\alpha^{-1}\subset G_L$. The quotient by its
2-morphisms is the actual transporter quotient
\[
 G_L\backslash
 \{\alpha\in G_K:\alpha G_{L'}\alpha^{-1}\subset G_L\}
    \longrightarrow \operatorname{Hom}_K(L,L'),
 \qquad [\alpha]\longmapsto\alpha^{-1}|_L .
\]
It is bijective by finite Galois theory. The inverse assigns an
extension of the given embedding to $\overline K$; any two such
extensions differ by the displayed left $G_L$-action.

In case II, let $L'=H^0(X',\mathcal O_{X'})$ for the complete
source. Its arithmetic fundamental group projects onto $G_{L'}$.
The target projection $G_L\hookrightarrow G_K$ is injective,
so the same equation determines $f$ from $\alpha$, kills the
geometric subgroup, and gives the same transporter quotient.
The actual morphism is the composite
$X'\to\operatorname{Spec}L'\to\operatorname{Spec}L$.
The literal pair condition requires $S'=\varnothing$; otherwise
both Hom sets are empty because neither $f^{-1}(\varnothing)=S'$
nor a map $I'\to\varnothing$ exists. The relaxation recorded
at scans 59--61 forgets the source boundary for a zero-dimensional
target, and its groupoid version does exactly the same.

In case III, conjugate $(f,\alpha)$ by an element of $E_U$
above $\alpha^{-1}$. It becomes $(s,1)$ with $p_Us=1$.
The remaining conjugations are exactly those by $P_U$, giving
the section classes, as in scans 62--64. Thus no dominant Hom
theorem proves all of this case: its exact comparison morphism
is $U(K)\to\operatorname{Sec}(E_U/G_K)/P_U$. For affine
targets this map misses the tangential packets, as the explicit
Kummer counterexample in Part B proves; the corrected version
has to exclude those packets. For proper targets its global
surjectivity is the remaining question just stated.

For IV(a), the trichotomy of Proposition~\ref{F-local-dichotomy}
has no cusp alternative because the target is proper. It gives
either an open geometric image, handled by Theorem A and the
reduction in Proposition~\ref{A:relative-hom}, or the unique
factorization through $G_K$ just written. These alternatives
are disjoint, since a constant map kills the geometric group
and a dominant map has open geometric image. The same
finite-$p$-adic base change used in the preceding proposition
proves this classification for fields finitely generated over
$\mathbf Q$.

For IV(c), $I'\ne\varnothing$ and $I=\varnothing$. The map
of finite index sets required by the original category does not
exist. Its geometric condition has the identical failure
$f^{-1}(\varnothing)=\varnothing\ne S'$. The later enlargement
in section 13 is implemented by filling the source punctures
on fundamental groups; its cohomological compatibility is
spelled out in Proposition~\ref{A:kummer-ambiguity}.

Finally IV(d) has both boundary sets nonempty. Its positive
integral local degrees force an open geometric image by
Proposition~\ref{A:positive-degree-open}, with the common
cyclotomic multiplier supplied by the boundary diagram.
Theorem A now applies over every sub-$p$-adic field.
Propositions~\ref{A:relative-hom} and
\ref{A:boundary-recovery} recover the unique map of pairs,
its local degrees, its nonrational boundary points, and every
2-morphism. Tamagawa's Theorem 0.3 gives the corresponding
affine Isom assertion over fields finitely generated over
$\mathbf Q$; taking an isomorphism and its inverse in the
same construction gives precisely his formulation.
\end{proof}

\begin{remark}[Refuted in its unrestricted reading: the scope of Lead 1]
\label{A:lead-one}
The dominant part of Conjectures 5.5 and 7.1 is proved by the
explicit chain of maps in Proposition~\ref{A:relative-hom}.
The assertion that the Hom theorems settle all of Conjecture 5.5
is false for the uncorrected affine point-source instance: a
tangential section gives an actual extra groupoid morphism.
After the section-9 correction, the proper point-source
instance is still exactly the global section map of case III,
and the all-source comparison is the proved equivalence of
Proposition~\ref{F-full-curve-reduction}. Thus the relation
between the two formulations is the displayed functor and
these explicit reductions; it does not remove that surviving
point-source assertion.
\end{remark}

\begin{proposition}[Proved: degrees, natural transformations, and the missing factor]
\label{A:two-morphisms}
With the convention above, a morphism between the connected boundary
systems consists of
\[
 f:E'\to E,\quad \tau:I'\to I,\quad f_j:E'_j\to E_{\tau(j)},
 \quad a_j\in E,\quad \alpha\in\Gamma
\]
subject to
\[
 pf=\operatorname{Int}(\alpha)p',\qquad
 \rho_{\tau(j)}f_j=\operatorname{Int}(a_j)f\rho'_j,
 \qquad \beta_j=p(a_j)\alpha,
\]
and
\[
 (k_{\tau(j)})^{-1}f_jk'_j
       =d_j\chi(\beta_j)\operatorname{id}_T,
       \qquad d_j\in\mathbf Z_{>0}.
\]
A natural transformation from these data to their primed versions is
given by $\mu\in E$ and $\mu_j\in E_{\tau(j)}$ with
\[
\begin{aligned}
 f'&=\operatorname{Int}(\mu)f,&
 f'_j&=\operatorname{Int}(\mu_j)f_j,\\
 \alpha'&=p(\mu)\alpha,&
 a'_j&=\rho_{\tau(j)}(\mu_j)a_j\mu^{-1}.
\end{aligned}                                                  \tag{A1}
\]
The last factor $\mu^{-1}$ is required.  It is absent from the boxed
formula on the actual scan 51, not just from its transcription.
\end{proposition}
\begin{proof}
A natural transformation between homomorphisms of one-object
groupoids is an element conjugating the first homomorphism into the
second.  Naturality on the boundary square gives
\[
 \rho_{\tau(j)}(\mu_j)a_j=a'_j\mu,
\]
which proves (A1).  Naturality on the base square gives the formula
for $\alpha'$.  Applying $p$ to the boundary square gives
$p_{\tau(j)}f_j=\operatorname{Int}(\beta_j)p'_j$.  The action on
$T$ of the base-change path $\beta_j$ is $\chi(\beta_j)$; the
local map $t\mapsto u(t')^{d_j}$ contributes multiplication by the
positive ramification index $d_j$.  These two operations give the
displayed degree formula.  Conversely those equations are exactly
the four naturality equations and the Kummer compatibility in the
groupoid prism, so they construct a morphism of that diagram.

There is a finite certificate for the missing factor.  Take all
groups in the boundary square to be $S_3$, take $\rho=f=f_j$ to be
the identity and $a_j=1$, and take $\mu=(12)$ and $\mu_j=1$.
Then $f'=\operatorname{Int}(\mu)$ and $f'_j=1_{S_3}$.
The corrected $a'_j=\mu^{-1}$ satisfies the square; the printed
$a'_j=1$ would assert $\operatorname{Int}(\mu)=1$.  The independent
program \texttt{code/A/exact\_checks.py} checks both equations on
all six elements.  This finite example tests the naturality formula,
and does not claim that its groups are geometric boundary systems.
\end{proof}

\begin{proposition}[Refuted: unrestricted rigidity of the preliminary category]
\label{A:abstract-rigidity}
The preliminary axioms for $\mathcal B$ on scans 30--31 do not imply
$\operatorname{Aut}(f,\alpha)=1$.  For geometric hyperbolic target
systems over a field finitely generated over $\mathbf Q$, this
automorphism group is trivial.  For any hyperbolic geometric target,
it is trivial for morphisms whose image on the geometric group is
open.  Thus the corresponding Hom groupoids are equivalent to
discrete sets, with unique isomorphisms between isomorphic objects.
\end{proposition}
\begin{proof}
For the first assertion take no boundary objects, take
$E=G_K\times C_2\to G_K$, and take the homomorphism
$f:G_K\to E$, $\sigma\mapsto(\sigma,1)$, with its identity
base-commutation isomorphism.  All the preliminary conditions hold.
Every element of the second factor is a natural automorphism over
$BG_K$, so the automorphism group contains $C_2$.

For a geometric target, an automorphism in the correctly pinned
Hom groupoid has $p(\mu)=1$ by (A1).  The element $\mu\in\pi$
centralizes $f(E')$.  The elementary geometric-centralizer argument
below proves it is one when $K$ is finitely generated over
$\mathbf Q$ and $p(f(E'))$ is open.  If the geometric image is
open, geometric slimness already gives this conclusion.  The
boundary components $\mu_j$ are then determined by (A1) and the
injectivity of $\rho_i$, and equal one.  In a groupoid, trivial
automorphism groups imply that any isomorphism, when present, is
unique.  This proves the precise discreteness assertion.
\end{proof}

\begin{lemma}[Proved: arithmetic geometric-centralizer rigidity]
\label{A:geometric-centralizer}
Let $U/K$ be hyperbolic with $K$ finitely generated over
$\mathbf Q$.  For every closed $H\subset E=\pi_1(U)$ whose image
in $G_K$ is open, one has $C_\pi(H)=1$.
\end{lemma}
\begin{proof}
Suppose $1\ne z\in C_\pi(H)$.  A finite quotient of $\pi$ in
which $z$ has nontrivial image, followed by the inverse image of the
cyclic subgroup generated by that image, gives an open subgroup
$P\subset\pi$ in which the image of $z$ in $P^{\rm ab}$ is
nonzero.  More precisely choose a prime $\ell$ and an open normal
$N\subset\pi$ with $N\cap\overline{\langle z\rangle}$ contained
in the subgroup of $\ell$th powers of its $\ell$-component, and put
$P=N\overline{\langle z\rangle}$.  The cyclic quotient $P/N$
detects a nonzero $\ell$-component.  The group $P^{\rm ab}\otimes
\mathbf Z_\ell$ is a free module, by the surface presentation of
the corresponding finite \'etale cover.  Its class of $z$ is
therefore a nonzero vector, not merely a torsion class.

The cover descends to a finite extension $L/K$.  Intersect $H$ with
its arithmetic covering group and replace $L$ by the finite
extension corresponding to the image of that intersection.  The
resulting group still centralizes $z$ and surjects onto $G_L$.
It therefore fixes this nonzero homology vector.  Spread the cover
and its smooth compactification and boundary over a normal
finite-type arithmetic model, and take a closed point of good
reduction whose residue characteristic differs from $\ell$.
Arithmetic Frobenius at that point acts on the boundary homology
with eigenvalues of complex absolute value $q$, and on the
compactification's homology with absolute value $q^{1/2}$.
Here $q>1$.  Neither part has eigenvalue one, contradicting the
fixed vector.  This uses only the usual weight-one theorem for the
Jacobian and the explicit cyclotomic boundary module.  It applies
also to a proper cover, when the boundary part is zero.
\end{proof}

\subsection{Full and primewise commutants; a new metabelian counterexample}

Write $U_0=\mathbf P^1_{\mathbf Q}\setminus\{0,1,\infty\}$,
$F=\pi_1(U_{0,\overline{\mathbf Q}})$ and
$F_\ell=F^{(\ell)}$.  A peripheral outer automorphism preserves
the unordered conjugacy classes of the three cyclic inertia groups.
Its pure subgroup preserves each of the three labels.  On a quotient
of $F$, the peripheral groups in this definition mean their actual
images in that quotient.

\begin{proposition}[Known: the full and every pro-prime tripod commutant]
\label{A:tripod-commutant}
For every number field $K$ and every prime $\ell$, including $2$,
\[
 C_{\operatorname{Out}(F)}(G_K)=S_3,
 \qquad
 C_{\operatorname{Out}(F_\ell)}(G_K)=S_3.
\]
The embeddings of $S_3$ are those induced by the six fractional
linear transformations permuting $0,1,\infty$.  Intersecting with
the peripheral outer automorphism groups leaves these equalities
unchanged.  More generally the same argument gives
$C_{\operatorname{Out}(\pi)}(G_K)=\operatorname{Aut}_K(U)$ in
the full and the geometrically pro-$\ell$ setting for every
hyperbolic $U/K$ over a sub-$\ell$-adic field.
\end{proposition}
\begin{proof}
For a center-free profinite group $D$ and a continuous outer action
$r:G_K\to\operatorname{Out}(D)$, the associated extension is
canonically
\[
 \operatorname{Aut}(D)\mathop{\times}_{\operatorname{Out}(D)}G_K.
                                                               \tag{A2}
\]
Indeed conjugation and projection map the given extension into
(A2); injectivity is $Z(D)=1$, and every lift of an outer class
differs from the conjugation action of a chosen lift by a unique
inner automorphism.  If $[a]\in\operatorname{Out}(D)$ commutes
with $r(G_K)$, the rule $(b,\sigma)\mapsto(aba^{-1},\sigma)$
is an automorphism of (A2) over $G_K$.  Changing $a$ by an inner
automorphism changes this rule by geometric inner conjugacy.

The geometric full surface groups and geometric pro-$\ell$
hyperbolic groups are center-free.  Apply Mochizuki's Theorem A,
equivalently Theorem 16.5, to this arithmetic isomorphism
\cite{A_Mochizuki1999}.  It is an open continuous homomorphism
over the specified base field.  Applying the same theorem to its
inverse shows that the resulting dominant morphism is a unique
$K$-automorphism.  The theorem explicitly allows every prime in
its geometrically pro-prime version; no odd-prime assumption
enters this argument.  For $U_0$, its automorphisms are exactly
the six transformations permuting its deleted set, all defined
over $\mathbf Q$.  Each induces the stated commuting outer
automorphism.  This proves both surjectivity and injectivity.
\end{proof}

\begin{proposition}[Proved: the abelian peripheral commutant]
\label{A:abelian-commutant}
On $F^{\rm ab}=\widehat{\mathbf Z}^2$ the peripheral commutant is
$\widehat{\mathbf Z}^{\times}\times S_3$.  In the whole linear
automorphism group it is $\operatorname{GL}_2(\widehat{\mathbf Z})$.
In particular the smallest derived length giving commutant $S_3$
is greater than one.
\end{proposition}
\begin{proof}
Choose the first two cusp classes $e_1,e_2$; the third is
$-e_1-e_2$.  A matrix preserving each of the three corresponding
primitive rank-one direct summands has the form
$\operatorname{diag}(a,b)$, and preservation of
$\widehat{\mathbf Z}(e_1+e_2)$ forces $a=b$.
Allowing permutations adds exactly the displayed copy of $S_3$;
its intersection with the scalar subgroup is trivial.
The arithmetic action is $\chi\operatorname{id}$ and hence
commutes with every matrix.  This gives the two asserted groups.
\end{proof}

\begin{proposition}[Proved: infinitely many arithmetic-centralizing metabelian involutions]
\label{A:metabelian-commutant}
Let $M=F/\overline{[F',F']}$ be the full profinite metabelian
tripod group.  Its pure peripheral outer automorphism group
contains a closed subgroup
\[
       \prod_{p\geq5\ \mathrm{prime}} C_2
       \ \subset\ C_{\operatorname{Out}_{\rm lac}(M)}(G_K)
                                                               \tag{A3}
\]
for every number field $K$.  This subgroup also centralizes the
geometric $S_3$, and all of its elements act trivially on every
maximal pro-prime quotient of $M$.  Consequently the metabelian
commutant is strictly larger than $S_3$, and the requested
smallest derived length, when defined, is at least three.
\end{proposition}
\begin{proof}
Put $A=M^{\rm ab}=\widehat{\mathbf Z}^2$,
$R=\widehat{\mathbf Z}[[A]]$, and $c=[x,y]$.
The profinite Alexander calculation gives $M'=Rc$, freely of
rank one as a continuous $R$-module.  Chen--Deligne prove both
this statement and the precise outer-automorphism identity
\[
 \operatorname{IOut}(M)\simeq R^{\times,\epsilon=1}/A,
 \qquad
 \operatorname{Out}(M)\simeq
       (R^{\times,\epsilon=1}/A)\rtimes\operatorname{GL}(A).
                                                               \tag{A4}
\]
The semidirect action is the natural action on $R$, not a
new arithmetic action; these are their Corollary 4.1.5,
Proposition 4.3.7 and Theorem 5.2.4
\cite{A_ChenDeligne2017}.  We now specify the units, actual
automorphisms, peripheral conjugators and continuity needed
to use (A4).

Fix a prime $p\geq5$.  In $A^{(p')}$ let
$\theta:A^{(p')}\to C_3$ send both $X$ and $Y$ to the same
generator $t$, and put $H=\ker\theta$.  Normalized averaging
over $H$ defines an idempotent $e_H$ in
$\mathbf Z_p[[A^{(p')}]]$: in every finite quotient $B$ its
image is $|H_B|^{-1}\sum_{h\in H_B}h$.  The denominators are
$p$-adic units, and these averages are compatible under
quotient maps.  Define
\[
 e_p=e_H\left(1-\frac{1+t+t^2}{3}\right)
\]
in the coefficient-$p$ factor of $R$, and set its other
coefficient components equal to zero.  This is a nonzero
continuous idempotent.  On its factor, the prime-to-$p$
parts of $X,Y$ are primitive cube roots of unity.  Hence
$X-1$, $Y-1$ and $XY-1$ are units on $e_pR$, whether the
quadratic cyclotomic factor splits over $\mathbf Z_p$ or not.

Set
\[
 r_p=\frac{-2e_p}{(X-1)(Y-1)},\quad
 s_p=\frac{X(Y-1)r_p}{XY-1},\quad
 f_p=c^{r_p},\quad h_p=c^{s_p},\quad u_p=1-2e_p.
\]
All divisions are in $e_pR$ and are then extended by zero
on its complement.  The rule
\[
       \phi_p(x)=x,\qquad \phi_p(y)=f_p^{-1}yf_p             \tag{A5}
\]
defines a continuous endomorphism by the universal property
of $M$.  Direct multiplication in the extension
$1\to Rc\to M\to A\to1$ gives
\[
 \phi_p(c)=c^{u_p},\qquad
 \phi_p(xy)=c^{X(Y-1)r_p}xy=h_p^{-1}xyh_p.
\]
Moreover $(1+u_p)r_p=(1+u_p)s_p=0$, so
$f_p\phi_p(f_p)=h_p\phi_p(h_p)=1$ and $\phi_p^2=1$.
The action on $e_pM'$ is $-1$, so $\phi_p$ has order
exactly two and preserves the three labelled peripheral
conjugacy classes.  Its determinant in (A4) is $u_p$.

Here is the arithmetic check.  The action of every
$\sigma\in G_K$ on $A$ is the scalar $\chi(\sigma)$.
Every scalar unit preserves $H$, and it permutes the two
nontrivial characters of $C_3$.  Thus
\[
      \sigma_R(e_p)=e_p,
      \qquad \sigma_R(u_p)=u_p.                            \tag{A6}
\]
For an arbitrary element $(v,g)$ of the right-hand side of
(A4), conjugation on $(u,1)$ is
$(v,g)(u,1)(v,g)^{-1}=(g(u),1)$, because the unit group
is abelian.  Applying this exact formula with the
arithmetic $g=\chi(\sigma)\operatorname{id}$ and (A6)
proves $[\phi_p]\rho(\sigma)[\phi_p]^{-1}=\rho(\sigma)$.
No claim that an arithmetic representative commutes in
$\operatorname{Aut}(M)$ is needed or used.
The generators of $S_3$ act on the character $(1,1)$
modulo $3$ by the same character or its negative:
permuting $x,y$ fixes it, and replacing $(x,y)$ by
$((xy)^{-1},y)$ sends it to $(-2,1)=(1,1)$ modulo $3$.
They preserve $e_p$ as well, proving the $S_3$ assertion.

For any subset $J$ of the primes at least five, the
coefficientwise element $e_J=\sum_{p\in J}e_p$ is an
idempotent of $R$.  The same formulas give $\phi_J$.
Orthogonality $e_pe_q=0$ for $p\ne q$ makes the product
law exactly symmetric difference of subsets.  Every
finite quotient of $M$ uses only finitely many coefficient
primes in its finite derived subgroup; consequently this
map from $\prod C_2$ to $\operatorname{Aut}(M)$ is
continuous.  Its composite to $\operatorname{Out}(M)$ is
injective.  Indeed, an inner automorphism has determinant
in $A$ by (A4), whereas $u_J$ is one on its entire
coefficient-$3$ factor.  The embedding
$A\hookrightarrow\mathbf Z_3[[A]]^\times$ is injective,
as is seen in every finite group algebra.  Thus
$u_J\in A$ forces that monomial to be one; its nonzero
$p$-component then forces $J=\varnothing$.
Compactness of $\prod C_2$ gives the asserted closed image.

Finally every pro-$\ell$ quotient kills $e_p$ for
$p\ne\ell$.  For $p=\ell$ it kills $A^{(p')}$, and
$1-(1+t+t^2)/3$ becomes zero.  Thus (A5) becomes the
identity at every prime.  Every nonidentity member of
(A3) is IA; the faithful action of $S_3$ on $A$ shows
that it is outside $S_3$.
\end{proof}

\begin{proposition}[Proved: an independent finite certificate for the metabelian construction]
\label{A:finite-metabelian}
The prime-$5$ involution of Proposition~\ref{A:metabelian-commutant}
has a noninner peripheral image on the group
\[
 G=\mathbf F_{25,\rm add}\rtimes C_3,
 \qquad (v,i)(w,j)=(v+\zeta^i w,i+j),
 \qquad \zeta^2+\zeta+1=0.
\]
Take $x=(0,1)$ and $y=(1,1)$.  They generate $G$, and
the involution is $(v,i)\mapsto(-v,i)$.
\end{proposition}
\begin{proof}
In this quotient
\[
 c=(\zeta-1,0),\quad r=1+\zeta,\quad s=\zeta,
 \quad f=(3+4\zeta,0),\quad h=(4+3\zeta,0).
\]
These identities give $\phi(x)=x$,
$\phi(y)=f^{-1}yf$, $\phi(xy)=h^{-1}xyh$, and both
norm equations $f\phi(f)=h\phi(h)=1$.  Conjugates of
$\zeta-1$ span $\mathbf F_{25}$ over $\mathbf F_5$, so
$x,y$ generate all $75$ elements.  An inner automorphism
acts on the normal additive subgroup by $1,\zeta$ or
$\zeta^2$, whereas $\phi$ acts by $-1$.  Hence it is
noninner.

The independent program \texttt{code/A/exact\_checks.py}
constructs the multiplication, checks all $75^3$
associativity triples, all $75^2$ multiplication pairs
for $\phi$, all possible inner conjugators, the generator
claim, and the displayed peripheral and norm equations.
It also constructs all $1200$ automorphisms from a
semilinear part $v\mapsto a v$ or $a v^5$ and a
translation of the complement, and verifies that every
commutator with $\phi$ is inner.  All assertions pass.
The infinite arithmetic-centralizer conclusion uses
(A4)--(A6); this finite calculation is its concrete
certificate of nontriviality.
\end{proof}

\begin{remark}[Open: the same-level commutant at larger finite derived length]
\label{A:solvable-open}
The computation above settles the abelian and metabelian failures,
and the full and pro-prime equalities have been proved.  The
smallest $m\geq3$ with same-level commutant $S_3$ is not determined
here.  A theorem reconstructing the isomorphism class of a curve
from its $m$-step quotient is insufficient: it supplies a map
from isomorphism classes, whereas the required assertion concerns
every element of an automorphism group.  The precise additional
map in Mochizuki's Theorem A$'$ forgets three derived steps.
Yamaguchi's strong finite-step statements likewise specify the
image in a lower-step quotient \cite{A_Yamaguchi2024}; they do
not assert that the kernel of this forgetting map is trivial.
The involutions just constructed make this distinction visible by
an actual element and its images.  The remaining task is to
control that kernel or produce an arithmetic-centralizing lift
through the next derived quotient. Proposition~\ref{A:solvable-projection-limit}
below determines the exact lower-step images and the
prosolvable limit; it does not determine this same-level kernel.
Theorem~\ref{A:third-derived-IA} below eliminates the entire
IA image from length three to two, and
Corollary~\ref{A:third-derived-projection} identifies the
whole projected peripheral commutant with the geometric
$S_3$. Thus the remaining same-level question is the
kernel of the third-to-second-derived map, not the
lifting of the metabelian unit family.
\end{remark}

\subsection{Cyclotomic recognition and the full profinite obstruction}

\begin{proposition}[Refuted: the literal maximal-subgroup criterion of scan 483]
\label{A:literal-normalizer}
For a hyperbolic curve, the condition
$p(N_E(A))\subset G_K$ open does not characterize inertia among
all closed nontrivial subgroups $A\subset\pi$.  It is satisfied
by $A=\pi$, by every characteristic open subgroup of $\pi$, and
by $\overline{[\pi,\pi]}$.  None of these assertions requires
a section or a rational cusp.
\end{proposition}
\begin{proof}
Each displayed subgroup is normal in $E$, so its normalizer is
$E$, whose arithmetic image is open.  The geometric group $\pi$
of a hyperbolic curve is nonabelian, whereas cusp inertia is
procyclic.  Thus $A=\pi$ already contradicts the asserted
containment in a loop group.  An open characteristic subgroup
exists by intersecting the finitely many subgroups of any given
bounded index in the finitely generated profinite group $\pi$.
Such a subgroup is again a hyperbolic geometric fundamental
group, so is nonabelian.  The closure of the commutator is
characteristic and nontrivial.  This proves each asserted
normalizer equality and the counterexample.
\end{proof}

\begin{proposition}[Known: the precise cyclotomic criterion and transporter]
\label{A:cyclotomic-normalizer}
Let $U/K$ be a hyperbolic curve over a number field.  For a
nontrivial closed subgroup $A\subset\pi$, the following three
conditions are equivalent:
\begin{enumerate}
\item $A$ is contained in a cusp inertia group $L$;
\item $A$ is procyclic, $p(N_E(A))$ is open, and the action
of $N_E(A)$ on $A$ is the cyclotomic action;
\item $A$ is procyclic and some closed subgroup of $N_E(A)$
has open arithmetic image and acts on $A$ cyclotomically.
\end{enumerate}
The containing inertia is unique.  The maximal subgroups in
this class are precisely the full cusp inertia groups, and
\[
 N_E(L)=E_{\widetilde s},\qquad
 \operatorname{Transp}_E(L^d,L)=E_{\widetilde s}\quad(d\geq1),
                                                               \tag{A7}
\]
where the transporter means $\{g:gL^dg^{-1}\subset L\}$.
This includes nonrational cusps.
\end{proposition}
\begin{proof}
Nakamura's Theorem 2.1.1 is formulated for elements and a
cyclotomic normalizer, for every full class of finite groups
\cite{A_Nakamura1994}.  Stix's Theorem 20 writes exactly the
three displayed conditions for arbitrary nontrivial closed
subgroups \cite{A_StixCuspidal}.  Its field hypothesis is
$\operatorname{Hom}_{G_L}(T_\ell B,T_\ell J)=0$ for tori
$B$ and abelian varieties $J$, over every finite extension
$L/K$.  For number fields this follows by choosing a prime
of good reduction away from $\ell$: after a further finite
extension the torus is split, its eigenvalues have absolute
value $q$, and those on $J$ have absolute value $q^{1/2}$.
Thus all hypotheses of that theorem are verified here.

For clarity, its finite-cover step detects a noncuspidal
element in the proper homology of a finite cover, after
quotienting out the boundary classes.  The cyclotomic
normalizer supplies a nonzero $\mathbf Z_\ell(1)$ line in
that proper homology.  The preceding Frobenius comparison
excludes it.  This is a statement about an actual finite
cover of the full group; it does not assume that maximal
pro-prime quotients jointly separate that group.

Distinct cusp inertia groups have trivial intersection;
the same assertion holds for their nontrivial closed
subgroups, including a subgroup supported at a single
prime.  One proof passes to a finite cover separating the
two lifted cusps and then takes an abelian cover ramified
at one and at a third cusp but not the other.  This is the
peripheral centralizer/separation input in Nakamura's
Lemma 2.1.2 and Stix's Corollary 7.  It proves uniqueness,
and forces any element normalizing $A\subset L$ to fix
the same lifted cusp.  The normalizer is therefore its
decomposition group and acts cyclotomically.  Local
Kummer theory proves the converse.  Maximality now gives
the assertion for full inertia groups.

For (A7), $gL^dg^{-1}\subset L$ has nontrivial intersection
with $L$, hence the lifted cusp of $gLg^{-1}$ is the
original one.  Thus $g\in E_{\widetilde s}$.  Every element
of this decomposition group preserves $L$ and its subgroup
$L^d$, giving the reverse inclusion.  This proof follows
the lifted cusp and so does not require its residue field
to equal $K$.
\end{proof}

\begin{proposition}[Proved: uniqueness without a previously fixed surface type]
\label{A:unique-loops}
For an affine hyperbolic curve over a number field, a
hyperbolic surface loop structure on its geometric full
profinite group that is invariant under an open subgroup
of $G_K$ and satisfies scan 84(c) is the geometric loop
structure.  In particular its genus and number of holes
equal those of the curve.  The same conclusion holds
over a field finitely generated over $\mathbf Q$.
\end{proposition}
\begin{proof}
Pass to a finite extension that fixes each of the finitely
many proposed cusp classes.  Condition (c) supplies, for
each proposed loop $L'_j$, a cyclotomic normalizing subgroup
with full arithmetic image.  Proposition
\ref{A:cyclotomic-normalizer} puts $L'_j$ inside a genuine
inertia group.  A peripheral group of any hyperbolic
profinite surface presentation is its own normalizer in
the geometric group, by commensurable terminality
\cite[Lemma 1.3.7]{A_Mochizuki2004}.  Since a containing
procyclic group normalizes it, this inclusion is equality.
Distinct proposed loop classes remain distinct, so the
proposed cusps form a subset of the actual cusp classes.

Write $(g,r)$ for the actual type and $(h,s)$ for the
proposed one.  Thus $s\leq r$, and abelianization gives
$2h+s-1=2g+r-1$.  If $h=0$, this yields
$s=2g+r\geq r$, and hence $s=r$, $g=0$.
If $h>0$, quotienting by the closed normal subgroup
generated by the proposed loops gives the proper
surface group of genus $h$.  If any actual cusp
remained, the same quotient computed from the actual
presentation would be a free profinite group and
would have $H^2(-,\mathbf F_\ell)=0$.  A proper
surface group of positive genus has
$H^2(-,\mathbf F_\ell)=\mathbf F_\ell$, a contradiction.
Thus $s=r$ in this case too and the rank equality gives
$h=g$.  The actual quotient maps and cohomology groups
establish the comparison; no identification of differently
presented groups is presumed.

Over a finitely generated field, spread out the curves
in the finite-cover weight argument in the preceding
proposition and specialize at a good finite-residue-field
point.  This verifies the identical no-map-from-a-torus
hypothesis.  The remainder of the proof is unchanged.
\end{proof}

\begin{proposition}[Known: the cyclotomic-kernel centralizer form]
\label{A:cyclotomic-centralizer}
Let $K$ be finitely generated over $\mathbf Q$, and put
$\Gamma^\circ=\ker\chi$.  For a germ of section $s$, the
group $C_\pi(s(\Gamma^\circ))$ is either trivial or one
full cusp inertia group.  Every cusp inertia group
occurs, after a finite extension, as such a centralizer.
For a section whose image is already contained in a
cyclotomic normalizer of $A\ne1$, this centralizer contains
$A$ and identifies its unique cusp.
\end{proposition}
\begin{proof}
The full finite-cover proof is Proposition~\ref{F-centralizer}.
Its extra arithmetic input is the absence of nontrivial
divisible torsion in the Tate module of an abelian variety
over $K(\mu_\infty)$, obtained by specialization and Ribet's
theorem.  Thus Stix's recognition theorem applies over
$K(\mu_\infty)$ to every nonidentity fixed element.
The peripheral separation argument puts all such
elements in the same loop, and normality of
$\Gamma^\circ$ shows that the section normalizes that
loop.  Since $\chi=1$ on $\Gamma^\circ$, its whole
inertia group is fixed, giving equality.  Conversely a
nonzero rational tangent at any rationalized cusp
gives a section through its local decomposition group;
the preceding equality gives exactly that inertia.
The last assertion follows from
$s(\Gamma^\circ)$ centralizing $A$ and this equality.
\end{proof}

\begin{proposition}[Proved: the proof-level comparison of the two cyclotomic forms]
\label{A:normalizer-implications}
The cyclotomic normalizer recognition proof, applied over
$K(\mu_\infty)$ with its Ribet hypothesis, gives the
section-centralizer form (1 bis).  Conversely the same
finite-cover centralizer proof applies to every closed
$B\subset E$ whose image is open in $\ker\chi$ and
which lies over that kernel.  This strengthened
centralizer assertion gives (1) for arbitrary arithmetic
normalizers, without a splitting hypothesis.
\end{proposition}
\begin{proof}
The first assertion is the proof of
Proposition~\ref{A:cyclotomic-centralizer}.  For the
second, use the strengthened assertion in
Proposition~\ref{F-centralizer}.  Its proof needs only
that the image of $B$ is open in $\ker\chi$: after
descending a finite geometric cover with group $P$, the
intersection of $B$ with its arithmetic covering group
still has open image, and its action on $P^{\rm ab}$
is the ordinary arithmetic action.  Its kernel acts
inner and hence acts trivially on the abelianization.
Thus the Ribet and boundary-recognition steps are
unchanged.

Let $A\ne1$ have an open arithmetic normalizer $N$
acting cyclotomically.  Put $B=N\cap p^{-1}(\ker\chi)$.
Then $p(B)=p(N)\cap\ker\chi$ is open in $\ker\chi$,
and $A\subset C_\pi(B)$.  The strengthened assertion
puts $A$ in its unique inertia group.  The converse
follows from the local decomposition extension.
This is a reduction between the proofs with their stated
arithmetic input.  The statement quantified only over
sections does not itself supply a section of an arbitrary
normalizer extension; none is used here.
\end{proof}

\begin{proposition}[Known: a full profinite replacement over a p-adic local field]
\label{A:padic-cusp-reconstruction}
Let $K_i/\mathbf Q_{p_i}$ be finite and $U_i/K_i$
hyperbolic.  Every continuous isomorphism
$\pi_1(U_1)\xrightarrow{\sim}\pi_1(U_2)$ preserves
the geometric subgroups, the types $(g_i,r_i)$, and
the full cusp inertia groups and their arithmetic
normalizers.  The inertia reconstruction uses all
finite covers and the following formula, with the
usual integral weight convention:
\[
 r_V-1=\dim(V_\ell^{\mathrm{wt}\,2})
            -\dim(V_\ell^{\mathrm{wt}\,0})
 \quad(r_V>0),
 \qquad V_\ell=H_1(V_{\overline K},\mathbf Q_\ell),
                                                               \tag{A8}
\]
where $\ell\ne p_i$.  This gives a reconstruction even
when a Tate curve contributes a cyclotomic line to
proper homology.
\end{proposition}
\begin{proof}
Mochizuki's Lemma 1.3.8 first recovers the arithmetic
quotient.  The local-field group theory in Proposition
1.2.1 recovers its residue cardinality and Frobenius.
The argument of Lemma 1.3.9 then gives (A8)
\cite{A_Mochizuki2004}.  Here are its precise weight
and covering steps.  Boundary homology is of weight
two.  Proper homology is self-dual with a Tate twist,
so its weight-zero and weight-two dimensions agree.
The weight-zero part is unaffected by removing
punctures.  Subtracting the two dimensions therefore
leaves exactly the boundary dimension $r_V-1$.
Quasi-unipotence of local inertia makes these weights
independent of the chosen lift of Frobenius.  The
free/proper alternative and
$\dim H_1=2g+r-1$ then recover the type.

For a Galois $\ell$-power cover $Z\to V$, the equation
\[
          r_Z=[Z:V](r_V-1)+1                              \tag{A9}
\]
recognizes total ramification over exactly one cusp.
Indeed every ramified cusp contributes a deficit
$d(1-1/e)$ to the unramified count $dr_V$, with
$e$ a nontrivial $\ell$-power.  Two such deficits
already sum to at least $d$, whereas (A9) has
deficit $d-1$.  Thus there is one, and its index
is $d$.  The same property on pro-$\ell$ covers,
together with commensurable terminality, recovers
their maximal inertia subgroups.  This works for
every $\ell$, including $p_i$, because the cusp
counts were already obtained using an auxiliary
prime.  Repeating on every finite \'etale cover
recovers the cusps of the completed universal
cover; their stabilizers are the actual full
inertia and decomposition groups.  Every map
used in this construction is induced by the
given isomorphism, so it is functorial.
\end{proof}

\begin{proposition}[Proved: the degree-five CM obstruction, independently checked]
\label{A:cm-degree-five}
In the rational free Lie algebra on $x,y$, put
$\omega=[x,y]$, let $F_0(x)=y$, $F_0(y)=0$, and set
\[
 \epsilon_4(x)=\operatorname{ad}_x^4y,\qquad
 \epsilon_4(y)=[y,\operatorname{ad}_x^3y]
                 -[[x,y],\operatorname{ad}_x^2y],
 \qquad T_4=[F_0,\epsilon_4].
\]
With standard Lyndon bracketing for $x<y$,
\[
 T_4(x)=-2b_{xxxyy}-b_{xxyxy},\qquad
 T_4(y)=2b_{xxyyy}+3b_{xyxyy},\qquad T_4(\omega)=0.
                                                               \tag{A10}
\]
In bidegree $(3,2)$, $[x,L_4]$ has dimension one,
whereas adjoining $T_4(x)$ raises the dimension to
two.  Thus a CM eigenline with first term $x$
cannot satisfy the conjugacy equation for a
nonzero degree-four arithmetic term proportional
to $T_4$.
\end{proposition}
\begin{proof}
Expand the definitions in the free associative
algebra.  The two Lyndon words in (A10) are
linearly independent because their least words
have coefficient one.  The degree-four
bidegree-$(2,2)$ Lie space has its single
Lyndon basis vector $b_{xxyy}$, and its bracket
with $x$ is $b_{xxxyy}$.  The coefficient
$-1$ of $b_{xxyxy}$ proves the rank assertion.
The program \texttt{code/A/exact\_checks.py}
independently expands these integral polynomials,
checks $T_4(\omega)=0$, and computes both exact
matrix ranks.  It also checks the kernel of
$T_4$ in each degree at most seven; the only
kernel is the boundary line in degree two.

For the conjugacy assertion, an IA arithmetic
element normalizing a procyclic lift with
nonzero primitive abelianization has exponent
one on that lift.  In degree five its first
outer error must therefore be the bracket
$[x,v]$ for an inner correction $v\in L_4$.
Terms of the lift above degree one contribute
only from degree six onward.  Equation (A10)
contradicts the necessary degree-five equation.
\end{proof}

\begin{proposition}[Known: the audited pro-prime CM normalizer theorem]
\label{A:cm-prime-normalizer}
Put $K=\mathbf Q(i)$, $E_0:y^2=x^3-x$,
$U=E_0\setminus\{O\}$, and
$P_\ell=\pi_1(U_{\overline K})^{(\ell)}$.
For every prime $\ell$, every nontrivial closed
procyclic subgroup of $P_\ell$ with open arithmetic
normalizer is contained in a cusp inertia group.
At $\ell=5$, a primitive CM eigenline lifts to
such a subgroup of the class-four nilpotent
quotient, but has no such lift at class five.
\end{proposition}
\begin{proof}
We spell out the reduction audited in LM8 G.3--G.13.
For a generator $u$ and its first lower-central
term $u_d$, the normalizer's character has
Frobenius absolute value $q^{d/2}$, since
$H_1(E_0)$ has weight one.  Choose an open
subgroup $N\overline{\langle u\rangle}$ in
which $u$ has nonzero abelian image, descend
that finite cover, and intersect with the
normalizer.  The same character now occurs
in the cover's homology.  Its weight is one
or two, so $d=1$ or $2$.

For an inert prime, and also for $\ell=2$,
$H_1(E_0,\mathbf Q_\ell)$ has no line stable
under an open arithmetic subgroup.  A positive
power of Frobenius at $5$ has eigenvalues
$(-1+2i)^n,(-1-2i)^n$; their ratio is not
a root of unity and $i\notin\mathbf Q_\ell$.
Thus $d=1$ is impossible.  At a split prime
the two possible lines are the CM eigenlines.
The order-four automorphism $[i]$ kills the
outer degree-four characters $(3,1),(1,3)$.
The remaining character $(2,2)$ is nonzero:
this is Ishii's Corollary 4.10, whose hypotheses
are an imaginary quadratic field of class
number one, full CM over that field, and a
split prime at least five
\cite{A_Ishii2025}.  All hold here at every
split prime.  The degree-four comparison
is injective on the three special character
lines, so that nonzero character is a multiple
of $T_4$.  Proposition~\ref{A:cm-degree-five}
excludes $d=1$.

Consequently $d=2$ and $u_2=a\omega$, $a\ne0$.
The equality
$a\eta(n)\omega=a\chi_\ell(p(n))\omega$
for every normalizer element gives
$\eta=\chi_\ell\circ p$ by cancellation in
$\mathbf Z_\ell$.  Nakamura's Theorem 2.1.1
applied to the full class of finite
$\ell$-groups now identifies actual inertia
in $P_\ell$; it is not merely a claim in a
Malcev completion.

For the exact stopping level at $5$, take
CM-equivariant symplectic logarithmic coordinates
through class four.  The needed denominators
are units in $\mathbf Z_5$; straightening the
boundary through degree five also uses only
these units.  The special degree-one and
degree-three modules vanish, and the
degree-two module is
$\mathbf Z_5\operatorname{ad}\omega$.
Hence the arithmetic action differs from
its diagonal homology action, on this
quotient, by an inner action of a power
of the boundary.  Correcting the arithmetic
lift makes $\exp(\mathbf Z_5x)$ stable.
Its primitive abelianization makes it
maximal procyclic.  The nonzero degree-four
term and (A10) exclude a class-five lift.
This verifies the obstruction and every
prime-specific hypothesis used in LM8's
normalizer reduction.  The exact
$\kappa_{11}/48$ normalization is immaterial
for this nonvanishing argument.
\end{proof}

\begin{proposition}[Proved: full profinite recognition from one detected prime]
\label{A:cm-detected-normalizer}
With the preceding CM curve, let
$\pi=\pi_1(U_{\overline K})$ and
$R=\bigcap_\ell\ker(\pi\to P_\ell)$.
A nontrivial closed procyclic subgroup
$C\subset\pi$ with open arithmetic
normalizer is contained in actual full
inertia as soon as its image in one
$P_\ell$ is nontrivial.
\end{proposition}
\begin{proof}
Its $\ell$-component is a copy of
$\mathbf Z_\ell$ and maps injectively to
its nonzero image in $P_\ell$.
Proposition~\ref{A:cm-prime-normalizer}
therefore makes the action on that actual
component cyclotomic.  Apply Nakamura's
theorem with the full class of finite
groups, now to this component inside
$\pi$.  It lies in a full inertia group
$I$.  The centralizer in $\pi$ of any
nonidentity element of $I$ is $I$, even
when its exponent is supported at only
one prime.  Since all of $C$ commutes
with its nontrivial $\ell$-component,
it follows that $C\subset I$.

For this one-puncture group the cited
centralizer assertion also has a short
reduction to a primitive peripheral
generator.  An index-two cover has two
cusps, and its geometric group contains
the original boundary generator as a
free basis member.  If $w$ centralizes
its nontrivial power, then $w^2$ lies
in that cover and in its peripheral
centralizer.  The boundary has zero
image in $\pi^{\rm ab}$, so
$2w^{\rm ab}=0$.  Torsion-freeness of
$\widehat{\mathbf Z}^2$ gives
$w^{\rm ab}=0$, whence $w$ itself
lies in the cover.  Nakamura's
primitive-generator centralizer lemma
then gives $w\in I$.
\end{proof}

\begin{remark}[Open: the common-kernel normalizer problem]
\label{A:cm-common-kernel}
For the full CM group the remaining problem is
the existence or exclusion of a nontrivial
closed procyclic $C\subset R$ whose normalizer
in the arithmetic full profinite group has
open image in $G_{\mathbf Q(i)}$. Proposition~\ref{A:deeper-common-kernel}
below further forces any such $C$ into the strictly smaller
group $R_{\rm nil}(F')\subset F''$. No such
subgroup has been constructed or excluded
here.  For a general affine hyperbolic curve
the same bare-normalizer problem also remains
unresolved here.  It is not settled by the
primewise theorem or by the cyclotomic-kernel
centralizer theorem.
The additional restrictions in
Propositions~\ref{A:full-support-normalizer} and
\ref{A:cm-normalizer-field} below settle cofinite prime support
and all normalizer fields containing no CM subfield. Over an
imaginary quadratic normalizer field the residual support
can contain only split primes. These restrictions apply
to the actual full profinite subgroup in this remark.

The actual attacks and their exact stopping
points are as follows.  Lower central
quotients and all maximal pro-prime quotients
leave precisely $R$.  Passing to a finite
cover detects every nonidentity element in
homology, but its line can have weight one;
the proof does not make that character
cyclotomic.  For a CM detecting factor the
restriction to $G_{K^{\rm ab}}$ is trivial,
so replacing $K(\mu_\infty)$ by $K^{\rm ab}$
invalidates the no-torus-to-abelian-variety
hypothesis: both $T_\ell E_0$ and
$\mathbf Z_\ell(1)$ become trivial.

The Prym refinement has also been checked
directly.  For $Y\to U$ finite \'etale,
an element in $R$ has proper homology in
$\ker(V_\ell J_{\overline Y}\to V_\ell E_0)$.
In genus one this kernel is zero, so a
weight-one residual detector requires genus
at least two.  Sending its Jacobian factor
to an elliptic curve does not automatically
give a punctured map.  A map
$f:\overline Y\to E'$ restricts to
$Y\to E'\setminus\{q\}$ precisely when
$\operatorname{Supp}f^{-1}(q)$ lies in the
existing boundary.  Riemann--Hurwitz then
forces $\deg f\leq\deg(Y/U)$.  The concrete
map $[3]:E_0\to E_0$ on the cover
$E_0\setminus E_0[2]\to U$ still surjects
onto all of $E_0$, although it detects
proper Tate homology.  Thus it supplies
no such punctured map.  Adding punctures
requires an arithmetic-equivariant lift;
the class-five obstruction above is an
explicit case in which that lift fails.

A further approach using a normalizer
extension as a source for the Hom theorem
stops at the same precise issue: it is
an infinite-index arithmetic subgroup,
not the fundamental group of a finite-type
source for which the required open Hom
theorem applies.  Its projection onto an
open Galois group supplies no section by
itself.  The finite metabelian example
above confirms, by an actual mixed-prime
element, why primewise tests leave
additional information.  It is not a
counterexample inside $\pi$ to the
normalizer assertion.  A genuinely new
argument on the residual Prym part or
a compatible full profinite construction
is still required.
\end{remark}

\subsection{Slimness from global fields and the positive-degree Hom reduction}

\begin{lemma}[Proved: an elementary local embedding lemma]
\label{A:local-embedding}
Let $k$ be a global field and let $M/k$ be a finitely
 generated separably generated extension.  There is a
nonarchimedean place $v$ of $k$ and an embedding
$M\hookrightarrow k_v$ fixing $k$.  The valuation on $M$
pulled back from $k_v$ has completion $k_v$.  In particular
no nonidentity $k$-automorphism of $M$ preserves this
valuation.
\end{lemma}
\begin{proof}
Choose a smooth integral $k$-variety with function field
$M$ and an etale coordinate chart of dimension
$d=\operatorname{trdeg}(M/k)$.  A rational point of the
nonempty open image of this chart gives a closed point
with finite separable residue field $L/k$.  Choose a
nonarchimedean place of $k$ splitting completely in the
Galois closure of $L/k$.  Chebotarev supplies such places.
The chosen point is then a $k_v$-point.  The etale
implicit-function theorem gives a neighborhood of it
mapping bijectively to an open ball in $k_v^d$.
Choose coordinates in that ball algebraically independent
over $k$.  Such coordinates exist successively because
$k$ and each finitely generated extension previously
chosen are countable whereas each nonempty open ball
is uncountable.  The resulting evaluation has zero
kernel on the function field and embeds $M$ in $k_v$.
For $d=0$ the splitting place directly supplies the
embedding.  Since $k$ is dense in $k_v$, the completion
of $M$ is $k_v$.  An automorphism preserving the valuation
extends continuously to that completion and fixes its
dense subfield $k$, so is the identity.
\end{proof}

\begin{proposition}[Proved: the exact finite-type slimness statement]
\label{A:slim}
For $K$ finitely generated over $\mathbf Q$, the
centralizer in $G_K$ of every open subgroup is trivial.
The same holds for a field finitely generated over
$\mathbf F_p$ of positive transcendence degree.  For a
finite field it is false: $G_{\mathbf F_q}=\widehat{
\mathbf Z}$ and every open subgroup has centralizer
$\widehat{\mathbf Z}$.
\end{proposition}
\begin{proof}
Choose a global subfield $k\subset K$ over which $K$
is separably generated: take $k=\mathbf Q$ in
characteristic zero and $k=\mathbf F_q(t)$ with $t$
in a separating transcendence basis in positive
characteristic.  Suppose $\sigma$ centralizes an open
subgroup and $\sigma\ne1$.  Choose a finite Galois
extension $M/K$ such that $G_M$ is contained in that
open subgroup and $g=\sigma|_M\ne1$.  Enlarge it to
contain $\mu_\ell$ for an odd prime $\ell\ne
\operatorname{char}K$.  Lemma~\ref{A:local-embedding}
gives a discrete valuation $w$ on $M$ moved by $g$.
Its value group is discrete and nonzero, so normalize
it to $\mathbf Z$.  Weak approximation for the two
distinct valuations $w$ and $w\circ g$ gives $a\in M^*$
with
\[
             w(a)=1,\qquad w(g(a))=0.
\]
On $H^1(G_M,\mu_\ell)=M^*/M^{*\ell}$, conjugation
by $\sigma$ acts by the scalar $\chi_\ell(\sigma)$:
indeed it fixes every element of $G_M$ by assumption
and acts on the coefficients by that scalar.  Naturality
of Kummer theory therefore gives
\[
       [g(a)]=\chi_\ell(\sigma)[a]
                   \quad\text{in }M^*/M^{*\ell}.
\]
Applying $w$ modulo $\ell$ contradicts $0=
\chi_\ell(\sigma)$ in $\mathbf F_\ell$, since the
cyclotomic value is a unit.  This proves the two
positive assertions using only local embeddings,
Kummer theory, weak approximation and the global-field
splitting-prime theorem.  The finite-field assertion
follows from its Frobenius description.
\end{proof}

\begin{proposition}[Proved: positive peripheral degrees force an open image]
\label{A:positive-degree-open}
Let $P'$ and $P$ be the profinite geometric groups of
connected hyperbolic punctured curves in characteristic
zero, with nonempty boundary.  Choose positively oriented
peripheral generators $c'_j$ and $c_i$ from surface
presentations.  A continuous homomorphism $h:P'\to P$
satisfying
\[
 h(c'_j)=b_j^{-1}c_{\tau(j)}^{d_j\rho}b_j,
 \qquad d_j\in\mathbf Z_{>0},\quad
 \rho\in\widehat{\mathbf Z}^{\times}
\]
has open image.  The same multiplier $\rho$ is required
at every cusp; this is exactly the multiplier supplied
by the base path and the Kummer condition in a
morphism of $\mathcal B$.
\end{proposition}
\begin{proof}
First suppose the target has at least two cusps.
In $P^{\rm ab}$ the peripheral classes have the unique
relation $\sum_i[c_i]=0$ and span a free direct summand
of rank one less than the number of cusps.  Apply $h$
to $\sum_j[c'_j]=0$.  Since $\rho$ is a unit, the
positive integer totals
\[
                   D_i=\sum_{\tau(j)=i}d_j
\]
are all the same integer $D$.  Some source cusp exists,
so $D>0$; in particular every target cusp is hit.

Put $H=h(P')$, which is closed by compactness, and
let $Q\subset P$ be any open subgroup containing $H$.
Its corresponding finite covering has cusps $y$ above
the original cusps.  Write $e_y$ for the ramification
index at $y$.  A source cusp determines a unique lifted
target cusp: its nontrivial image lies in the
intersection of $Q$ with a conjugate of the original
inertia, and distinct lifted inertia groups have
trivial intersection.  For this lifted cusp,
$e_y\mid d_j$ because $\rho$ is a unit, and the local
exponent becomes $(d_j/e_y)\rho$.  The same boundary
relation in $Q^{\rm ab}$ forces the totals
\[
             D_y'=\sum_{j\mapsto y}d_j/e_y
\]
to one integer $D_Q>0$.  This equality also includes
cusps not initially known to be hit: their zero total
would force all totals to be zero, contradicting a
positive source exponent.  For any original cusp $i$,
\[
 D=\sum_{y\mid i} e_yD_Q
   =D_Q[P:Q],
\]
since the ramification indices in its finite fiber sum
to the covering degree.  Consequently $[P:Q]\le D$.
A closed subgroup whose overgroups of finite index all
have bounded index has finite index: it is the
intersection of those overgroups and the indices of
finite quotients of $P/H$ are bounded.  Hence $H$ is
open.

If the target has one cusp, its genus is positive.
Choose a degree-two abelian cover, using a nonzero
homology character $P\to\mathbf Z/2$.  Its peripheral
generator is a product of commutators, so maps to zero;
the covering therefore has two cusps.  Restrict $h$ to
the finite-index preimage of the covering group in
$P'$.  This is the group of a finite covering of the
source, with nonempty boundary.  A lifted source
exponent is $d_j e'_j/e_y$, a positive integer, where
$e'_j$ is the source ramification index.  The preceding
argument proves openness of the restricted image and
hence of the original image.
\end{proof}

\begin{proposition}[Known: the exact dominant relative Hom-form in the boundary category]
\label{A:relative-hom}
Let $K$ be sub-$p$-adic and let the source and target
be hyperbolic curves.  A morphism of the groupoid
systems of $\mathcal B$ with open geometric image is
induced by a unique dominant $K$-morphism of curves
satisfying its specified boundary condition.  Its local
integers $d_j$ are the actual ramification indices.
For case IV(d), where both boundaries are nonempty,
all morphisms of $\mathcal B$ have open geometric
image, so this proves the entire boundary-decorated
case.  The result applies to full profinite groups and,
for the open-image assertion of the Hom theorem, to
geometrically pro-$p$ groups.
\end{proposition}
\begin{proof}
After a finite extension splitting all cusps, use
$\alpha$ and a target arithmetic conjugation to pin the
base map.  In the boundary formula of
Proposition~\ref{A:two-morphisms}, write
$a_j=e_jb_j$ with $e_j$ in the target cusp decomposition
group having the same arithmetic projection as $a_j$
and $b_j$ geometric.  Conjugation by $e_j$ removes
$\chi(p(a_j))$ from the exponent.  Thus the geometric
map has the displayed form in
Proposition~\ref{A:positive-degree-open}, with the
single factor $\chi(\alpha)$, which is one in the
pinned presentation.  This proves openness in case
IV(d), including the one-cusp case.  Openness descends
through the finite base extension.

In all stated open cases Mochizuki's Theorem A
\cite{A_Mochizuki1999} gives the unique dominant map.
For smooth curves a dominant map extends uniquely to
the smooth projective compactifications.  Its local
map at a cusp is $t\mapsto u(t')^e$, with $e>0$;
Kummer theory therefore sends its oriented inertia
generator to the $e$th power, with the same base-path
factor as in Proposition~\ref{A:two-morphisms}.
Equality of the induced homomorphisms makes $e=d_j$
in $\widehat{\mathbf Z}$ and hence in $\mathbf Z$.
A source cusp cannot map to an interior point, since
its inertia would then be killed, contrary to $d_j>0$.
An interior source point cannot map to a deleted target
point because the map is defined on the source open
curve.  Thus the inverse image of the target boundary
is exactly the source boundary, as required by the
literal category.  The boundary decomposition map,
its conjugator and its Kummer map are recovered up to
the unique natural transformation of
Proposition~\ref{A:boundary-recovery} below.  This
proves full faithfulness in the specified category,
not merely equality of its underlying group maps.
\end{proof}

\subsection{Recovery of the Kummer decoration and the surface multipliers}

\begin{proposition}[Proved: recovery of the boundary morphism without rational cusps]
\label{A:boundary-recovery}
For geometric boundary systems over a field finitely
generated over $\mathbf Q$, the geometric-conjugacy
class of the global homomorphism $f:E'\to E$ in a
morphism of $\mathcal B$ determines that morphism up
to its unique natural transformation.  No rationality
assumption on the cusps is needed.  For every cusp
and every positive integer $d$,
\[
                   \operatorname{Transp}_E(L_i^d,L_i)=E_i.
\]
\end{proposition}
\begin{proof}
The transporter equality was proved in
Proposition~\ref{A:cyclotomic-normalizer}: any
transporter gives two cusp inertia groups with a
nontrivial intersection, so their lifted cusps coincide;
the resulting stabilizer is exactly the decomposition
group $E_i$.  This argument uses the decomposition
extension over $K(i)$ and does not require $K(i)=K$.

The image of each source inertia is nontrivial by its
positive degree.  It consequently determines its
unique target cusp inertia, up to the relevant
conjugacy, and hence the map of closed cusp orbits.
Since $E'_j$ normalizes that source inertia, its image
normalizes the nontrivial subgroup of the determined
target inertia.  The transporter assertion puts it
in the associated decomposition group.  Thus, once a
conjugator $a_j$ is chosen, the local homomorphism is
forced to be $\rho_{\tau(j)}^{-1}
\operatorname{Int}(a_j)f\rho'_j$.  Two choices of
$a_j$ differ on the left by an element of
$E_{\tau(j)}$; this is exactly $\mu_j$ in (A1).
Their arithmetic projections change the cyclotomic
factor by that same conjugation, so comparison of
$k_i$ gives the same positive integer $d_j$.

Two base paths for the same $f$ differ by an element
centralizing the open image of $p'$.  Proposition~
\ref{A:slim} says that they coincide.  Changing $f$
by geometric conjugation is the global $\mu$ in
(A1).  These statements construct the required natural
transformation, and Lemma~\ref{A:geometric-centralizer}
proves its uniqueness.  The argument simultaneously
recovers the residue-field homomorphisms, since they
are the projections of the forced local maps; no
extra choice appears at a nonrational cusp.
\end{proof}

\begin{proposition}[Proved: the exact ambiguity of the Kummer data]
\label{A:kummer-ambiguity}
A loop structure with cyclotomic compatibility determines
its equivariant Kummer isomorphisms up to one independent
factor in $\widehat{\mathbf Z}^{\times}$ for each
arithmetic orbit of geometric cusps.  The set of such
choices is a torsor under that product.  Requiring a
common positively oriented surface presentation reduces
this to the diagonal factor
$\widehat{\mathbf Z}^{\times}$ when the boundary is
nonempty.

For a proper hyperbolic curve the replacement is its
trace orientation
\[
 \operatorname{tr}:H^2(\pi,\widehat{\mathbf Z}(1))
                       \xrightarrow{\sim}\widehat{\mathbf Z}.
\]
Equivariant orientations form a
$\widehat{\mathbf Z}^{\times}$-torsor.  A dominant
map of proper curves has pullback compatible with the
traces by multiplication by its positive degree.
For a map from an affine curve to a proper curve this
condition is imposed on the filled source group
$\bar\pi'=\pi'/\overline{\langle\!\langle L'_j
\rangle\!\rangle}$; equivalently it is the relative
orientation of the source group with its boundary.
\end{proposition}
\begin{proof}
Two continuous isomorphisms from the rank-one module
$T$ to a fixed inertia group differ uniquely by a unit
of $\widehat{\mathbf Z}$.  Every such multiplication
commutes with the cyclotomic action.  Transport by
arithmetic conjugation forces exactly equality of
these scalars within a geometric cusp orbit, and
imposes no relation between different orbits.  This
proves the first torsor assertion by an explicit action
and its inverse.

In abelianization the peripheral classes have the
single primitive relation $\sum_i[c_i]=0$.
If replacement generators $c_i^{u_i}$ extend to a
surface presentation, their same relation forces
$u_1=\cdots=u_\nu$ for $\nu\ge2$.  For one cusp
there is already one scalar.  Conversely every common
unit is realized by Proposition~\ref{A:all-multipliers}
below, proving both necessity and sufficiency.

For a proper curve the comparison with a compact
oriented surface, or etale Poincare duality, identifies
$H^2(\pi,T)$ with $\widehat{\mathbf Z}$ and makes
the trace equivariant for the trivial action on the
right.  Its generators differ by precisely the stated
units.  Pullback of the fundamental cohomology class
under a finite map multiplies its trace by the degree:
this follows on every finite coefficient ring from
the pullback of a closed point divisor, whose total
multiplicity is the degree, and then by passage to the
inverse limit.

A map to a proper curve kills each source cusp
inertia, because it extends at that cusp.  It hence
factors through $\bar\pi'$, the fundamental group
of the smooth compactification.  Apply the preceding
degree identity to that factor.  In relative group
cohomology the same orientation is the trace on
$H^2(\pi',\{L'_j\};T)$; the boundary homomorphism
from $\bigoplus_jH^1(L'_j,T)$ has trace the sum of
the local residues.  The surface presentation proves
this by the single boundary relation.  Filling the
cusps identifies its degree-two orientation with the
proper orientation, giving the asserted compatibility.
This extension of the category is necessary: in the
literal category of scan 35 an index map from a
nonempty source boundary to an empty target boundary
does not exist.
\end{proof}

\begin{proposition}[Proved: every profinite unit is a surface multiplier]
\label{A:all-multipliers}
Let $2g-2+\nu>0$ and $\nu\ge1$.  In the profinite
surface presentation
\[
 P_{g,\nu}=\left\langle a_1,b_1,\ldots,a_g,b_g,
 c_1,\ldots,c_\nu\ \middle|\
 \prod_{j=1}^g[a_j,b_j]\prod_{i=1}^\nu c_i=1
                          \right\rangle^{\wedge},
\]
for every $\rho\in\widehat{\mathbf Z}^{\times}$
there is an automorphism taking every $c_i$ to a
conjugate of $c_i^\rho$.  Its images of the handle
generators, together with those conjugates, form a
privileged system.  When $g\ge1$ the same original
handle generators can be retained exactly for
$\rho=1$.  When $g=0$ all units are possible and
there are no handle generators to retain.
\end{proposition}
\begin{proof}
Realize the surface group by a curve and a reduced
rational divisor of degree $\nu$ over $\mathbf Q$.
For $g=0$ use the projective line.  For $g\ge1$ use
$y^2=x^{2g+1}-1$, which is smooth of genus $g$ and
has a rational point at infinity.  Distinct rational
$x$-values away from the branch locus give reduced
divisors of degree two; their union, with the point
at infinity when needed, gives every positive $\nu$.
The cusps need not individually be rational.  The
cyclotomic character of $G_{\mathbf Q}$ is surjective,
so an element with value $\rho$ acts on each inertia
as the $\rho$th power and permutes the cusp labels.
Compose its outer action with a surface homeomorphism
realizing the inverse permutation of the cusps.
That homeomorphism preserves the orientation and has
multiplier one.  Any representative of the resulting
outer automorphism has the asserted conjugacy property.
Transporting it through the complex surface comparison
proves the statement about the profinite presentation
alone; no arithmetic structure remains in its conclusion.

Here is an explicit multiplier $-1$.  Put
$Q=\prod_j[a_j,b_j]$ and $C=\prod_i c_i$.  The
substitution
\[
 R(a_j)=b_{g+1-j},\qquad R(b_j)=a_{g+1-j},\qquad
 R(c_i)=c_{\nu+1-i}^{-1}
\]
sends $Q$ to $Q^{-1}$ and $C$ to $C^{-1}$.
Since $C=Q^{-1}$, it respects the relation and
$R^2=1$.  Adjacent Hurwitz substitutions
\[
 h_i(c_i)=c_{i+1},\quad
 h_i(c_{i+1})=c_{i+1}^{-1}c_i c_{i+1}
\]
fix all other generators and preserve $C$.
Their product
$B=h_1(h_2h_1)\cdots(h_{\nu-1}\cdots h_1)$,
with the rightmost factor applied first, reverses
the cusp labels.  Thus $BR$ takes every $c_i$ to
a conjugate of $c_i^{-1}$ and supplies the requested
explicit element.

For the restriction on the handles, pass to the
class-two quotient in which all $c_i$ are central.
There $Q$ has class
$\omega=\sum_j a_j\wedge b_j$ in
$\bigwedge^2\widehat{\mathbf Z}^{2g}$, and the
relation is $C=Q^{-1}$.  The vector $\omega$ is
primitive and has no nonzero annihilator in
$\widehat{\mathbf Z}$.  Keeping all handle generators
fixed leaves $Q$ fixed, whereas replacing each cusp
by a conjugate of its $\rho$th power sends $C$ to
$C^\rho$.  Their relation therefore gives
$(1-\rho)\omega=0$, and hence $\rho=1$.
The identity realizes that value.  For $g=0$ the first
part already proves every value with no handles.
\end{proof}

\begin{proposition}[Proved: the geometric systems satisfy all five conditions]
\label{A:five-conditions}
For a smooth hyperbolic curve over a field finitely
generated over $\mathbf Q$, its system satisfies
(a)--(e) of scans 82--88, in their literal sense:
a section after a finite base extension; split local
decomposition extensions; cyclotomic compatibility;
the privileged surface presentation; and, in every
finite cover, purity of weight one and independence
of $\ell$ on homology after killing boundary homology.
Condition (c) is already a consequence of the local
Kummer equivariance and the surjectivity of the local
projection after splitting the cusp fields.  When the
boundary is nonempty, (b) implies (a).
\end{proposition}
\begin{proof}
A closed point of the nonempty open curve has finite
residue field over the base and gives a section after
that extension; it is an interior point and therefore
also gives the extra admissibility intended in (a).
At a cusp, a uniformizer defines the local section
and the Kummer description gives (b).  After the cusp
field is made rational, any $g\in E$ has the same
arithmetic image as some $e\in E_i$.  Writing
$g=\alpha e$ with $\alpha\in\pi$ gives
\[
 \operatorname{Int}(g)k_i(t)=
       \operatorname{Int}(\alpha)k_i(\chi(p(g))t),
\]
which is exactly (c), with the conjugator genuinely
geometric.  Finite base change gives the general
orbit form.  A section of $E_i$ composed with its
inclusion in $E$ gives the implication (b)$\Rightarrow$
(a), without its additional admissibility adjective.

An embedding of the base field in $\mathbf C$ and
the comparison theorem give (d) and its oriented
presentation.  Every finite arithmetic covering
subgroup comes from a finite etale cover; its
compactification has homology equal to the quotient
by the sum of its cusp homology.  This homology is
the Tate module of its Jacobian.  Spread out to a
good-reduction model.  At every closed point of good
reduction its characteristic polynomial is integral,
independent of $\ell$ different from the residue
characteristic, and all eigenvalues have absolute
value the square root of the residue cardinality.
These are exactly the requested clauses of (e).
\end{proof}

\begin{proposition}[Refuted: two proposed redundancies among the five conditions]
\label{A:condition-countermodels}
In the abstract group-data setting, (a), (b), (c),
and (d) do not imply (e).  Nor do (a), (b), (c), and
(e) imply (d).
\end{proposition}
\begin{proof}
For the first assertion take the proper genus-two
profinite surface group $P$, the empty boundary,
and $E=P\times G_K$ over a number field.  The product
section gives (a), (b) and (c) are vacuous, and the
surface presentation gives (d).  On $P^{\rm ab}$
every Frobenius is the identity, so its eigenvalues
have absolute value one and violate weight one.

For the second, take $P=T=\widehat{\mathbf Z}(1)$,
$E=T\rtimes G_K$, and one declared cusp with
$E_1=E$, $L_1=P$, and $k_1$ the identity.  The
splitting and the cyclotomic action give (a)--(c).
In every finite covering the inertia spans its entire
geometric abelianization, so the quotient required
in (e) is zero; its purity and independence statements
hold.  A surface presentation with one cusp has
abelian rank $2g$, whereas $P$ has rank one.  No
integer $g$ gives (d).
\end{proof}

\begin{remark}[Open: the remaining independence lattice]
\label{A:condition-lattice-open}
The implications and two countermodels just proved
settle those specific edges of the requested lattice.
A proof deciding every remaining implication among
(a)--(e), with the precise ambient axioms fixed, has
not been obtained here.  In particular no assertion
that a general outer arithmetic action virtually lifts
to an effective action is being inferred from the
weight condition.  The existence of a section after
finite extension is extra data in that question.
\end{remark}

\subsection{The added two-section stabilizer task}

To define the action in this task, retain the standard faithful
embedding $G_{\mathbf Q}\hookrightarrow\operatorname{Out}(F)$
and put
\[
 \mathcal A=\operatorname{Aut}(F)
       \mathop{\times}_{\operatorname{Out}(F)}
                   \widehat{\mathrm{GT}}_0.
\]
Its kernel is $F$, since $Z(F)=1$.  The pullback over
$G_L$ is the arithmetic fundamental group of the
tripod over $L$.  For an algebraic point $z$, write
$[S_z]$ for the $F$-conjugacy class of the image of
its section over its field of definition, together
with that projected arithmetic domain.  Conjugation
by a lift in $\mathcal A$ gives an action of
$\widehat{\mathrm{GT}}_0$ on these ambient subgroup
classes: two lifts differ by $F$, so the action is
independent of the lift.  This definition does not
assume that every element normalizes $G_{\mathbf Q}$.

\begin{proposition}[Proved: the stabilizer of the two specified point sections]
\label{A:two-section-stabilizer}
Let $z$ be one chosen root of $z^2-z+1=0$ and
$L=\mathbf Q(z)=\mathbf Q(\sqrt{-3})$.  For the
ambient action just defined, the simultaneous stabilizer
of $[S_{1/2}]$ and $[S_z]$ is exactly
\[
                         G_L.
\]
It is contained in $\ker\Delta$ for the character
of LM7\_PUSH E.65.  It does not contain $C^\varphi$:
complex conjugation belongs to $C^\varphi$ and
exchanges the two distinct root sections.
\end{proposition}
\begin{proof}
First we verify the arithmetic normalizer used in the
argument.  An outer automorphism commensurating the
faithful image of $G_{\mathbf Q}$ gives an isomorphism
between absolute Galois groups of two number fields.
The prescribed-isomorphism theorem of Uchida
\cite{A_Uchida1979} realizes it by an automorphism
$\sigma$ of $\overline{\mathbf Q}$.  Multiplication
by $\sigma^{-1}$ then centralizes $G_M$ for one
number field $M$.  Proposition~\ref{A:tripod-commutant}
identifies this centralizer with the geometric $S_3$.
Thus the full outer commensurator is
$G_{\mathbf Q}S_3$, and its subgroup preserving each
cusp label is exactly $G_{\mathbf Q}$: a nonidentity
member of that $S_3$ moves a label.  This supplies
the complete reduction behind LM3b D.8.

An element fixing $[S_{1/2}]$ normalizes its projection,
which is all of $G_{\mathbf Q}$.  The preceding result
puts it in $G_{\mathbf Q}$.  Arithmetic conjugation
transports $[S_z]$ to $[S_{\sigma z}]$.  Since $L$
is normal, both sections have domain $G_L$.
Injectivity of the point-section map makes the classes
equal exactly when $\sigma z=z$.  For these two
points this injectivity can be seen already by Kummer
theory: their ratio is a nontrivial third root of
unity, which is not a cube in $L$, since $L$ does
not contain a ninth root of unity.  Hence their
classes under the coordinate unit $t$ differ in
$H^1(L,\mu_3)$.  The stabilizer is therefore $G_L$.
The same conclusion holds for germs: after any finite
extension a sufficiently high power of three still
detects the two distinct roots.  If the phrase means
an unordered pair of these two classes with their
domains, it gives the same answer.  Interchanging
these domains would give a commensurator and hence
an arithmetic element, which cannot conjugate
$G_{\mathbf Q}$ onto its proper subgroup $G_L$.

We use only the arithmetic-vanishing part of the
character theorem E.65, not the existence assertion
E.58.  Its proof and the supplied referee's repair
were checked as follows.  On multiplier one, the two
frame signs are the restrictions to the same fourth
root $r=(3/4)^{1/4}$.  At $1/2$, the Legendre curve
is the quadratic twist by $2$ of $y^2=x^3-x$; at
$z$, it is the twist of $y^2=x^3+1$ by
$1/\sqrt{-3}$.  Over $\mathbf Q^{\rm ab}$ the first
twist is trivial, and the second twisting character
is $\sigma(r)/r$, because
$(-3)^{1/4}=\zeta_8\sqrt2\,r$.  The rational
three-torsion line of $y^2=x^3+1$ gives that same
sign; the projective three-torsion of $y^2=x^3-x$
gives the sign from its division polynomial.
The two signs therefore add to zero in $\mathbf F_2$.
The extension formula
$\Delta(a)=\delta(a\sigma^{-1})$, for an arithmetic
$\sigma$ with $\chi(\sigma)=\lambda(a)$, then gives
$\Delta|_{G_{\mathbf Q}}=0$.  In particular it
vanishes on the stabilizer just computed.

For the last assertion, write
$B=\widehat{C_2*C_3}=\langle s,r\mid s^2=r^3=1\rangle$.
The exact map of cote 140-3, scan 229, is
\[
 \varphi(c_0)=sr,\qquad \varphi(c_1)=s,
              \qquad\varphi(c_\infty)=r^{-1}.
\]
It is induced by the open immersion over $\mathbf Q$
from $\mathbf P^1_w\setminus\{0,1,\infty\}$ into
the orbifold with cusp $0$ and stacky points of orders
$2$ and $3$ at $1$ and $\infty$.  This orbifold is
$[U_0/S_3]$, with quotient coordinate
\[
        w=\frac{1728}{j}
          =\frac{27\lambda^2(1-\lambda)^2}
                  {4(\lambda^2-\lambda+1)^3}.
\]
The local ramification orders and the peripheral
product relation give precisely the three displayed
images, not merely a map with the same target.
Functoriality over $\mathbf Q$ supplies exterior
commutation for every arithmetic automorphism.
If its initial conjugator $b$ lies in $B$, the
pure-conjugator convention of equation (5) is obtained
as follows.  Changing the representative by
$\operatorname{Int}(a)$ with $a\in F$ changes $b$
to $\varphi(a)ba^{-1}$.  The map
$F\xrightarrow{\varphi}B\to S_3$ is onto, so
choose $a$ making the latter conjugator lie in $F$.
This proves membership in exactly the subgroup
$C^\varphi$ defined there.  In particular complex
conjugation belongs to it.  It exchanges $z$ and
$\bar z$, whose section classes were just shown to
be distinct, and so is outside $G_L$.
\end{proof}

\begin{remark}[Proved: the different stabilizer of the full root orbit]
\label{A:root-orbit-stabilizer}
Preserving $[S_{1/2}]$ and the unordered set of
\emph{both} root classes instead gives
$G_{\mathbf Q}$, by the same proof.  This is a
different specified set, related to the preceding
stabilizer by the exact index-two inclusion
$G_L\subset G_{\mathbf Q}$.  The result above answers
the task for one chosen root.  No conclusion that
$C^\varphi\subset G_{\mathbf Q}$ is being inferred
for this alternative task.
\end{remark}

\subsection{A further reduction of the full profinite common-kernel problem}

For a profinite group $D$, write
$R_{\rm nil}(D)=\bigcap_{j\ge1}\gamma_j(D)$,
where the lower central subgroups are closed.  This is
also the kernel of its maximal pronilpotent quotient.
For the geometric group $F$ of $E_0\setminus\{O\}$
put $P_n=\ker(F\to F^{\rm ab}/nF^{\rm ab})$.
The cover with group $P_n$ is
$E_0\setminus E_0[n]\xrightarrow{[n]}E_0\setminus\{O\}$.

\begin{lemma}[Proved: a weight-two normalizer test in an isogeny cover]
\label{A:isogeny-weight-test}
Let $E/K$ be an elliptic curve over a number field,
$D$ the maximal pro-$\ell$ group of $E$ minus a
nonempty finite divisor, and $B\ne1$ a procyclic
subgroup of $D$ normalized by a closed arithmetic
subgroup with open image in $G_K$.  If the image of
$B$ in $T_\ell E$ is zero, the normalizer acts on
$B$ cyclotomically after a finite base extension.
\end{lemma}
\begin{proof}
Make all punctures rational by a finite extension.
Use free pro-$\ell$ generators $x,y,c_1,\ldots,c_{r-1}$,
with the remaining cusp determined by the surface
relation.  Give $x,y$ weight one and each $c_i$
weight two in the integral Magnus expansion.  The
Magnus embedding separates the group, so a nonidentity
generator $b$ of $B\simeq\mathbf Z_\ell$ has a
first nonzero weighted degree $d$.  Its zero image in
$T_\ell E$ makes $d\ge2$.  The arithmetic action
preserves this filtration.  Its action on weight-one
generators is $T_\ell E$; its action on every cusp
generator in weight two is $\chi_\ell$.  Consequently
the eigenvalues of a good Frobenius on weight $d$
all have complex absolute value $q^{d/2}$: this
follows already in the finite free module of
homogeneous noncommutative polynomials, where the
eigenvalues are products of those on the generators.
Inner conjugation acts trivially on the first
nonzero weighted term.  If the normalizer character
is $\theta$, that term is a $\theta$-eigenvector.

There is also a finite pro-$\ell$ cover detecting
$b$ in homology.  Choose an open normal subgroup
$M$ for which $b$ survives in $D/M$, and use
$J=M\overline{\langle b\rangle}$.  The cyclic
quotient $J/M$ detects its homology class; the
abelianization of the finite covering group is free
over $\mathbf Z_\ell$, so this class remains
nonzero after tensoring with $\mathbf Q_\ell$.
Descend the cover and intersect its arithmetic group
with the original normalizer.  The projection remains
open and its homology vector again has character
$\theta$.  Good Frobenius eigenvalues on this covering
homology have absolute values $q^{1/2}$ or $q$.
Choose one good Frobenius common to both calculations.
It forces $d=1$ or $2$, and hence $d=2$.

In degree two, a group element with zero linear part
has leading term
\[
                a[X,Y]+\sum_i b_i C_i.
\]
For completeness, the group-like identity in the
Magnus Hopf algebra makes its quadratic part primitive.
The coefficients of $X\otimes X$ and $Y\otimes Y$
kill the $X^2,Y^2$ terms, and the coefficient of
$X\otimes Y$ makes the two mixed coefficients
opposite.  This argument uses torsion-freeness of
$\mathbf Z_\ell$ and works also at $\ell=2$.
On $[X,Y]$ the action is
$\det(T_\ell E)=\chi_\ell$ by the Weil pairing,
and every $C_i$ has that same character.  The
nonzero degree-two vector is torsion-free, so
$\theta=\chi_\ell$ on the indicated open Galois
subgroup, as asserted.
\end{proof}

\begin{proposition}[Proved: the residual normalizer lies deeper than the common pro-prime kernel]
\label{A:deeper-common-kernel}
Let $F=\pi_1((E_0\setminus\{O\})_{\overline K})$
for $E_0:y^2=x^3-x$ and $K=\mathbf Q(i)$.
A nontrivial closed procyclic subgroup with open
arithmetic normalizer is either contained in cusp
inertia or contained in
\[
                         R_{\rm nil}(F').
\]
The latter group is strictly smaller than
$R_{\rm nil}(F)=\ker(F\to\prod_\ell F^{(\ell)})$.
More precisely,
\[
 R_{\rm nil}(F')
   =\bigcap_{n\ge1,\ \ell\ \mathrm{prime}}
                   \ker(P_n\longrightarrow P_n^{(\ell)}),
                                                               \tag{A7}
\]
where the intersection is taken inside $F'$.
Thus every subgroup still left by this problem dies
in every maximal pro-prime quotient of every
multiplication-isogeny covering of the punctured
elliptic curve, and in particular in the full
metabelian quotient of $F$.
\end{proposition}
\begin{proof}
Proposition~\ref{A:cm-detected-normalizer} already
handles every subgroup detected by one $F^{(\ell)}$.
Suppose the remaining subgroup $A$ lies in
$R_{\rm nil}(F)\subset F'$.  It lies in every $P_n$.
If its image in some $P_n^{(\ell)}$ is nontrivial,
its $\ell$-component maps injectively to that
nontrivial procyclic image.  Its proper homology class
maps to zero under
\[
 T_\ell E_0\xrightarrow{[n]}T_\ell E_0,
\]
because $A$ dies in $F^{(\ell)}$.  This map is
injective, even when $\ell\mid n$, since the Tate
module is torsion-free.  Its proper homology class
is therefore zero.  Lemma~\ref{A:isogeny-weight-test}
shows that the normalizer character on this
$\ell$-component is cyclotomic after finite extension.
The full profinite recognition theorem in
Proposition~\ref{A:cyclotomic-normalizer}, including
its one-prime essential-cyclotomic version, puts that
component in one full inertia group.  Since $A$
centralizes its component, the inertia centralizer
puts all of $A$ in that inertia.  A nontrivial
inertia component survives in $F^{(\ell)}$, a
contradiction.  Thus $A$ lies in the right side of
(A7).

Here is a finite-quotient proof of (A7).  Every
map $P_n\to P_n^{(\ell)}$ kills
$R_{\rm nil}(F')$, since its restriction to $F'$
has pro-$\ell$ image.  Conversely an element of
$F'$ not in $R_{\rm nil}(F')$ survives in a finite
quotient $G$ of $F/R_{\rm nil}(F')$; the group
$G'$ is nilpotent.  Choose $n$ divisible by the
exponent of $G/G'$.  The image of $P_n$ is exactly
$G'$: it is contained there by that choice of $n$
and contains it because $F'\subset P_n$.
A Sylow projection of the finite nilpotent group
$G'$ detects the selected element and factors
through $P_n^{(\ell)}$.  This proves the reverse
inclusion.

Finally the explicit quotient
$G=\mathbf F_{25}\rtimes C_3$ of
Proposition~\ref{A:finite-metabelian} has
$G'=\mathbf F_{25}$ abelian, and
$\gamma_\infty(G)=\mathbf F_{25}\ne1$.
For a continuous finite quotient, the image of an
intersection of the nested closed lower central
groups is the intersection of their images, by
compactness.  Thus the image of
$R_{\rm nil}(F)$ in $G$ is $\mathbf F_{25}$,
whereas the image of $R_{\rm nil}(F')$ is
$\gamma_\infty(G')=1$.  This proves strictness by
an actual quotient, not by comparison of names.
\end{proof}

\subsection{Absolute Hom and the precise local comparison}

\begin{proposition}[Proved: an open arithmetic map over number fields preserves the geometric subgroup]
\label{A:absolute-geometric}
Let $X/K$ and $Y/L$ be hyperbolic curves over number fields.
Every continuous open homomorphism $f:E_X\to E_Y$ satisfies
$f(P_X)\subset P_Y$. Thus it induces a continuous open
homomorphism $\alpha:G_K\to G_L$.
\end{proposition}
\begin{proof}
The subgroup $H=f(E_X)$ is open in $E_Y$. Its projection
$J=p_Y(H)$ is an open subgroup of $G_L$, hence the absolute
Galois group of a number field. The group
$N=p_Yf(P_X)$ is closed, normal in $J$, and topologically
finitely generated because $P_X$ is. Mochizuki's
Theorem 2.4 in \cite{A_Mochizuki2006} states that a
finitely generated closed normal subgroup of the absolute
Galois group of a number field is trivial. Its field is
exactly the finite extension corresponding to $J$, so it
applies to $N$ and gives $N=1$. Factoring $p_Yf$ through
$E_X/P_X=G_K$ gives $\alpha$; its image is $J$, hence open.
This is the exact argument of LM3 D.11, with the requisite
normality taken in $J$, rather than assumed in all of $G_L$.
\end{proof}

\begin{proposition}[Known using the supplied Hom theorem: absolute open Hom over number fields]
\label{A:absolute-hom}
For hyperbolic curves $X/K$, $Y/L$ over number fields, the
natural map from dominant morphisms of schemes $X\to Y$ to
continuous open homomorphisms $E_X\to E_Y$, modulo inner
automorphisms of $E_Y$, is bijective. The arithmetic input
beyond the classical relative theorem is Theorem 1.1 of the
supplied October 2026 paper \cite{A_OpenAI2026Hom}; it remains
an unrefereed input, not an independently verified theorem
of this return. Its requested use of Uchida's local
correspondence for nonsurjective open homomorphisms is valid.
\end{proposition}
\begin{proof}
By Proposition~\ref{A:absolute-geometric}, an open $f$ induces
an open $\alpha:G_K\to G_L$. The supplied Theorem 1.1,
applied to the algebraic closures as solvably closed Galois
extensions, produces a unique equivariant embedding
$\bar j:\overline L\to\overline K$. An embedding between
algebraic closures of number fields is onto: its image is
algebraically closed and $\overline K$ is algebraic over it.
Hence $\alpha$ is injective and identifies $G_K$ with an
open subgroup of $G_L$. Its restriction is an embedding
$j:L\to K$.

Form the actual base change $Y_j=Y\times_{L,j}K$. Its
arithmetic group is
\[
 E_{Y_j}=E_Y\times_{G_L,\alpha}G_K,
 \qquad
 \widetilde f(e)=(f(e),p_X(e)).
\]
Projection $E_{Y_j}\to E_Y$ is injective and has open image,
since $\alpha$ does. Thus $\widetilde f$ is open. More
explicitly, injectivity of $\alpha$ gives
$f(P_X)=f(E_X)\cap P_Y$, an open subgroup of $P_Y$.
Mochizuki's Theorem A \cite{A_Mochizuki1999} now produces the
unique dominant $K$-morphism $X\to Y_j$ inducing
$\widetilde f$, and composition with $Y_j\to Y$ gives the
required scheme morphism. Conversely a dominant scheme
morphism restricts on algebraic constants to an embedding
$L\to K$ and has open arithmetic fundamental-group image.
The local inertia diagram of Proposition~\ref{A:typed-functor}
recovers all its Kummer maps and positive ramification degrees.

For uniqueness, conjugating $f$ in $E_Y$ changes $\bar j$
by an element of $G_L$ and leaves $j|_L$ unchanged. Once $j$
is fixed, any remaining arithmetic conjugator centralizes the
open image of $G_K$ in $G_L$, and is trivial by slimness.
The remaining conjugation is geometric, exactly the equivalence
relation in the relative Hom theorem. This proves bijectivity
with the stated inner-automorphism convention.

Here is the requested source check. On page 595 of
\cite{A_Uchida1981}, after stating the conjecture, Uchida
explicitly fixes a continuous homomorphism with \emph{open
image}; surjectivity is not imposed. On pages 595--596,
case (ii) of the local trichotomy associates a target prime
to a source prime whose pro-$\ell$ inertia survives, for
odd $\ell$, and puts its decomposition group in the target
decomposition group. Proposition 1, page 597, proves that
all but finitely many target primes occur. Its first paragraph
also proves compatibility with replacement of the target by
any finite extension inside the solvably closed extension.
These are exactly the hypotheses and conclusions used in
Theorem 2.4 of \cite{A_OpenAI2026Hom}. In particular, passage
to an open image changes the finite target field and is
already allowed in Uchida's argument. The finite exceptional
set in Proposition 1 is for a fixed $\ell$; no uniformity in
$\ell$ is being imported from it. This checks the step named
in the task, without treating it as a check of every later
argument in the 2026 main theorem.
\end{proof}

\begin{proposition}[Refuted in dimension zero; Known for local curve isomorphisms]
\label{A:local-absolute-isom}
There are finite extensions $k_1,k_2$ of $\mathbf Q_3$ with
isomorphic absolute Galois groups and no field isomorphism.
One explicit pair is
\[
 k_1=\mathbf Q_3(\zeta_3,\sqrt[3]{3}),\qquad
 k_2=\mathbf Q_3(\zeta_3,\sqrt[3]{4}).
\]
For hyperbolic curves over finite $p$-adic extensions, however,
the supplied Mochizuki--Tsujimura RNS preprint, Theorem D
or Theorem 3.12, gives the absolute Isom bijection
\[
 \operatorname{Isom}(X,Y)\ \xrightarrow{\sim}\
 \operatorname{OutIsom}(E_X,E_Y).
\]
These two assertions are related by the actual quotient maps
$E_X\to G_{k_1}$ and $E_Y\to G_{k_2}$: the first assertion
alone does not produce an isomorphism of their source groups.
\end{proposition}
\begin{proof}
The polynomials $t^3-3$ and $(t+1)^3-4$ are Eisenstein at
$3$. Their cubic fields have discriminant valuations $5$ and
$3$, respectively. Each splitting field has group $S_3$ and
unique maximal abelian subfield $\mathbf Q_3(\zeta_3)$,
because its cubic discriminant has square class $-3$.
Both splitting fields have degree six. The theorem on page 2
of Jarden--Ritter \cite{A_JardenRitter1979} therefore gives
$G_{k_1}\simeq G_{k_2}$: for local fields containing
$\zeta_p$, its exact criterion is equality of absolute degree
and of the maximal abelian subfield over $\mathbf Q_p$.

The discriminant exponents of the splitting fields are $11$
and $7$. To check this directly, decompose the regular
representation of $S_3$ as
$1\oplus\operatorname{sgn}\oplus2\operatorname{std}$.
The quadratic character has conductor exponent one; the
standard character has conductor exponent equal to the
cubic discriminant exponent, since the cubic permutation
representation is $1\oplus\operatorname{std}$. The
conductor-discriminant formula gives $1+2\cdot5=11$ and
$1+2\cdot3=7$. All the polynomial and finite character
calculations are independently checked in
\texttt{code/A/local\_field\_certificate.py} and its output.
An abstract field isomorphism between finite $p$-adic fields
preserves their canonical henselian valuation, hence is
continuous and fixes $\mathbf Q_p$; the unequal discriminants
exclude it.

For the positive curve assertion, use the full set of primes
as $\Sigma$ in the supplied RNS Theorem 3.12
\cite{A_MochizukiTsujimuraRNS}. We checked the exact
reduction used there. Theorem 3.11(i) first reconstructs the
geometric subgroup and the arithmetic quotient. Finite
residue fields satisfy its cyclotomic-kernel hypothesis;
Theorem 2.17 gives $\Sigma$-resolution of nonsingularities.
Theorem 3.11(ii) consequently makes the isomorphism preserve
closed-point decomposition groups. The last line of the
proof of Theorem 3.12 applies the point-theoretic absolute
reconstruction theorem, [AbsTopII], Corollary 2.9, to obtain
the unique scheme isomorphism. These steps apply to full
profinite groups and retain the stated preprint status.
An isomorphism of the quotients $G_{k_i}$ in the first part
is not asserted to lift through $E_X\to G_{k_1}$ and
$E_Y\to G_{k_2}$; if such a lift existed for hyperbolic
curves, this theorem would supply a field isomorphism,
contradicting the discriminant calculation.
\end{proof}

\begin{remark}[Open: unrestricted absolute open Hom over local fields]
\label{A:local-absolute-hom-open}
The preceding theorem handles isomorphisms and, by passing
to an open target subgroup, injective open homomorphisms
between local arithmetic curve groups: the target subgroup
is the group of its finite etale cover, so the given map is
an isomorphism onto that group and Theorem 3.12 applies.
For a general open homomorphism with nontrivial kernel, the
number-field proof above stops at an exact missing step.
Its projected geometric image is a finitely generated normal
subgroup of an open local Galois group; the number-field
vanishing theorem no longer applies, because that open
local Galois group is itself finitely generated. The local
Isom reconstruction uses both an isomorphism and its inverse
in Theorem 3.11(i), and supplies no morphism for this step.
No counterexample of hyperbolic curves, and no proof of this
remaining unrestricted Hom assertion, is obtained here.
The finite dimension-zero example is explicitly not substituted
for such a counterexample.
\end{remark}

\subsection{Kernels beyond number fields and beyond the full profinite group}

\begin{proposition}[Proved: constant curves and the unavoidable geometric kernel]
\label{A:constant-kernel}
Let $K/k$ be a regular finitely generated extension of a number
field, of positive transcendence degree, and let $U_0/k$ be a
hyperbolic curve. For $U=U_0\times_k K$ the kernel of the full
outer action is exactly
\[
 \ker(G_K\to\operatorname{Out}(P_U))
       =\ker(G_K\to G_k)=G_{K\overline k}.
\]
It is nontrivial. More generally, for every hyperbolic curve
$U/K$, with $K$ finitely generated over $\mathbf Q$ of positive
transcendence degree, the full outer-action kernel is nontrivial
and contains the inertia of every divisorial valuation at which
a smooth proper compactification with etale boundary extends
smoothly. It is contained in $G_{K\overline{\mathbf Q}}$.
\end{proposition}
\begin{proof}
For the constant curve the actual homotopy exact sequences
identify its outer representation with the composite
\[
 G_K\longrightarrow G_k\xrightarrow{\rho_{U_0/k}}
                  \operatorname{Out}(P_{U_0}).
\]
The second arrow is injective by Hoshi--Mochizuki, Theorem C,
and the first is surjective because $K/k$ is regular. This
proves the asserted equality, including its field-theoretic
identification by Galois theory.

For the general assertion take a smooth geometrically integral
model $S/k$ of $K$, shrinking it so that the compactification
$X/K$ and its reduced boundary extend to a smooth proper
curve $\mathcal X/S$ and a finite etale divisor
$\mathcal D\subset\mathcal X$. The homotopy sequence for
$\mathcal X\setminus\mathcal D$ gives a continuous outer
monodromy map $\rho_S:\pi_1(S)\to\operatorname{Out}(P_U)$
and the exact factorization
\[
          \rho_{U/K}=\rho_S\circ q_S,
       \qquad q_S:G_K\twoheadrightarrow\pi_1(S).
\]
At a codimension-one point of $S$, residue characteristic is
zero. Smooth specialization for the proper pair therefore
identifies the entire geometric fundamental group, and its
inertia is killed by $q_S$. Such inertia is nontrivial:
a root of a uniformizer gives a ramified finite extension of
the henselian discretely valued field, with nontrivial image
of inertia. A positive-dimensional smooth model has such a
codimension-one point after taking an affine open. This also
proves nontriviality in the constant case by an explicit
ramified extension. Finally Hoshi--Mochizuki, Corollary 6.2(ii)
\cite{A_HoshiMochizuki2011}, applied with the full set of primes,
puts the kernel inside the kernel of $G_K\to G_{\mathbf Q}$,
which is $G_{K\overline{\mathbf Q}}$.
\end{proof}

\begin{proposition}[Proved: the remaining kernel maps exactly to the congruence kernel]
\label{A:family-kernel-quotient}
For the model $S$ and the monodromy map in the preceding
proposition, let $N_S$ be the closed normal subgroup of
$G_K$ generated by divisorial inertia at the codimension-one
points of $S$. Then
\[
 N_S=\ker q_S,\qquad
 \ker\rho_{U/K}/N_S\ \xrightarrow{\sim}\ \ker\rho_S.
\]
For a universal hyperbolic curve over a fine level cover of
$\mathcal M_{g,r}$, the group on the right is precisely the
kernel of the geometric mapping-class monodromy
\[
 \widehat\Gamma_{g,r}[N]\longrightarrow
                         \operatorname{Out}(\widehat\pi_{g,r}).
\]
Here a level $N\ge3$ and a number field of definition are
chosen so that the moduli space and universal family are
schemes; the assertion is restricted to its positive-dimensional
components. Thus this is an actual quotient of the requested
function-field kernel, with its geometric meaning specified.
\end{proposition}
\begin{proof}
A finite extension of the function field is represented by
its finite normalization over $S$. Since $S$ is regular,
purity of the branch locus says that this normalization is
etale exactly when it is unramified at all codimension-one
points. The finite quotients of $G_K/N_S$ are consequently
exactly the finite etale covers of $S$, proving the first
identity. The surjectivity of $q_S$ and
$\rho_{U/K}=\rho_Sq_S$ give the second displayed isomorphism;
its map sends the class of $g$ to $q_S(g)$ and its kernel is
exactly $N_S$.

In the universal-family case the geometric fundamental group
of the fine level moduli space is the indicated profinite
mapping-class group, and its action on the fibre is its
usual geometric outer monodromy. Hoshi--Mochizuki,
Corollaries 6.4 and 6.5 \cite{A_HoshiMochizuki2011}, identify
the kernel of the full arithmetic universal monodromy with
this geometric kernel: an arithmetic monodromy element that
also lies in the geometric image has trivial image in the
number-field Galois group. Restriction to the chosen fine
level cover preserves the same assertion. Applying the
already proved quotient isomorphism now gives exactly the
claimed group, with no inference from a name or analogy.
\end{proof}

\begin{proposition}[Proved: every single-prime outer action has nontrivial kernel]
\label{A:proprime-kernels}
Let $K$ be a number field or a finite extension of
$\mathbf Q_p$, and let $U/K$ be hyperbolic. For every prime
$\ell$, including $\ell=2$, the outer action
\[
 G_K\longrightarrow\operatorname{Out}(P_U^{(\ell)})
\]
has nontrivial kernel. At every place of good reduction of
the compactified pair with residue characteristic different
from $\ell$, its inertia is contained in this kernel.
For a constant curve over a regular finitely generated
extension $K/k$, its kernel is the exact inverse image of
the corresponding kernel over $k$ under $G_K\to G_k$.
\end{proposition}
\begin{proof}
The group $P_U^{(\ell)}$ is finitely generated and pro-$\ell$.
The kernel of the action of its automorphism group on the
finite Frattini quotient is pro-$\ell$: on every finite
characteristic $\ell$-group quotient the analogous kernel is
an $\ell$-group, and these quotients are cofinal. Passing to
outer automorphisms preserves the conclusion that the outer
automorphism group has an open pro-$\ell$ subgroup. This is
also Proposition 4.1 of \cite{A_Mochizuki2006}.

No open subgroup of $G_K$ is pro-$\ell$. In the local case
it surjects onto $\widehat{\mathbf Z}$ through the maximal
unramified extension. In the number-field case its field
admits a cyclotomic $\mathbf Z_q$-extension for every prime
$q\ne\ell$; the intersection with that tower changes only
its finite initial portion, so the quotient is still
$\mathbf Z_q$. An injective map into the displayed outer
automorphism group would make an open subgroup of $G_K$
pro-$\ell$, a contradiction. This proves nontriviality for
all primes without an exception at two.

The inertia assertion is the prime-to-residue-characteristic
smooth specialization theorem for the proper pair; its
geometric pro-$\ell$ group is unchanged, so the inertia
outer action is trivial. For a constant curve the exact
base-change composite used in Proposition~\ref{A:constant-kernel}
also holds after maximal pro-$\ell$ completion. Taking its
kernel gives the stated inverse-image formula.
\end{proof}

\begin{proposition}[Known: the exact tripod kernel at odd regular primes]
\label{A:regular-prime-kernel}
Let $U_0=\mathbf P^1_{\mathbf Q}\setminus\{0,1,\infty\}$.
For an odd regular prime $\ell$, let $\Omega_\ell$ be the
maximal pro-$\ell$ extension of $\mathbf Q(\mu_\ell)$
unramified outside $\ell$. Then
\[
 \ker\bigl(G_{\mathbf Q}\to
          \operatorname{Out}(P_{U_0}^{(\ell)})\bigr)
                   =G_{\Omega_\ell}.
\]
\end{proposition}
\begin{proof}
Sharifi's Theorem 1.1 \cite{A_Sharifi2002} identifies the
fixed field of this kernel with $\Omega_\ell$ from the
freeness of the rational graded Lie algebra generated by
the odd Soule classes. Brown's Theorem 1.1 and its stated
consequence for the Deligne--Ihara action
\cite{A_Brown2012} supply this freeness. Thus both inputs
are theorems in the present application. The identification
is also stated with precisely these hypotheses in the
introduction of \cite{A_IshiiKernels2026}. Taking the subgroup
that fixes the common fixed field yields the displayed kernel
by the Galois correspondence. No step here extends the odd
regular-prime hypothesis to $\ell=2$ or to irregular primes.
\end{proof}

\begin{remark}[Open: the general kernel questions left after these explicit computations]
\label{A:kernel-open}
For a nonconstant curve over a field of positive
transcendence degree, the quotient in
Proposition~\ref{A:family-kernel-quotient} is not computed
in general here; universal families make the remaining
congruence kernel an exact instance of the problem. For a
general single-prime action, including the tripod at
$\ell=2$ and at irregular primes, nontriviality and the
specialization subgroups in Proposition~\ref{A:proprime-kernels}
do not determine the entire kernel. The Sharifi--Brown
calculation above handles its stated odd regular-prime
instances. The January 2026 once-punctured CM kernel paper
\cite{A_IshiiKernels2026} assumes a further elliptic
Deligne--Ihara assertion in its main Theorem A; its
application is therefore not an unconditional computation
of the remaining kernels, and is not used as one here.
\end{remark}

\subsection{Further consequences and limits of the finite-derived commutant}

\begin{corollary}[Refuted: the exact same-level strong metabelian assertion]
\label{A:strong-metabelian-counterexample}
For $U=\mathbf P^1_{\mathbf Q}\setminus\{0,1,\infty\}$,
the natural map
\[
 \operatorname{Aut}_{\mathbf Q}(U)\longrightarrow
 \operatorname{Aut}_{G_{\mathbf Q}}(E_U/F'')/\operatorname{Inn}(F/F'')
\]
is not surjective. Thus the instance $S_{2,0,U,U}$ of
Conjecture 0.1(2) in \cite{A_Yamaguchi2024}, whose
Remark 0.2 explicitly names $(2,0)$ as the best expected
pair, has a counterexample in the full profinite setting.
The counterexamples form the closed product of involutions
in Proposition~\ref{A:metabelian-commutant}.
\end{corollary}
\begin{proof}
We verify the precise passage from an outer commutant to an
arithmetic extension automorphism. The center of $M=F/F''$
is trivial. Indeed, a central element has abelian image
$a\in A$ acting trivially on the free $R$-module $M'$, so
$a=1$. It then lies in $M'$, and multiplication by $X-1$
kills it. The latter is a non-zero-divisor in $R$, as
Chen--Deligne's decomposition into coefficient power-series
rings shows; this is their Corollary 4.1.6. Thus the element
is trivial.

Consequently the map
\[
 E_U/F''\longrightarrow
 \operatorname{Aut}(M)\times_{\operatorname{Out}(M)}G_{\mathbf Q},
 \qquad e\longmapsto(\operatorname{Int}(e)|_M,p(e))
\]
is a continuous isomorphism. Its kernel is $Z(M)$; for
surjectivity, choose any lift of the specified Galois element
and correct its automorphism by the required element of $M$.
For the actual involution $\phi_J$ constructed above, the map
\[
             (a,\sigma)\longmapsto
                   (\phi_Ja\phi_J^{-1},\sigma)
\]
preserves this fibre product because $[\phi_J]$ centralizes
the outer Galois action. It is therefore an actual continuous
$G_{\mathbf Q}$-automorphism of the arithmetic metabelian
extension, restricting to $\phi_J$ geometrically. Changing
$\phi_J$ by an inner automorphism changes this map by the
allowed geometric inner automorphism.

Every geometric automorphism of the tripod belongs to $S_3$,
whose action on $M^{\rm ab}$ is faithful. The constructed
nonidentity classes act identically there and are noninner,
so none is geometric. The exact finite certificate of
Proposition~\ref{A:finite-metabelian} detects this failure.
The cited conjecture uses full geometrically metabelian
arithmetic groups in characteristic zero and its $n=0$
target is the same quotient; hence it is exactly the map
just disproved, not a statement obtained after forgetting
additional derived layers.
\end{proof}

\begin{lemma}[Known with independent proof: center-freeness of all finite-derived tripod groups]
\label{A:solvable-centers}
For $F$ free profinite of rank two and every $m\ge2$, put
$S_m=F/F^{(m)}$, with closed derived subgroups. Then
$Z(S_m)=1$. Its maximal prosolvable quotient
$S_\infty=\varprojlim_m S_m$ is also center-free.
The finite-step assertion is included in Theorem B of
Yamaguchi's May 2026 preprint \cite{A_Yamaguchi2026Centers};
the following transfer argument proves the case used here
independently.
\end{lemma}
\begin{proof}
The case $m=2$ was proved in the preceding corollary.
Inductively, a central element $z\in S_m$ maps trivially
to $S_{m-1}$, hence lies in
$N=F^{(m-1)}/F^{(m)}$. Let $H$ range over the open normal
subgroups of $F$ containing $F^{(m-1)}$. Since
$F^{(m)}\subset H'$, the image of $z$ in $H^{\rm ab}$
is defined and invariant under the finite group $F/H$.
Its image in $F^{\rm ab}$ is zero.

Write $i:H^{\rm ab}\to F^{\rm ab}$ for inclusion and
$\operatorname{Ver}:F^{\rm ab}\to H^{\rm ab}$ for the
continuous transfer defined using finitely many coset
representatives. Its defining product gives
\[
 \operatorname{Ver}(i(v))=\sum_{q\in F/H}q(v)
 \qquad(v\in H^{\rm ab}).
\]
For the invariant image $z_H$ this says
$[F:H]z_H=\operatorname{Ver}(i(z_H))=0$.
The abelianization of $H$ is a free profinite abelian group
by profinite Nielsen--Schreier, hence has no nonzero
integer torsion. Therefore $z_H=0$ integrally at every
prime.

Finally
$\bigcap_H H'=F^{(m)}$. To verify the nontrivial inclusion,
an element of $F^{(m-1)}$ outside $F^{(m)}$ survives in
some finite quotient of $F/F^{(m)}$. In that quotient the
image of $F^{(m-1)}$ is abelian. Taking $H$ to be its
inverse image gives an open normal subgroup for which the
selected element is not in $H'$. The converse inclusion
is immediate. Hence $z=1$. A central element of
$S_\infty$ has trivial image in every $S_m$, proving the
last assertion.
\end{proof}

\begin{proposition}[Known with explicit reduction: lower-derived commutants and the prosolvable limit]
\label{A:solvable-projection-limit}
Let $C_m$ be the commutant of the tripod's $G_K$-image in
the peripheral outer automorphism group of $S_m$, for a
number field $K$, and let $C_\infty$ be the corresponding
commutant on $S_\infty$. Then, for every $m\ge5$,
\[
 \operatorname{image}\bigl(C_m\to\operatorname{Out}(S_{m-3})\bigr)
                   =S_3,
       \qquad C_\infty=S_3.
\]
In particular no nonidentity metabelian involution of
Proposition~\ref{A:metabelian-commutant} lifts to an element
of $C_5$. These exact projections do not determine whether
$C_3$, $C_4$, or any $C_m$ itself equals $S_3$.
\end{proposition}
\begin{proof}
Center-freeness identifies the arithmetic extension with
$\operatorname{Aut}(S_m)\times_{\operatorname{Out}(S_m)}G_K$,
exactly as in the preceding corollary. Thus every member of
$C_m$ supplies a $G_K$-automorphism of that extension.
The truncated Hom theorem, Mochizuki's Theorem A-prime or
18.1 \cite{A_Mochizuki1999}, in its full profinite form
recorded in Theorem 0.4 of \cite{A_Yamaguchi2024}, gives a
dominant morphism $U\to U$ agreeing with that automorphism
on $S_{m-3}$. This is precisely why three derived layers
are lost; the statement of that theorem contains this
forgetting map.

A dominant self-map of the tripod has degree one. Indeed
its extension to $\mathbf P^1$ has degree $d$, the three
target punctures have inverse image contained in the three
source punctures, and the total ramification above them
is at least $3d-3$. Riemann--Hurwitz bounds total
ramification by $2d-2$, so $d\le1$. The map therefore
belongs to $S_3$. Conversely all of $S_3$ acts over $K$
and commutes with Galois, proving equality of the displayed
image. At $m=5$ the target is $S_2$, where our nontrivial
involutions are outside $S_3$; the nonlifting assertion
follows through that exact homomorphism.

For $\alpha\in C_\infty$, fix $n\ge2$ and apply the
preceding argument at $m\ge\max(5,n+3)$. It makes the
image of $\alpha$ in $\operatorname{Out}(S_n)$ a member
of $S_3$. Its faithful action on the common abelianization
makes this member independent of $n$. Compose with its
inverse, so that $\alpha$ is inner on every $S_n$.
For an actual representative of $\alpha$, the sets of
conjugators in $S_\infty$ realizing it on $S_n$ are
nonempty closed nested subsets of the compact group
$S_\infty$. Their intersection is nonempty, and an element
there realizes $\alpha$ on every quotient and hence on
$S_\infty$. Thus its outer class is trivial, proving
$C_\infty=S_3$ without asserting the same at a finite
level.
\end{proof}

\begin{proposition}[Proved: a full profinite unit family in the metabelian commutant]
\label{A:metabelian-unit-family}
The pure peripheral arithmetic commutant on $M=F/F''$
contains a closed subgroup isomorphic to
\[
                    \prod_p\mathbf Z_p^\times
                         =\widehat{\mathbf Z}^{\times}.
\]
Every member acts trivially on $M^{\rm ab}$ and on every
maximal pro-prime quotient of $M$, and centralizes $S_3$.
The family includes nontrivial coefficient-$2$ and
coefficient-$3$ members. This copy acts on the commutator
module; its image on the abelianization is the identity.
Theorem~\ref{A:third-derived-IA} below proves that no
nonidentity member of this entire family lifts to the
third-derived arithmetic commutant over any number field.
The preceding coefficient-five examples are retained
as explicit numerical certificates of this obstruction.
\end{proposition}
\begin{proof}
For every $p\ne3$, including $p=2$, use the idempotent
$e_p$ of Proposition~\ref{A:metabelian-commutant}.
The proof there only needs $3$ to be invertible in
$\mathbf Z_p$, so it applies at two as well. For $p=3$
take $B=(\mathbf Z/4)^2$, the quotient of $A^{(3')}$
on the two indicated generators. In $\mathbf Z_3[B]$
take the sum of the character idempotents indexed by
\[
 \mathcal S=\{(a,b)\in(\mathbf Z/4)^2:
                     a\ne0,\ b\ne0,\ a+b\ne0\}.
\]
This set has six elements. It is invariant under
$(a,b)\mapsto(b,a)$, $(a,b)\mapsto(-a-b,b)$, and
$(a,b)\mapsto(-a,-b)$. The first two generate the
peripheral $S_3$ action, and the third is the only
nonidentity scalar unit action modulo four. The Fourier
sum is defined over $\mathbf Z_3$: conjugate characters
occur in pairs and its coefficients are in
$\frac18\mathbf Z$. Lift it to
$\mathbf Z_3[[A^{(3')}]]$ by normalized averaging over
the kernel of $A^{(3')}\to B$, and then to the
coefficient-$3$ factor of $R$. Call the resulting
idempotent $e_3$. It has augmentation zero. On all
six character factors the prime-to-three parts of
$X,Y,XY$ are nonidentity fourth roots of unity;
therefore $X-1,Y-1,XY-1$ are units on $e_3R$.

For $t=(t_p)_p\in\prod_p\mathbf Z_p^\times$ put,
coefficient by coefficient,
\[
 u(t)=1+\sum_p(t_p-1)e_p,\qquad
 r(t)=\frac{u(t)-1}{(X-1)(Y-1)},\qquad
 s(t)=\frac{X(Y-1)r(t)}{XY-1}.
\]
The denominators are inverted only on the displayed
idempotent factors. Each component of $u(t)$ is $t_p$
on $e_p$ and one on its complement, hence is a unit.
Set $f(t)=c^{r(t)}$, $h(t)=c^{s(t)}$, and define
\[
 \phi_t(x)=x,\qquad \phi_t(y)=f(t)^{-1}yf(t).
\]
The same explicit multiplication as before gives
$\phi_t(c)=c^{u(t)}$ and
$\phi_t(xy)=h(t)^{-1}xyh(t)$, proving preservation
of all three labelled peripheral classes. For two
parameters the conjugator exponent in the composite is
\[
 r(t)+u(t)r(v)=\frac{u(t)u(v)-1}{(X-1)(Y-1)}=r(tv).
\]
Thus $\phi_t\phi_v=\phi_{tv}$ and
$\phi_{t^{-1}}$ is its inverse. This proves that these
are actual continuous automorphisms, without relying
on a determinant-only existence assertion.

All $e_p$ are fixed by the scalar cyclotomic action and
by $S_3$. Formula (A4) then proves arithmetic and
$S_3$ commutation for their outer classes. The
parameters $t_p\in\mathbf Z_p$ are fixed coefficients
for that natural group-ring action. Their images in
every maximal pro-prime group are the identity, since
the corresponding nontrivial prime-to-$p$ character
idempotent has zero image there.

For injectivity, suppose $u(t)$ is a monomial from
$A$, as an inner automorphism requires. For each $p$
apply the exact coefficient and group quotient map
\[
 R\longrightarrow\mathbf Z_p[[A^{(p)}]].
\]
It sends $u(t)$ to one, so it sends that monomial to
one. The embedding of $A^{(p)}$ into this completed
group ring is faithful, and hence its $p$-component
is trivial. This holds for every $p$, so the monomial
is one. The nonzero torsion-free $\mathbf Z_p$-module
$e_pR$ now gives $t_p=1$ for every $p$. Thus the outer
map is injective. Every finite quotient uses finitely
many coefficient primes and finitely many residue
digits of their parameters, proving continuity.
Compactness then gives a closed image.

The independent certificate
\texttt{code/A/metabelian\_units\_certificate.py}
checks the coefficient-three Fourier idempotent
exactly over $\mathbf Q[B]$, its three required
inverses and its symmetry. It also checks explicit
noninner peripheral images on
$((\mathbf Z/4\mathbf Z)[\zeta_3])_{\rm add}\rtimes C_3$
and $((\mathbf Z/9\mathbf Z)[i])_{\rm add}\rtimes C_4$,
of orders $48$ and $324$, respectively. In the first
the parameter is $-1$; in the second it is $2$.
All $2304$ and $104976$ multiplication pairs, respectively,
the generation claims, all candidate inner conjugators,
and the three peripheral conjugacy equations pass.
\end{proof}

\subsection{A full profinite normalizer restriction from complex-conjugation primes}

\begin{lemma}[Proved: descent of a power of the normalizer character]
\label{A:normalizer-character-descent}
Let $U/K$ be an affine hyperbolic curve over a number
field, let $A\ne1$ be a closed procyclic subgroup of
$F=P_U$, and let $N\subset E_U$ be a closed subgroup
normalizing $A$, with $p(N)=G_L$ open in $G_K$.
Fix a prime $\ell$ for which $A_\ell\ne1$ and write
$\theta_\ell:N\to\mathbf Z_\ell^\times$ for the action
on $A_\ell\simeq\mathbf Z_\ell$. Put
$d=\ell-1$ for odd $\ell$, and $d=1$ for $\ell=2$.
There is a continuous character
\[
 \lambda_\ell:G_L\longrightarrow\mathbf Z_\ell^\times,
 \qquad \lambda_\ell\circ p=\theta_\ell^d.
\]
Moreover either $A$ lies in one cusp inertia group, or,
over some finite extension $L_1/L$, there is a character
$\psi:G_{L_1}\to\mathbf Z_\ell^\times$ occurring as
a nonzero line in the Tate module of the Jacobian of
a finite cover of $U$, with
\[
              \lambda_\ell|_{G_{L_1}}=\psi^d.
\]
In the latter case $\psi$ is pure of Frobenius weight one
and has Hodge--Tate weights in $\{0,1\}$, with the
convention that the cyclotomic character has
Hodge--Tate weight one.
\end{lemma}
\begin{proof}
The free profinite group $F$ is torsion-free, so its
nontrivial pro-$\ell$ subgroup $A_\ell$ is
$\mathbf Z_\ell$. Put $D=N\cap F$. Its action on
$A_\ell$ has trivial pro-$\ell$ part. To see this,
take a pro-$\ell$ Sylow subgroup $S$ of $D$ containing
the normal subgroup $A_\ell$. Closed pro-$\ell$
subgroups of a free profinite group are free pro-$\ell$:
their $\ell$-cohomological dimension is at most one,
which is the free pro-$\ell$ criterion. For $g\in S$,
the group generated by $A_\ell$ and $g$ has a normal
procyclic subgroup and procyclic quotient, hence is
metabelian. A free pro-$\ell$ group on at least two
generators has finite nonmetabelian $\ell$-group
quotients. Therefore this subgroup is procyclic and
$g$ centralizes $A_\ell$. The image of a Sylow
subgroup is a Sylow subgroup of the image, proving
the assertion about the pro-$\ell$ part.

For odd $\ell$ the remaining subgroup of
$\mathbf Z_\ell^\times$ has exponent dividing
$\ell-1$; for $\ell=2$ it is trivial. Thus
$\theta_\ell^d$ kills $D$ and descends through the
actual compact quotient $N/D=G_L$, giving
$\lambda_\ell$ and its continuity. No splitting of
$N\to G_L$ is asserted.

Choose an open characteristic subgroup $V\subset F$
which does not contain a topological generator of
$A_\ell$, and put $H=VA_\ell$. The finite cyclic
quotient $H/V$ shows that $A_\ell$ has nonzero image
in $H^{\rm ab}$. This abelianization is torsion-free,
so its restriction to $A_\ell\simeq\mathbf Z_\ell$
is injective. Descend the geometric cover defined by
$H$ to a finite extension and let $E_H\subset E_U$
be its open arithmetic fundamental group. Its
geometric kernel is exactly $H$. The open subgroup
$N\cap E_H$ has open arithmetic image, say
$G_{L_1}$ after further finite base extension.
Its conjugation on $H^{\rm ab}$ is the ordinary
arithmetic action: the geometric kernel lies in $H$
and acts inner. Consequently the displayed line is
a genuine $G_{L_1}$-stable line with character
$\psi=\theta_\ell|_{N\cap E_H}$, and
$\lambda_\ell|_{G_{L_1}}=\psi^d$.

In geometric homology the boundary kernel of
$H^{\rm ab}\to T J$ is a quotient of the permutation
module on the cusps tensored with
$\widehat{\mathbf Z}(1)$. If the line has zero image
in $T_\ell J$, split the finitely many cusp fields
by a finite extension. Its character then equals
$\chi_\ell$. The exact cyclotomic-normalizer
criterion, Proposition~\ref{A:cyclotomic-normalizer},
puts $A_\ell$ inside an inertia group $I$.
Every element of $A$ normalizes $A_\ell$, and that
same proposition gives $N_F(A_\ell)=I$. Hence
$A\subset I$.

In the other case the image of the line in
$T_\ell J\otimes\mathbf Q_\ell$ is injective.
The good-reduction Frobenius eigenvalues of this
Jacobian have weight one. Tate's Hodge--Tate theorem
for abelian varieties gives weights zero and one,
and the line inherits one of these at each local
place. These are exactly the asserted properties
of $\psi$; they do not depend on good reduction
at $\ell$ itself.
The power character $\lambda_\ell$ is consequently
Hodge--Tate after finite extension and hence before
it. It is also unramified outside finitely many
places: outside the bad-reduction and ramification
places for the finite extension, its inertia group
lies in $G_{L_1}$ and $\psi^d$ is unramified there.
\end{proof}

\begin{lemma}[Proved: one-dimensional characters at a complex-conjugation prime]
\label{A:conjugation-prime-character}
Let $M/\mathbf Q$ be a finite Galois number field.
Fix a complex conjugation $c\in\operatorname{Gal}(M/\mathbf Q)$,
allowing $c=1$ when $M$ is totally real. Let $\ell$
be unramified in $M$ with Frobenius conjugacy class
equal to that of $c$. Every Hodge--Tate character
$\lambda:G_M\to\mathbf Q_\ell^\times$ is a finite-order
character times an integral power $\chi_\ell^a$.
When $M$ is totally real the conclusion holds for
every prime $\ell$, without the Frobenius condition.
\end{lemma}
\begin{proof}
We spell out the maps in this use of the theory of
abelian $\ell$-adic representations. Tate's local
algebraicity theorem for a one-dimensional Hodge--Tate
representation gives a $\mathbf Q_\ell$-defined
algebraic character
\[
 h:(\operatorname{Res}_{M/\mathbf Q}\mathbf G_m)_{\mathbf Q_\ell}
                              \longrightarrow\mathbf G_{m,\mathbf Q_\ell}
\]
which describes $\lambda$ on sufficiently small
local units at the places over $\ell$. See
\cite[Chapter III, Theorem on p.~III--7 and Proposition
on p.~III--9]{A_SerreAbelian}. With the convention
calibrated by $\chi_\ell$, write its character as
$h=\sum_\tau n_\tau[\tau]$, $n_\tau\in\mathbf Z$.

Choose a modulus for $\lambda$ as in Chapter III,
Proposition 2 on p.~III--11 of that source. Global
reciprocity then says that $h$ is one on a subgroup
of finite index in the global unit group: on the
other local unit factors the character is trivial,
and the product on a principal idele is one. The
equality $\prod_\tau\tau(u)^{n_\tau}=1$ is algebraic,
so it remains true after passing to the complex
embedding fixed above. Dirichlet's unit theorem
therefore gives
\[
                  n_\tau+n_{c\tau}=2a
                  \quad\text{independent of }\tau.
\]
For a totally real field it gives directly
$n_\tau=a$ for every $\tau$. This is also the
character-lattice calculation in Chapter II,
Proposition on p.~II--30 of the cited source.

In the other case choose a place over $\ell$ whose
Frobenius is the particular element $c$; the
Chebotarev conjugacy condition permits this choice.
The algebraic character $h$ is defined over
$\mathbf Q_\ell$, so its coefficient vector is fixed
by the local decomposition group, which contains
$c$. Hence $n_{c\tau}=n_\tau$. The preceding equality
forces $n_\tau=a$ for all $\tau$, and $a$ is an
integer since the $n_\tau$ are integers. Thus $h$
is exactly the $a$th power of the norm character.

The quotient $\lambda\chi_\ell^{-a}$ consequently
has trivial algebraic character on every local
unit group over $\ell$. It has a modulus at every
other place as well. The global class-field map
therefore factors this quotient through a finite
ray class group. Its image is finite, proving the
claim. This argument neither presumes that the
Frobenius values are rational numbers nor replaces
$\mathbf Q_\ell$-valuedness by a stronger coefficient
condition.
\end{proof}

\begin{proposition}[Proved: full-support subgroups and a positive-density exclusion]
\label{A:full-support-normalizer}
In the situation of Lemma~\ref{A:normalizer-character-descent},
let $M/\mathbf Q$ be a finite Galois extension
containing $L$, and let $\mathcal P_M$ be the set of
primes unramified in $M$ whose Frobenius conjugacy
class is the class of a complex conjugation.
Then either $A$ is contained in one cusp inertia
group, or
\[
                       A_\ell=1
                     \qquad(\ell\in\mathcal P_M).
\]
The set $\mathcal P_M$ has positive Dirichlet density.
In particular every such $A$ with nontrivial
pro-prime component at all but finitely many primes
is contained in a unique inertia group. This includes
every $A\simeq\widehat{\mathbf Z}$. If $L$ is totally
real, every nontrivial $A$ satisfying the stated
normalizer condition is contained in a unique
inertia group, with no restriction on its prime
support.
\end{proposition}
\begin{proof}
Suppose $A$ is noncuspidal and $A_\ell\ne1$ for a
prime in $\mathcal P_M$. The preceding descent
lemma supplies $\lambda_\ell$ and, after finite
extension, a weight-one character $\psi$ with
$\lambda_\ell=\psi^d$ on that open Galois subgroup.
Hodge--Tate decompositions descend through finite
extensions, so $\lambda_\ell|_{G_M}$ is Hodge--Tate.
This descent is stated in
\cite[Chapter III, Appendix, Corollary 2 on
p.~III--32]{A_SerreAbelian}.
Lemma~\ref{A:conjugation-prime-character} gives
\[
              \lambda_\ell|_{G_M}
                         =\varepsilon\chi_\ell^a
\]
with $\varepsilon$ finite and $a\in\mathbf Z$.
Pass to a common finite extension also killing
$\varepsilon$. There the exact character equality is
$\psi^d=\chi_\ell^a$. Comparison at a prime of good
reduction different from $\ell$ gives $2a=d$,
because $\psi$ has Frobenius weight one. Comparison
of Hodge--Tate weights at a place above $\ell$
gives $a\in\{0,d\}$. These equalities contradict
$d>0$. Hence $A_\ell$ must be trivial.

Chebotarev gives the Dirichlet density among
rational primes
\[
 \lim_{s\downarrow1}
 \frac{\sum_{\ell\in\mathcal P_M}\ell^{-s}}
      {\log(1/(s-1))}
          =\frac{|\operatorname{Cl}(c)|}{[M:\mathbf Q]}>0.
\]
A subgroup with cofinite prime
support therefore cannot be in the noncuspidal
case. If $L$ is totally real, choose its totally
real Galois closure for $M$. The last assertion of
Lemma~\ref{A:conjugation-prime-character} supplies
the same contradiction at any prime with
$A_\ell\ne1$. Every nontrivial torsion-free procyclic
group has such a prime. Uniqueness of the containing
inertia follows from
Proposition~\ref{A:cyclotomic-normalizer}.
\end{proof}

This restriction applies in particular to the common-kernel
subgroups of Remark~\ref{A:cm-common-kernel}: their prime
components are tested inside actual finite-cover homology,
even when their images in every maximal pro-prime quotient
of the original curve vanish. A remaining noncuspidal
subgroup must have prime support disjoint from the stated
positive-density set, in addition to lying in
$R_{\rm nil}(F')$ in the CM example. The argument does not
exclude a single-prime subgroup at a prime outside that set.

\begin{proposition}[Proved: the necessary CM fields in the remaining normalizer case]
\label{A:cm-normalizer-field}
Retain the notation of
Lemma~\ref{A:normalizer-character-descent}. If $A$ is
noncuspidal, its arithmetic normalizer field $L$
contains a CM subfield. Moreover, for every prime
$\ell$ with $A_\ell\ne1$ and every finite Galois
extension $M/\mathbf Q$ containing $L$, there is a
CM subfield $E\subset M$ admitting a field embedding
$E\hookrightarrow\mathbf Q_\ell$. In particular,
when $L$ contains no CM subfield, every nontrivial
procyclic subgroup satisfying the stated open
normalizer condition is cuspidal. When $L$ itself
is imaginary quadratic, a remaining noncuspidal
subgroup can have nontrivial prime components only
at rational primes that split in $L$.
\end{proposition}
\begin{proof}
Choose a nontrivial component $A_\ell$ and the
character $\lambda_\ell$ constructed above. Let
$G=\operatorname{Gal}(M/\mathbf Q)$ and
$H=\operatorname{Gal}(M/L)$. The algebraic character
of $\lambda_\ell|_{G_M}$ has a coefficient function
$n:G\to\mathbf Z$ satisfying $n(gh)=n(g)$ for
$h\in H$: restriction from $G_L$ pulls back the
character by the norm from $M$ to $L$. Write
\[
 n(g)=a+b(g),\qquad
 a=|G|^{-1}\sum_{g\in G}n(g).
\]
The unit-character calculation of
\cite[Chapter II, section 3.1, pp.~II--29--30]{A_SerreAbelian}
gives
\[
 b(cg)=-b(g)
 \quad\text{for every complex conjugation }c\in G.
\]
This is the left action, the same action as in the
preceding lemma. The function $b$ is nonzero.
Otherwise all local algebraic exponents are one
integer $a$, so $\lambda_\ell|_{G_M}$ is a finite
character times $\chi_\ell^a$. Its equality to a
power $\psi^d$ after finite extension gives both
$2a=d$ and $a\in\{0,d\}$, the contradiction already
proved above.

Let $J\lhd G$ be generated by all products $c_1c_2$
of complex conjugations. Every such product fixes
$b$, so $b$ is left $J$-invariant. Since $J$ is
normal, it is also right $J$-invariant. Together
with right $H$-invariance this makes $b$ right
$JH$-invariant. All complex conjugations have the
same image modulo $J$, and that image is central.
No complex conjugation lies in $JH$: if one did,
left multiplication by it would act on $b$ as
right multiplication modulo $J$, hence would fix
$b$, contradicting $b\ne0$ and $cb=-b$.

It follows that $F=M^{JH}\subset L$ has a
nontrivial involution induced by complex conjugation.
All complex conjugations induce this same involution;
its fixed field is totally real, whereas none of
the embeddings of $F$ is real. Thus $F$ is a
totally imaginary quadratic extension of a totally
real field, proving the asserted CM subfield of $L$.
This proves explicitly the field-theoretic content
of the norm reduction in Chapter II, section 3.3
of the cited source.

For the additional local condition put
$P=\operatorname{Stab}_G(b)$ for the left action.
It contains $J$ and contains no complex conjugation.
The same centrality argument proves that
$E=M^P$ is a CM field. Since the algebraic
character is defined over $\mathbf Q_\ell$, the
local decomposition subgroup $D_v\subset G$ at
the chosen embedding into $\overline{\mathbf Q}_\ell$
fixes $b$. Hence $D_v\subset P$. The restriction
of this embedding to $E$ takes values in
\[
                        M_v^{D_v}=\mathbf Q_\ell,
\]
giving the required field embedding. This argument
also covers primes ramified in $M$.

Finally, for imaginary quadratic $L$ take $M=L$.
Its only CM subfield is $L$ itself. A quadratic
number field embeds in $\mathbf Q_\ell$ precisely
when it has a degree-one completion there, which
in the quadratic case means that $\ell$ splits.
This proves the last assertion, including the
exclusion of ramified primes in this case.
\end{proof}

\subsection{A third-derived obstruction detected by two explicit covers}

\begin{proposition}[Proved: the coefficient-five involution does not lift to derived length three]
\label{A:third-derived-obstruction}
Let $\phi_5$ be the actual metabelian involution in
Proposition~\ref{A:metabelian-commutant}, with coefficient-five
determinant $1-2e_5$, where the cubic character defining $e_5$
takes both cusp generators to the same generator of $C_3$.
There is no outer automorphism of $S_3=F/F'''$ that commutes
with $G_{\mathbf Q}$ and projects to $[\phi_5]$ on $S_2$.
In fact there is no such lift commuting with
$G_{K_0}$, where
\[
 E:y^2=x^3+16,\qquad
 K_0=\mathbf Q(\mu_{15},E[5]).
\]
Thus this nonidentity metabelian class is already excluded
at the third derived level. More generally, for the unit
family $\phi_t$ of Proposition~\ref{A:metabelian-unit-family},
the existence of a commuting third-derived lift over either
of these fields forces $t_5\equiv1\pmod5$.
These assertions do not determine the entire commutant
at that level.
\end{proposition}
\begin{proof}
We construct a finite metabelian quotient and two of its
intermediate covers. The argument uses the full homology
of these covers, which is retained at the third derived
level.

The first cover is the restriction over $U_0$ of
\[
 E\xrightarrow{[5]}E\longrightarrow\mathbf P^1,
       \qquad t=(y+4)/8.
\]
The last map is a cyclic cubic cover: putting $w_1=x/4$
gives $w_1^3=t(t-1)$. Its group is
$G_1=E[5]\rtimes C_3$, where the order-three action is
$(x,y)\mapsto(\zeta_3x,y)$. It is irreducible on $E[5]$,
since $T^2+T+1$ is irreducible over $\mathbf F_5$.
Its abelian cubic character at the first two cusps is
$(1,1)$.

For the second cover take the Fermat curve
\[
 C:X^5+Y^5+Z^5=0,
 \quad r=X^5/Z^5,\quad
 w_2=\frac{r-\zeta_3}{r-\zeta_3^2},\quad
 t=\frac{w_2^3}{w_2^3-1}.
\]
The cyclic permutation $(X,Y,Z)\mapsto(Y,Z,X)$ acts
on $r$ by $r\mapsto(-1-r)/r$. Its action on $w_2$ is
multiplication by a primitive cube root of unity.
Together with the fifth-root diagonal transformations
it gives
\[
 G_2=\mathbf F_5^2\rtimes C_3,\qquad
 T(a,b)=(b-a,-a).
\]
Here the cubic intermediate extension is
$w_2^3=t/(t-1)$, so the two cusp exponents are $(1,-1)$.
The diagonal action and the cyclic permutation are defined
over $K_0$.

Let $Y\to U_{0,K_0}$ be the fibre product of these two
finite etale covers. It is connected with constant deck
group
\[
 G=G_1\times G_2
   =(V_1\times V_2)\rtimes(C_3\times C_3),
 \qquad |G|=5625,
 \quad V_1=E[5],\quad V_2=\mathbf F_5^2.
                                                               \tag{A18}
\]
Indeed the two cubic Kummer classes are independent:
their valuation vectors at $0,1$ form a matrix of
determinant $-2$, nonzero modulo $3$. Each $G_i$ has
only $C_3$ as a nontrivial proper quotient, by the
irreducibility just checked. A common nontrivial quotient
would therefore give a common cubic subextension, which
the independence excludes. This also proves connectedness
after reduction at the prime used below. In particular
$G$ is metabelian and center-free.

Let $q:F\twoheadrightarrow G$ be the geometric quotient.
On $V_1\times V_2$ the unit $1-2e_5$ acts as
\[
                 S(v_1,v_2)=(-v_1,v_2).                       \tag{A19}
\]
To verify an actual induced automorphism, rather than only
an action on a module, extend $S$ by the identity on the
displayed complement $C_3^2$. Its value on $q(x)$ can be
corrected by a unique conjugation in $V_1\times V_2$:
$q(x)-1$ is invertible on both irreducible summands.
Call the resulting automorphism $\gamma$. It fixes
$q(x)$ and has the same value on $q([x,y])$ as $q\phi_5$.
The values of $\gamma(q(y))$ and $q(\phi_5(y))$ have the
same image in $C_3^2$. Their quotient is in $V_1\times V_2$;
equality of their commutators with $q(x)$, and the same
invertibility of $q(x)-1$, makes this quotient zero.
Consequently
\[
                         q\phi_5=\gamma q.
\]
The correcting inner conjugation is trivial on $V_1\times V_2$,
so $\gamma$ still acts there by (A19).

We now specify two subgroups of order five using a place
of $K_0$ over $211$. This prime splits completely in $K_0$.
Indeed $211\equiv1\pmod{15}$, and direct enumeration of
$E(\mathbf F_{211})$ gives $225$ points and all $25$ points
of $E[5]$ rational over $\mathbf F_{211}$. The reduction
is good and the five-torsion is unramified; thus its
Frobenius acts identically on the generators of $K_0$.
Choose the place with $\zeta_3=14$ and $\zeta_5=55$ in
the residue field. A point of $E[5](K_0)$ is uniquely
specified by its reduction at this place. Put
\[
 P_0=(2,92),\quad u_0=[2]P_0=(103,117),\quad v_0=(3,1),
 \qquad
 J=\langle(u_0,v_0)\rangle,
 \quad J_-=\langle(-u_0,v_0)\rangle.                          \tag{A20}
\]
Thus $\gamma(J)=J_-$. Both quotient covers $Y/J$ and
$Y/J_-$ have degree $1125$ over $U_0$.

Here is the complete finite rule for counting their points
over $\mathbf F_{211}$. For $t\in\mathbf F_{211}\setminus\{0,1\}$
the Frobenius class can meet $J$ only if both $t(t-1)$ and
$t/(t-1)$ are cubes. There are exactly $24$ such values:
\[
\begin{aligned}
 &12,13,28,64,65,68,87,88,89,90,97,98,\\
 &114,115,122,123,124,125,144,147,148,184,199,200.
\end{aligned}                                                   \tag{A21}
\]
For each take any cube roots $w_1,w_2$, and set
\[
 \begin{aligned}
 Q&=(4w_1,8t-4)\in E(\mathbf F_{211}),& P&=[9]Q,\\
 r&=(w_2\zeta_3^2-\zeta_3)/(w_2-1),&
 v&=\bigl(\log_{55}(r^{42}),
             \log_{55}((-1-r)^{42})\bigr)\in\mathbf F_5^2.
 \end{aligned}                                                   \tag{A22}
\]
The logarithms here mean the unique exponents in
$\{0,1,2,3,4\}$, so this is a finite formula involving no
choice of analytic logarithm.

The actual elliptic Kummer class in the $[5]$-torsor is
conjugate under $C_3$ to $[2]P$. To check the factor,
the Frobenius endomorphism $\pi$ has trace $-13$ and
norm $211$. This ordinary CM reduction has endomorphism
ring $\mathbf Z[\zeta_3]$, and solving these two equations
gives $\pi=1+15\zeta_3$ or $1+15\zeta_3^2$. For a point
$R$ with $[5]R=Q$ its Kummer class is
\[
 (\pi-1)R=3\zeta_3^iQ
             =2\zeta_3^i[9]Q,
\]
since $[15]Q=0$. The Fermat Kummer class is exactly $v$:
Frobenius multiplies fifth roots of $r,-1-r$ by
$r^{42},(-1-r)^{42}$. Thus a representative, up to
conjugacy, is $\delta_t=([2]P,v)$.

For any $\delta=(u,v)\in V_1\times V_2$, its number of
fixed cosets in $G/J$ is
\[
 \#(G/J)^\delta
   =125\sum_{i,j=0}^2
       \mathbf1_J(\zeta_3^iu,T^jv).                            \tag{A23}
\]
For a fixed complement coordinate there are $125$ cosets;
translation fixes them exactly when its conjugate lies
in $J$. This proves (A23) directly. The analogous formula
holds for $J_-$. Substitution of the $24$ values (A21)
in (A22)--(A23) gives
\[
\begin{array}{c|rr}
 t&\#(G/J)^{\delta_t}&\#(G/J_-)^{\delta_t}\\ \hline
 12&125&0\\
 200&125&0\\
 \text{the other 22 values in (A21)}&0&0.
\end{array}
\]
The other $185$ rational base points have nontrivial cubic
Frobenius and contribute zero to both counts. Hence
\[
       \#(Y/J)(\mathbf F_{211})=250,
       \qquad \#(Y/J_-)(\mathbf F_{211})=0.                     \tag{A24}
\]
The integer-only program
\texttt{code/A/third\_derived\_certificate.py} verifies
every enumeration above, all nine pairs of cubic-root
choices for every row, and the coset formula by direct
subgroup membership. Its complete output is saved beside
it. The independently written
\texttt{code/F/A\_m3\_independent.py} gives the same two
counts using $\delta(Q)=3Q$ directly.

Suppose now that a lift as stated existed over $K_0$.
By Lemma~\ref{A:solvable-centers}, an outer class commuting
with $G_{K_0}$ induces an actual automorphism of the
arithmetic extension with geometric kernel $S_3$.
An inner correction makes its geometric projection to
$S_2$ equal to the specified representative $\phi_5$.
The finite quotient just constructed is preserved.
There is no undetermined arithmetic twist in this
statement. Indeed the geometrically connected constant
$G$-torsor supplies an epimorphism
$q_{\rm ar}:\pi_1(U_{0,K_0})\to G$ whose geometric
restriction is $q$. The pair $(q_{\rm ar},p)$ identifies
the quotient by $\ker(q|_F)$ with $G\times G_{K_0}$.
An automorphism of this product over $G_{K_0}$ restricting
to $\gamma$ on $G$ sends $(1,\sigma)$ to
$(z_\sigma,\sigma)$ with $z_\sigma\in Z(G)$. Since
$Z(G)=1$, it is exactly $\gamma\times\operatorname{id}$.
It therefore takes the arithmetic subgroup for $Y/J$
to that for $Y/J_-$.

Write $H=q^{-1}(J)$ and $H_-=q^{-1}(J_-)$. Because $G$
is metabelian, $F''\subset H,H_-$, and therefore
\[
                       F'''\subset H',H_-'.
\]
Abelianizing their images in $S_3$ consequently gives
the full groups $H^{\rm ab}$ and $H_-^{\rm ab}$.
The assumed arithmetic automorphism gives a
$G_{K_0}$-equivariant isomorphism between them, hence
between $H^1_{\rm et}$ of the two geometric covers with,
for example, $\mathbf Q_7$ coefficients.

Both covers have smooth marked reduction at $211$.
Their ramification is tame, its order divides $5625$,
and the three branch sections remain distinct; the
normalizations give smooth compactifications with etale
boundary after unramified extension. Smooth proper base
change together with the boundary sequence therefore
identifies the preceding representations with the
cohomology of their reductions. Poincare duality gives
isomorphic compactly supported $H^1$, while the compactly
supported $H^0$ vanishes and $H^2=\mathbf Q_7(-1)$ for
both geometrically connected affine curves. The trace
formula forces equal point counts, contradicting (A24).
This proves the assertion over $K_0$, and a fortiori
over $\mathbf Q$.

For the final assertion, the other coefficient primes of
$\phi_t$ have zero image in the derived subgroup of (A18).
On that quotient its action is $t_5$ on $V_1$ and one on
$V_2$, and the same argument defining $\gamma$ gives the
actual induced automorphism. If $t_5\not\equiv1\pmod5$,
some power of its residue in the cyclic group
$\mathbf F_5^\times$ equals $-1$. The same power of a
putative commuting lift therefore induces (A19) on
$V_1\times V_2$. Every step from (A20) to the contradiction
uses only this finite action. It applies verbatim to that
power and rules it out. Hence $t_5\equiv1\pmod5$.
\end{proof}

\begin{corollary}[Proved: a second coefficient-five digit is forced at derived length three]
\label{A:third-derived-second-digit}
For the unit family $\phi_t$ of
Proposition~\ref{A:metabelian-unit-family}, a lift to the
arithmetic commutant on $F/F'''$ over $\mathbf Q$, or over
$K_0=\mathbf Q(\mu_{15},E[5])$, necessarily satisfies
\[
                            t_5\equiv1\pmod{25}.
\]
This is the explicit second-digit certificate. The
structural Theorem~\ref{A:third-derived-IA} below
excludes every nonidentity member of the family at once;
the whole third-derived commutant remains undetermined.
\end{corollary}
\begin{proof}
The preceding proposition already forces
$t_5\equiv1\pmod5$. Replace its two covers by
\[
 E\xrightarrow{[25]}E\xrightarrow{(y+4)/8}\mathbf P^1,
 \qquad X^{25}+Y^{25}+Z^{25}=0,
\]
using $r=X^{25}/Z^{25}$ and the same rational functions
$w_2=(r-\zeta_3)/(r-\zeta_3^2)$ and $t=w_2^3/(w_2^3-1)$.
Over $K_{25}=\mathbf Q(\mu_{75},E[25])\supset K_0$ their fibre product
has constant metabelian deck group
\[
 G_{25}=
 \bigl(E[25]\rtimes C_3\bigr)
       \times\bigl((\mathbf Z/25)^2\rtimes C_3\bigr),
             \qquad |G_{25}|=3515625.
\]
For completeness, connectedness follows from the same
independent cubic subfields. Any nontrivial quotient
common to the two factors would be solvable and hence
have a nontrivial abelian quotient; the abelianization
of either factor is $C_3$. Thus such a common quotient
would identify their cubic subfields. The order-three
actions are fixed-point-free on both two-dimensional
$\mathbf Z/25$ modules, so $Z(G_{25})=1$.

The unit-family automorphism acts on the two derived
modules by $(t_5,1)$. The proof constructing $\gamma$
in the preceding proposition works over $\mathbf Z/25$:
$X-1$ remains invertible, and it gives an actual
automorphism of this finite quotient. All coefficient
primes other than five have zero image here.

The prime $q=5701$ is prime and splits completely in
$K_{25}$. Direct enumeration gives
\[
 \#E(\mathbf F_q)=5625=75^2,\qquad
 E[25]\subset E(\mathbf F_q),\qquad
 [75]E(\mathbf F_q)=0,
\]
and $q\equiv1\pmod{75}$. Frobenius has trace $77$ and
norm $5701$, hence is $1-75\zeta_3$ or
$1-75\zeta_3^2$. Accordingly the elliptic
$[25]$-Kummer class at $Q=(4w_1,8t-4)$ is conjugate
to $-3Q$. Choose $\zeta_3=75$ and $\zeta_{25}=303$
in $\mathbf F_q$. Reduction identifies the two
geometric torsion modules with the displayed finite
modules. Put
\[
 P_0=(25,2322)\in E[25],\qquad v_0=(0,23),\qquad
 J_a=\langle(aP_0,v_0)\rangle\quad(a\in(\mathbf Z/25)^\times).
\]
Both $P_0$ and $v_0$ have exact order $25$.
These intermediate covers have degree $140625$.

The finite point-count rule is exactly (A22)--(A23)
with $v$ taken modulo $25$, the exponent $42$
replaced by $(5701-1)/25=228$, the elliptic class
replaced by $-3Q$, and the coset multiplicity $125$
replaced by $25^4/25=15625$. Its output is
\[
\begin{array}{c|rrrrr}
 a&1&6&11&16&21\\ \hline
 \#(Y_{25}/J_a)(\mathbf F_{5701})&15625&0&0&0&15625.
\end{array}
\]
Only the fibre $t=8$ contributes to the first entry;
only $t=4104$ contributes to the last. There are $624$
base values with trivial image in both cubic quotients
and $5075$ other base values. The certificate
\texttt{code/A/third\_derived\_25\_certificate.py}
checks all $5699$ base values, all $5616$ pairs of
cubic-root choices in the relevant fibres, and all
$100$ generator choices for the five displayed
order-$25$ subgroups. Its output is saved with the code.
The independent program
\texttt{code/F/A\_m3\_25\_independent.py}, using $-3Q$
directly and separate subgroup-membership calculations,
gives the same first two counts and the unique fibre $t=8$.

If $t_5\not\equiv1\pmod{25}$, its residue, already in
$1+5\mathbf Z/25$, is a nonidentity element of a
cyclic group of order five. A suitable positive power
has residue $6$. The same power of a putative commuting
lift takes the subgroup $J_1$ to $J_6$. Restricting to
$G_{K_{25}}$ and using the center-free constant torsor
identifies the two arithmetic intermediate-cover
subgroups exactly. As before $F'''\subset H'$,
so their full first cohomology representations would
be isomorphic. Good tame reduction and the trace
formula would give equal point counts at $5701$,
contradicting the first two entries of the table.
This proves the congruence.
\end{proof}

\subsection{Fibre isolation eliminates the entire IA image at derived length three}

\begin{theorem}[Proved: the third-derived arithmetic commutant has no nontrivial metabelian IA image]
\label{A:third-derived-IA}
Let $K$ be any number field and
\[
 C_3(K)=C_{\operatorname{Out}(F/F''')}(G_K).
\]
An element of $C_3(K)$ inducing the identity on
$F^{\rm ab}$ induces an inner automorphism of $M=F/F''$.
Equivalently, the actual homomorphism
\[
 \operatorname{image}\bigl(C_3(K)\longrightarrow
                      \operatorname{Out}(M)\bigr)
       \longrightarrow\operatorname{GL}_2(\widehat{\mathbf Z})
                                                               \tag{A25}
\]
is injective. No peripheral hypothesis is needed here.
In particular the entire nonidentity unit family of
Proposition~\ref{A:metabelian-unit-family}, at every
coefficient prime and every depth, admits no commuting
third-derived lift over any number field.

The theorem does not assert that
$\operatorname{Out}(F/F''')\to\operatorname{Out}(M)$
has trivial kernel on the arithmetic commutant.
\end{theorem}
\begin{proof}
Put $A=M^{\rm ab}$, $R=\widehat{\mathbf Z}[[A]]$,
$I=\ker(R\xrightarrow\epsilon\widehat{\mathbf Z})$,
and $c=[x,y]$. We use $M'=Rc$ and the exact IA identity
(A4). The proof has six steps, including the arithmetic
normalizations and the point-count comparison.

\emph{1. An exactly invariant representative.}
Let the given third-derived class project to an IA class
of $M$, represented by a unit $u\in R^\times$ with
$\epsilon(u)=1$. There is a canonical continuous map
\[
 \partial:1+I\longrightarrow I/\overline{I^2}
       \simeq\widehat{\mathbf Z}^2=A,
       \qquad \partial(u)=u-1\pmod{\overline{I^2}}.
\]
The displayed identification sends $X-1,Y-1$ to the
two standard basis vectors. It follows by expanding
$X^aY^b$ modulo $I^2$: its image is
$1+a(X-1)+b(Y-1)$. This expansion extends continuously
from the integral group algebra. Conversely the ring
$\widehat{\mathbf Z}\oplus A$ with square-zero second
summand realizes these images, proving injectivity as
well as surjectivity. In particular $\partial$ is
additive on augmentation-one units, and the restriction
$A\to I/\overline{I^2}$ is the identity.

Arithmetic conjugation acts on IA outer classes by the
natural scalar action $\sigma_R$ induced by $\chi(\sigma)$
on $A$; its IA component commutes because IA outer
automorphisms form the abelian group in (A4).
Centralization therefore says
$u/\sigma_R(u)\in A$. Put $v=\partial u$ and multiply
$u$ by the monomial $a_{-v}\in A$. This is an allowed
inner change of representative. The resulting unit
$u_0$ has derivative zero and still satisfies
$u_0/\sigma_R(u_0)\in A$. That monomial has derivative
zero, so it is one. Hence
\[
       \sigma_R(u_0)=u_0\quad(\sigma\in G_K).                 \tag{A26}
\]
Choose an actual IA automorphism $\phi$ with this
determinant, and choose an inner adjustment of the
third-derived representative projecting exactly to
$\phi$. Such an adjustment lifts through the surjection
$F/F'''\to M$.

\emph{2. Exact identity on the arithmetic abelian quotient.}
Center-freeness supplies the arithmetic automorphism
$\beta$ of the extension with geometric kernel $F/F'''$.
Write $E_j=\pi_1(U_{0,K})/F^{(j)}$. We claim that its
induced automorphism of $E_1$ is the identity, not just
the identity on its geometric kernel.
For an arithmetic lift of $\sigma$, let $b_\sigma$ be
its conjugation action on $M$. On $Rc$, this action has
the form $r\mapsto\sigma_R(r)v_\sigma$ for a unit
$v_\sigma\in R^\times$. The automorphism $\phi$ acts
there by multiplication by $u_0$. Its commutator with
$b_\sigma$ thus acts by
$u_0/\sigma_R(u_0)=1$.
Since their outer classes commute, this commutator is
$\operatorname{Int}(d_\sigma)$ for some $d_\sigma\in M$.
Conjugation by $d_\sigma$ acts on $Rc$ by its monomial
$d_\sigma^{\rm ab}\in A\subset R^\times$. Faithfulness
of this monomial embedding gives $d_\sigma\in M'$.
Consequently $\beta$ changes each arithmetic lift only
by an element of $M'$, and it fixes the geometric
abelianization. This proves
\[
                         \beta|_{E_1}=\operatorname{id}.
                                                               \tag{A27}
\]

\emph{3. A finite quotient with different cyclic subgroups.}
Suppose the projected IA outer class is nontrivial.
Then $u_0\notin A$. The closed cyclic subgroups
$C=\overline{\langle c\rangle}$ and
$D=\overline{\langle\phi(c)\rangle}$ are not conjugate
in $M$. Indeed $C=\widehat{\mathbf Z}c$ and
$D=\widehat{\mathbf Z}u_0c$ inside $Rc$. A conjugator
acts by a monomial $a\in A$. Equality of these two
cyclic subgroups would give $u_0=\lambda a$ with
$\lambda\in\widehat{\mathbf Z}^{\times}$.
Applying augmentation gives $\lambda=1$, contrary to
$u_0\notin A$.

There is therefore a finite quotient of $M$ in which
their images are not conjugate. To justify this
finite detection, for every open normal subgroup $N$
consider the closed subset of $M$ of conjugators
identifying their images modulo $N$. If all these
sets were nonempty, their finite-intersection property
and compactness would give an actual conjugator in
$M$, a contradiction. Since $M$ is topologically
finitely generated, it has finitely many open subgroups
of any given index. Intersecting the finite $\phi$-orbit
of such an $N$ makes it $\phi$-stable without losing
the distinction. We obtain a finite metabelian quotient
\[
 q:M\twoheadrightarrow G,\qquad q\phi=\gamma q,
 \qquad J=\langle q(c)\rangle,
 \quad J'=\gamma(J)
\]
with $J,J'$ nonconjugate. For $g=q(c)$, the two
permutation characters satisfy
\[
                 \#(G/J)^g>0,\qquad \#(G/J')^g=0.           \tag{A28}
\]
The second equality follows because containment of the
generator $g$ in a conjugate of $J'$ would make the
cyclic groups $J,J'$ of equal order conjugate.

\emph{4. Arithmetic descent and removal of the finite twist.}
After a finite extension of $K$, the geometric quotient
is a geometrically connected constant $G$-torsor
$Y\to U_0$, with all deck transformations defined.
The map $(q_{\rm ar},p)$ identifies the corresponding
arithmetic quotient with $G\times G_K$.
The induced automorphism is of the form
\[
                 (h,\sigma)\longmapsto
                         (\gamma(h)z_\sigma,\sigma),
                 \qquad z_\sigma\in Z(G).
\]
Here $z:G_K\to Z(G)$ is a continuous homomorphism:
the $G$-factor is constant, and multiplication in the
product proves this formula directly. Replace $K$ by
the finite extension killing $z$, and then, if needed,
by its Galois closure over $\mathbf Q$. From now on
$K/\mathbf Q$ is Galois, (A27) still holds, and the
finite arithmetic action is exactly
$\gamma\times\operatorname{id}$. No center-free
assumption on $G$ is required in this step.

\emph{5. A Frobenius class and a cover isolating its fibre.}
Hilbert irreducibility gives $a\in U_0(K)$ such that
$Y_a=\operatorname{Spec}L$ is a field with
$\operatorname{Gal}(L/K)=G$; the precise specialization
statement is \cite[Proposition 3.3.1 and Theorem 3.4.1]{A_SerreTopics}.
Take the normal closure $L^*/\mathbf Q$ and lift
$g\in\operatorname{Gal}(L/K)$ to
$h\in\operatorname{Gal}(L^*/K)$.
Chebotarev applied to $L^*/\mathbf Q$
\cite[Chapter V, Theorem 3.23]{A_MilneCFT}
gives arbitrarily large rational primes $q$ and a prime
of $L^*$ with Frobenius $h$. Such a $q$ splits in $K$,
and at the resulting degree-one place $v$ the Frobenius
of $Y_a$ is $g$. Avoid the finitely many primes of bad
marked reduction, primes dividing $|G|$, and primes
at which $a$ or $1-a$ fails to be a unit. Write
$a_0\in\mathbf F_q\setminus\{0,1\}$ for its reduction.

Put $N=q-1$ and extend constants to $K'=K(\mu_N)$.
The chosen place has a degree-one extension to $K'$,
because $q\equiv1\pmod N$. Consider the actual finite
etale abelian cover
\[
 B:\qquad z_1^N=t/a,\qquad
               z_2^N=(1-t)/(1-a)                            \tag{A29}
\]
of $U_{0,K'}$. Its arithmetic monodromy factors through
$E_1$, so (A27) preserves this cover exactly, including
the two constant Kummer twists. Over $\mathbf F_q$,
the fibre $B_t$ has $N^2$ rational points at $t=a_0$
and no rational points at any other $t\in U_0(\mathbf F_q)$:
the $N$-th power of every nonzero element is one.

Form the two finite etale covers
\[
 Z=(Y/J)\times_{U_0}B,
       \qquad Z'=(Y/J')\times_{U_0}B.
\]
Their geometric monodromy is a subgroup of a product
of a metabelian group and an abelian group, hence is
metabelian. The two exact arithmetic actions already
proved make $\beta$ identify their finite-cover data.
Counting their rational points at the chosen place gives
\[
 \#Z(\mathbf F_q)=N^2\#(G/J)^g>0,
       \qquad \#Z'(\mathbf F_q)=N^2\#(G/J')^g=0.              \tag{A30}
\]

\emph{6. Why the third derived quotient forces equality.}
Neither fibre product is assumed geometrically connected.
For every geometric connected component its fundamental
group is an open subgroup $H\subset F$ containing $F''$,
because the entire geometric monodromy is metabelian.
Consequently $F'''\subset H'$ and the image of $H$ in
$F/F'''$ retains the full abelianization $H^{\rm ab}$.
The arithmetic automorphism $\beta$ gives an equivariant
bijection of the component sets and equivariant
isomorphisms of these abelianizations, with the conjugacy
identifications on stabilizers. It therefore identifies
both $H^0_{\rm et}$ and $H^1_{\rm et}$ of the geometric
curves $Z,Z'$ as $G_{K'}$-representations, including the
permutation of components.

All covers in this comparison have good marked reduction
at the chosen place: $Y$ was chosen good, and (A29) is
tame there, with constant units and branch sections only
at $0,1,\infty$. Finite products and the smooth marked
normalizations preserve this property: etale locally at
a branch section the covers have Kummer equations
$u^e=t$ with $e$ prime to the residue characteristic,
and their normalized products have the same form with
the least common multiple of the indices. Choose a
coefficient prime different from $q$. Base change and
Poincare duality identify the compactly supported first
cohomology representations as well. Compactly supported
$H^0$ vanishes, and $H^2_c$ is the component permutation
module twisted by $-1$ for both curves. Thus the trace
formula makes the two point counts equal, contradicting
(A30).

The projected IA class is therefore trivial. Applying
this to the quotient of any two elements having the
same abelian action proves the injectivity of (A25).
The unit family acts identically on $A$ and is injective
in IA outer automorphisms, so every nonidentity member
is excluded, at all primes and to all depths. The proof
does not compare different automorphisms inside the
kernel of the third-to-second-derived projection.
\end{proof}

\begin{proposition}[Proved: the peripheral metabelian commutant has no extra scalar multiplier]
\label{A:metabelian-multiplier}
For every number field $K$, the image of
$C_{\operatorname{Out}_{\rm lac}(M)}(G_K)$ on
$M^{\rm ab}$ is exactly the geometric $S_3$.
In particular a pure peripheral outer automorphism
commuting with $G_K$ has common multiplier one.
This statement leaves the nontrivial IA commutant of
Proposition~\ref{A:metabelian-unit-family} present at
the metabelian level.
\end{proposition}
\begin{proof}
Compose with a geometric permutation of the three cusps
to make the class pure. By
Proposition~\ref{A:abelian-commutant}, its abelian action
is $\rho\operatorname{id}$ for a unit
$\rho\in\widehat{\mathbf Z}^{\times}$. Center-freeness
gives an arithmetic automorphism $\beta$ of $E_2$.
We first normalize its action on $E_1$ exactly.

The geometric normalizer of $\langle x\rangle$ in $M$
is $\langle x\rangle$. Here is a proof with all
profinite coefficients retained. If an element commutes
with $x$, the free class-two quotient shows that its
$y$-coordinate in $M^{\rm ab}$ is zero. It is therefore
$x^a c^r$ for $a\in\widehat{\mathbf Z}$, $r\in R$,
and commutation gives $(X-1)r=0$. Multiplication by
$X-1$ is injective on $R$. Indeed in a coefficient
factor $\mathbf Z_p[[A]]$, an $X$-invariant element has
constant coefficients on the finite $X$-orbits in
every finite abelian group quotient. Refine the order
of the first cyclic factor by $p^j$ while retaining
the second factor. Projection sums $p^j$ equal
coefficients. Modulo $p^j$ the projected element is
zero. This holds for every coefficient modulus and
group quotient, so the element is zero. Thus
$C_M(x)=\langle x\rangle$. An element normalizing
$\langle x\rangle$ induces exponent one on it by
abelianization, and hence centralizes it. The same
argument applies to $y$.

Let $s_0,s_1$ be the tangential sections at $0$ pointing
towards $1$, and at $1$ pointing towards $0$,
respectively. Their images in $E_1$ are the same
section. This can be checked on the actual abelian
Kummer covers $t^{1/n},(1-t)^{1/n}$: the parameter
has leading coefficient one at its own cusp and the
other function has value one. All these leading
coefficients admit the compatible rational roots one.
Choose a path whose abelian image identifies these
root labels. Both sections then have abelian Kummer
coordinates $(0,0)$. This identifies
\[
              E_1=A\rtimes G_K.
\]
The arithmetic normalizers are exactly
$\langle x\rangle\rtimes s_0(G_K)$ and
$\langle y\rangle\rtimes s_1(G_K)$: subtract the
specified tangential lift of an arithmetic element
and use the geometric normalizer equality above.

Write the induced automorphism of $E_1$ as
\[
               (v,g)\longmapsto(\rho v+c(g),g).
\]
Peripheral preservation takes each of those two
arithmetic normalizers to a geometric conjugate of
itself. If the conjugator at cusp $i$ has abelian
coordinate $a_i$, comparison of the second coordinate
at cusp $0$ and of the first coordinate at cusp $1$
gives
\[
 c_2(g)=(1-\chi(g))(a_0)_2,
       \qquad c_1(g)=(1-\chi(g))(a_1)_1.
\]
Thus $c$ is the actual coboundary
$(1-\chi)a$, with $a=((a_1)_1,(a_0)_2)$.
An inner adjustment with abelian coordinate $-a$
now makes
\[
                   \beta|_{E_1}(v,g)=(\rho v,g).             \tag{A31}
\]

Take any odd rational prime $q$ splitting completely
in a Galois closure of $K$, put $N=q-1$, and extend
constants to $K'=K(\mu_N)$. There is a degree-one
place of $K'$ over $q$, since $q\equiv1\pmod N$.
For $a\in\{2,\ldots,q-1\}$ consider
\[
 B_a:\quad z_1^N=t/a,\qquad z_2^N=(1-t)/(1-a).
\]
Relative to the split coordinates (A31), its arithmetic
quotient is $(v,g)\mapsto v-(\kappa_N(a)(g),
\kappa_N(1-a)(g))$. If $r$ is the residue of $\rho$
modulo $N$, equation (A31) therefore sends its
arithmetic cover subgroup to that of
\[
 B'_a:\quad w_1^N=t/a^r,\qquad
                  w_2^N=(1-t)/(1-a)^r.                      \tag{A32}
\]
This verifies the constant twists as well as the
geometric monodromy; it is not merely a comparison
of untwisted abelian covers.

The monodromy of either cover is geometrically abelian.
Every geometric component subgroup $H$ thus contains
$F'$, and $F''\subset H'$. Already $E_2$ retains
its full first homology. Hence $\beta$ identifies the
$H^0$ and $H^1$ Galois representations of $B_a,B'_a$.
Both have good marked reduction at $q$: their
projective equations are diagonal Fermat equations
of degree $N$, with unit coefficients and with
$N$ invertible in the residue field. Their boundaries
are etale there. The same duality and trace-formula
argument as in the preceding theorem forces equal
affine point counts.

The first cover has exactly $N^2$ rational points,
all over $t=a$. The second has $N^2$ rational points
if
\[
                     a^r+(1-a)^r=1\quad\text{in }\mathbf F_q,
                                                               \tag{A33}
\]
and zero otherwise. Thus (A33) holds for every
$a\in\mathbf F_q\setminus\{0,1\}$.
The map $f(0)=0$, $f(x)=x^r$ on $\mathbf F_q^\times$
is multiplicative and bijective, since $r$ is a unit
modulo $q-1$. Equation (A33), applied to
$a=x/(x+y)$, proves additivity whenever
$x,y,x+y$ are nonzero. The zero cases are immediate;
when $x+y=0$ use $f(-1)=-1$, since $r$ is odd.
Therefore $f$ is an automorphism of the prime field
$\mathbf F_q$ and is the identity. It follows that
\[
                         \rho\equiv1\pmod{q-1}.
\]
For every positive integer $n$, Chebotarev supplies
such an odd prime $q$ also splitting in $\mathbf Q(\mu_n)$.
Then $n\mid q-1$, so $\rho\equiv1\pmod n$.
This proves $\rho=1$ in $\widehat{\mathbf Z}$.
Undoing the initial geometric permutation gives exactly
the displayed $S_3$ image.
\end{proof}

\begin{corollary}[Proved: the third-to-second-derived peripheral commutant is exactly geometric]
\label{A:third-derived-projection}
For every number field $K$,
\[
 \operatorname{image}\left(
 C_{\operatorname{Out}_{\rm lac}(F/F''')}(G_K)
          \longrightarrow\operatorname{Out}(F/F'')\right)=S_3.
\]
The same equality holds for the projection from every
derived length $m\ge3$ to length two. In particular
all nongeometric metabelian peripheral commutant classes
are obstructed at the next derived step.
The same-level commutant at length three is reduced
to its kernel over this geometric image; that kernel
has not been proved trivial here.
\end{corollary}
\begin{proof}
By Proposition~\ref{A:metabelian-multiplier}, the
abelian image of a projected class is that of one
geometric permutation. Divide by that permutation.
The quotient has trivial abelian action, and
Theorem~\ref{A:third-derived-IA} makes its metabelian
outer image trivial. Conversely all six geometric
permutations lift over $K$ at every derived level.
For $m>3$ first project to length three. This gives
the stated equality and does not assert anything
further about the kernel of that projection.
\end{proof}

\subsection{An explicit unrestricted last-layer construction and its peripheral obstruction}

\begin{proposition}[Proved: unrestricted third-derived scalar classes fail the peripheral test]
\label{A:last-layer-unrestricted}
For every odd prime $\ell$, the kernel
\[
 \ker\bigl(\operatorname{Out}(F/F''')\to\operatorname{Out}(M)\bigr)
       \cap C_{\operatorname{Out}(F/F''')}(G_{\mathbf Q})
\]
contains a closed copy of $\mathbf Z_\ell^\times$.
No nonidentity member of this copy preserves the
conjugacy class of the cusp subgroup $\langle x\rangle$.
Thus this explicit construction does not give an element
of the peripheral kernel left open above.
\end{proposition}
\begin{proof}
Write $L=F''/F'''$, $D=M'$, and
$\Lambda=\widehat{\mathbf Z}[[M]]$.
We use the profinite Fox relation-module sequence,
with its actual boundary map,
\[
 0\longrightarrow L\longrightarrow\Lambda^2
   \xrightarrow{\partial}\Lambda
   \xrightarrow\epsilon\widehat{\mathbf Z}\longrightarrow0,
 \qquad \partial(a,b)=a(x-1)+b(y-1).                         \tag{A34}
\]
This is the Blanchfield--Lyndon sequence, recalled
with precisely this formula in
\cite[Proposition 1.1]{A_Yamaguchi2026Centers} from
Ihara's profinite Fox calculus. More explicitly the
continuous Fox derivation $\delta:F\to\Lambda^2$
satisfies $\delta(x)=(1,0)$, $\delta(y)=(0,1)$,
$\delta(fg)=\delta(f)+\bar f\delta(g)$, and
$\partial\delta(f)=\bar f-1$.
The resulting Magnus map identifies $F/F'''$ with
\[
 \{(v,g)\in\Lambda^2\rtimes M:\partial v=g-1\}.             \tag{A35}
\]
Indeed its kernel is $F'''$ by (A34). For a prescribed
$g$, subtract the Fox vector of any lift of $g$;
(A34) says that every other vector with the same
boundary differs by an element of $L$, proving the
asserted equality of the images.

Let $D^{(\ell')}$ be the prime-to-$\ell$ factor of
the abelian group $D=Rc$. Reduction modulo two of
the augmentation $Rc\to\widehat{\mathbf Z}$ gives
an $M$-invariant surjection
$\lambda:D^{(\ell')}\to C_2$. Put $H=\ker\lambda$.
Normalized averaging over this pro-prime-to-$\ell$
group defines $e_H\in\mathbf Z_\ell[[M]]$.
For $z\in D^{(\ell')}$ with $\lambda(z)\ne1$, set
\[
                   e=e_H(1-z)/2.
\]
Regard $e$ as zero in the other coefficient factors
of $\Lambda$. It is a central idempotent of augmentation
zero: $z^2\in H$, and conjugation by $M$ fixes the
character $\lambda$. It is also arithmetic-invariant.
Indeed arithmetic acts on the augmentation of $D$
by $\chi^2$, which is one modulo two. Every finite
averaging denominator is an $\ell$-adic unit; no
averaging over a pro-$\ell$ kernel is used.

Choose explicitly $z=c^{\eta}$, where
$\eta\in\widehat{\mathbf Z}$ is zero at $\ell$
and one at every other prime, and take the same
word $\widetilde z=[x,y]^\eta$ in $F$. Since $ez=-e$,
\[
                    w=-\tfrac12e\delta(\widetilde z)
          \quad\text{satisfies}\quad \partial w=e.
\]
For $t\in\mathbf Z_\ell^\times$, put
$u_t=1+(t-1)e$ and define on (A35)
\[
 \Phi_t(v,g)=
 \bigl(u_tv+(1-t)(g-1)w,\ g\bigr).                         \tag{A36}
\]
Centrality of $e$ and the crossed-homomorphism identity
for $(g-1)w$ prove directly that this is a homomorphism
preserving the boundary condition. Also
$\Phi_t\Phi_s=\Phi_{ts}$ and
$\Phi_t^{-1}=\Phi_{t^{-1}}$. Thus these are actual
continuous automorphisms, identity on $M$, acting by
$t$ on $eL$ and by one on $(1-e)L$.

We check arithmetic commutation in the outer group.
First
\[
                         H^1(M,eL)=0.                       \tag{A37}
\]
To see this directly, a continuous cocycle restricts
on $H$ to a homomorphism from a pro-prime-to-$\ell$
group to the pro-$\ell$ group $eL$, hence is zero.
It descends to $M/H$, where $z$ is a central
involution acting by $-1$. If $b$ is such a cocycle,
the two expressions for $b(zg)=b(gz)$ give
$2b(g)=(1-g)b(z)$. This is the principal cocycle
with parameter $b(z)/2$, proving (A37).

Arithmetic preserves $eL$ and commutes with its scalar
$t$. The commutator of $\Phi_t$ with any arithmetic
automorphism therefore acts identically on $M$ and
on $L$, and is identity modulo $eL$. Such an
automorphism is a continuous cocycle $M\to eL$:
compare the image of any lift of $g\in M$ with that
lift; independence follows because $L$ is fixed.
By (A37) it is inner. Hence $[\Phi_t]$ commutes
with the arithmetic outer action.

Finally we give an exact detector of both nontriviality
and failure at the cusp. Put
\[
                    n_x=e(1,0)-(x-1)w\in eL.
\]
Map $M$ to the dihedral group of order eight using
the integral matrices
\[
 x\longmapsto\begin{pmatrix}1&0\\0&-1\end{pmatrix},
 \qquad
 y\longmapsto\begin{pmatrix}0&-1\\1&0\end{pmatrix}.
\]
The commutator maps to $-1$, so $H$ maps to the
identity and $e$ maps to the identity matrix.
The first row is fixed by $x$. Apply this row to
the first coordinate of $\Lambda^2$ in the
coefficient-$\ell$ factor. The resulting continuous
linear map kills $(x-1)eL$ and takes $n_x$ to
$(1,0)\in\mathbf Z_\ell^2$. In particular
\[
          (t-1)n_x\notin(x-1)eL\quad(t\ne1).               \tag{A38}
\]
The automorphism $\Phi_t$ changes $x$ by precisely
$(t-1)n_x$. If its image were conjugate to $x$, the
conjugator's image in $M$ would centralize $x$ and
thus be a power of $x$, by the proof of
Proposition~\ref{A:metabelian-multiplier}. Multiplying
the conjugator by the inverse power of $x$ reduces
it to $L$. The conjugation difference would then
lie in $(x-1)L$, contradicting (A38) after applying
$e$. A conjugate generator of $\langle x\rangle$
cannot use any other exponent, since $\Phi_t$ is
identity on abelianization. Thus it fails even this
single peripheral condition.

The same row detector shows that $t$ acts nontrivially
on $eL$ for $t\ne1$. An inner automorphism inducing
the identity on the center-free group $M$ has its
conjugator in $L$, and acts identically on $L$.
Therefore $t\mapsto[\Phi_t]$ is injective.
Continuity follows from (A36); compactness gives
a closed image.
\end{proof}

\begin{proposition}[Proved: no normal prime-to-prime averaging block removes a peripheral divisor]
\label{A:haar-block-obstruction}
Let $\ell$ be any prime, let $N$ be a closed normal
pro-prime-to-$\ell$ subgroup of $M$, and put
$\Lambda_\ell=\mathbf Z_\ell[[M]]$.
Then $N\subset M'$. The natural map
\[
 \mathbf Z_\ell[[N]]\longrightarrow
 \Lambda_\ell/(x-1)\Lambda_\ell
       =\mathbf Z_\ell[[\langle x\rangle\backslash M]]       \tag{A39}
\]
is injective. In particular, for any nonzero
idempotent $e\in\mathbf Z_\ell[[N]]$ whose image is
central in $\Lambda_\ell$, multiplication
by $x-1$ is not surjective on $e\Lambda_\ell$.
It cannot be inverted there. The same holds for
$y$ and $xy$.

Consequently a scalar construction using Haar
idempotents of normal pro-prime-to-$\ell$ subgroups
cannot reproduce the simultaneous peripheral
inversions used in the metabelian construction.
This is an obstruction to that specified construction,
not a determination of the remaining peripheral kernel.
\end{proposition}
\begin{proof}
Let $D_\ell$ be the pro-$\ell$ factor of $M'$.
Both $N$ and $D_\ell$ are normal, so
$[N,D_\ell]\subset N\cap D_\ell=1$.
The action of $M$ on
$D_\ell=\mathbf Z_\ell[[A]]c$ is its faithful
monomial action through $A$. Hence every element
of $N$ has trivial image in $A$, proving
$N\subset M'$.

Since the abelianization embeds
$\langle x\rangle$ into $A$, one has
$N\cap\langle x\rangle=1$. The map of profinite
sets $N\to\langle x\rangle\backslash M$ is therefore
a continuous injection, hence a closed embedding.
The induced map of free profinite
$\mathbf Z_\ell$-modules is injective: a finite
clopen partition of a closed subset of a profinite
space extends to a finite clopen partition of the
whole space. Thus every nonzero finite coefficient
of a measure on $N$ remains detectable in the
ambient permutation module. This proves (A39).
The equality on the right of (A39) is precisely
the coinvariant quotient for left multiplication
by the procyclic group $\langle x\rangle$;
the image of $x-1$ is closed by compactness.

In particular the image of $e$ in that quotient
is nonzero. If $x-1$ were surjective on
$e\Lambda_\ell$, one could write
$e=(x-1)v$, giving zero image, a contradiction.
Replace $x$ by $y$ or $xy$ to obtain the other
assertions: each displayed cyclic subgroup
injects into $A$ by its primitive abelian vector.
\end{proof}

\subsection{The remaining normalizer cannot admit a virtual arithmetic section}

\begin{proposition}[Proved using the supplied local theorem: virtually split procyclic normalizers are cuspidal]
\label{A:split-normalizer}
Let $U/K$ be an affine hyperbolic curve over a number
field, with arithmetic group $E$ and geometric group
$\pi$. Let $A\subset\pi$ be a nontrivial closed
procyclic subgroup whose arithmetic normalizer has
open image. Then the following conditions are equivalent:
\begin{enumerate}
\item $A$ is contained in a cusp inertia group;
\item for some finite extension $L/K$, the map
\[
 N_E(A)\cap p^{-1}(G_L)\longrightarrow G_L
\]
has a continuous section.
\end{enumerate}
Consequently every noncuspidal normalizer remaining
in Propositions~\ref{A:full-support-normalizer} and
\ref{A:cm-normalizer-field} fails to admit a section
after every finite arithmetic base extension.
The implication from the second condition uses the
supplied local section theorem through
Proposition~\ref{F-local-dichotomy}; its stated
unrefereed status is retained.
\end{proposition}
\begin{proof}
For the first implication, pass to a finite extension
making the cusp rational. A nonzero rational tangent
gives a section of its decomposition group. The
cyclotomic action preserves every closed subgroup
of its procyclic inertia, so this section takes
values in $N_E(A)$.

Conversely, replace $K$ by the finite extension
in the second condition and write
$s:G_K\to N_E(A)$ for its actual section.
Choose a prime $\ell$ with $A_\ell\ne1$.
The torsion-freeness of the free profinite group
$\pi$ gives $A_\ell\simeq\mathbf Z_\ell$.
Choose a characteristic open subgroup $V\subset\pi$
such that a generator of $A_\ell$ has nontrivial
image in $\pi/V$, and set $H=V A_\ell$.
Then $H$ is open and is normalized by $s(G_K)$.
The actual open arithmetic subgroup
\[
                         Hs(G_K)\subset E
\]
therefore defines a geometrically connected finite
\'etale cover $Y\to U$ over $K$. Its geometric
group is $H$. The cyclic quotient $H/V$ detects
$A_\ell$ in $H^{\mathrm{ab}}$.
Since $H^{\mathrm{ab}}$ is torsion-free, the resulting
map
\[
 A_\ell\longrightarrow H^{\mathrm{ab}}_{\ell}
       =H_1(Y_{\overline K},\mathbf Z_\ell)
                                                        \tag{A40}
\]
is injective. It is equivariant for the precise
action defined by $s$: conjugation by this section
is the ordinary arithmetic action on the
abelianization of the cover.

Let $\overline Y$ be the smooth compactification
and $J=\operatorname{Jac}(\overline Y)$. The
boundary exact sequence is
\[
 0\longrightarrow W\longrightarrow H_1(Y_{\overline K},\mathbf Z_\ell)
       \longrightarrow T_\ell J\longrightarrow0.
                                                        \tag{A41}
\]
If the composite of (A40) with $T_\ell J$ is zero,
then after splitting the finitely many cusps over
a finite extension its image lies in a direct
sum of copies of $\mathbf Z_\ell(1)$. The section
$s$ thus acts cyclotomically on $A_\ell$.
Proposition~\ref{A:cyclotomic-normalizer} puts
$A_\ell$ in one cusp inertia $I$.
Every element of $A$ centralizes every element
of $A_\ell$, and the full geometric centralizer
of any nontrivial element of $I$ is $I$.
Hence $A\subset I$.

It remains to treat a nonzero composite
$i:A_\ell\hookrightarrow T_\ell J$.
It is injective because $T_\ell J$ is torsion-free.
Faltings's semisimplicity theorem over this number
field, in its exact rational Tate-module form
\cite[\S5, Satz 3]{Faltings1983}, supplies a
$G_K$-equivariant retraction of
$i\otimes\mathbf Q_\ell$. Clearing denominators
in the two finite free lattices gives a nonzero
continuous $G_K$-homomorphism
\[
                 r:T_\ell J\longrightarrow A_\ell.
                                                        \tag{A42}
\]
The action on the target of (A42) is exactly
conjugation by $s$, by (A40).

Use the zero point of $J$ to identify its actual
arithmetic extension with
\[
 E_J=\left(\prod_qT_qJ\right)\rtimes G_K.
\]
The formula
\[
 f:E_J\longrightarrow E,
       \qquad f(v,g)=r(v_\ell)s(g)                       \tag{A43}
\]
is a continuous arithmetic homomorphism: (A42)
is precisely the equivariance needed for its
multiplication law. Its geometric image is a
nonzero closed subgroup of $A_\ell$, hence is
procyclic. It is not open in the nonabelian
geometric group of a hyperbolic curve.
Proposition~\ref{F-local-dichotomy}, whose source
may be any smooth geometrically connected
variety, puts this image in one cusp inertia
$I\subset\pi$. Finally $A$ centralizes that
nonzero image. The same full inertia centralizer
identity gives $A\subset I$.
This proves the second implication and the
stated obstruction to every virtual splitting.
\end{proof}

\begin{proposition}[Proved: the remaining normalizer section problem has prime-to-prime geometric kernel]
\label{A:normalizer-prime-complement}
Keep the number-field and affine hypotheses above.
Let $A\subset\pi$ be noncuspidal, procyclic and
nontrivial, with open arithmetic normalizer image,
and choose $\ell$ with $A_\ell\ne1$.
There exist an open subgroup $N_0\subset N_E(A)$,
a finite extension $L/K$ with $p(N_0)=G_L$,
and closed subgroups $P,C,B\subset N_0$ such that
\[
 \begin{gathered}
 N_0\cap\pi=P\times C,\qquad
 1\ne P\subset A_\ell,\quad P\simeq\mathbf Z_\ell,
       \quad C\text{ is pro-prime-to-}\ell,\\
 N_0=P\rtimes B,\qquad
 1\longrightarrow C\longrightarrow B
       \xrightarrow{p}G_L\longrightarrow1.
 \end{gathered}                                             \tag{A44}
\]
In particular the quotient $N_0/C$ has an actual
arithmetic section. The still required section of
$N_0\to G_L$ is the explicit lifting problem
through the last extension in (A44).
Using Proposition~\ref{A:split-normalizer}, that
extension has no section after any finite extension
of $L$; in particular $C\ne1$.
The structure and quotient splitting in (A44) use
Faltings semisimplicity and the cyclotomic recognition
theorem. Only the final non-splitting conclusion
also uses the supplied local section theorem.
\end{proposition}
\begin{proof}
Initially write $N=N_E(A)$ and $D=N\cap\pi$.
The action $\theta:N\to\mathbf Z_\ell^\times$
on $A_\ell$ has finite prime-to-$\ell$ image
on $D$, by Lemma~\ref{A:normalizer-character-descent}.
Choose an open subgroup of
$\mathbf Z_\ell^\times$ meeting that finite image
only in one, and replace $N$ by its inverse
image. Then $D$ centralizes $A_\ell$.
Let $P^*$ be a pro-$\ell$ Sylow subgroup of $D$.
It contains the normal subgroup $A_\ell$.
As a closed pro-$\ell$ subgroup of the free
profinite group $\pi$, it is free pro-$\ell$.
A noncyclic free pro-$\ell$ group has trivial
center, whereas $A_\ell$ is central in $P^*$.
Thus $P^*\simeq\mathbf Z_\ell$ and $A_\ell$
is open in $P^*$.

Here is the precise shrinking needed to make the
Sylow subgroup normal; torsion-freeness alone
does not identify the different Sylow subgroups.
The compact set $P^*\setminus A_\ell$ avoids
the identity. Choose an open normal subgroup
$N_1\triangleleft N$ disjoint from this set.
Put $D_1=N_1\cap D$ and $P_1=P^*\cap D_1$.
Then $D_1$ is normal in $D$, $P_1$ is a Sylow
subgroup of $D_1$, and
\[
                 P_1\subset A_\ell\subset Z(D).
\]
Conjugacy of Sylow subgroups now makes $P_1$
the unique Sylow subgroup of $D_1$.
It is nontrivial, because it is open in $P^*$.
The normal central Sylow subgroup has a unique
prime-to-$\ell$ complement $C_1$, and
$D_1=P_1\times C_1$.
For completeness this splitting can be checked
at every finite quotient: the central
$\ell$-subgroup extension of a group of order
prime to $\ell$ splits by the averaging
calculation of $H^2$, and the complement is
unique since a homomorphism from a prime-to-$\ell$
group to an $\ell$-group is zero. These unique
complements are compatible under the finite
quotient maps. Their inverse limit is $C_1$.
Uniqueness also makes $C_1$ characteristic in
$D_1$ and normal in $N_1$.

Choose a characteristic open $V\subset\pi$
detecting a generator of $P_1$, and set
$H=V P_1$. Descend its finite geometric cover
to a finite number-field extension and choose
its actual arithmetic subgroup $E_Y\subset E$,
with $E_Y\cap\pi=H$. Intersect $N_1$ with
$E_Y$. Its geometric kernel is
\[
 (P_1\times C_1)\cap H
       =(P_1\cap H)\times(C_1\cap H).
                                                        \tag{A45}
\]
Indeed closed subgroups of a product of a
pro-$\ell$ and a pro-prime-to-$\ell$ group
split into their two factors: the two primary
components of each element are limits of its
integer powers. The same assertion holds
after every subsequent finite arithmetic
restriction. We may therefore keep the notation
$N,D=P\times C$ for this intersection, with
$P\subset P_1$ nontrivial. Its arithmetic
image is still open; replace the field by
that image's fixed field.

The detector construction gives a nonzero,
hence injective, map
$P\to H_1(Y_{\overline K},\mathbf Z_\ell)$.
Its arithmetic action is the action of $N$:
the kernel $D\subset H$ acts trivially on this
abelianization, so the action factors through
the displayed arithmetic quotient.
If its image in the compactified Tate module
were zero, after splitting the boundary it
would be cyclotomic. The third condition of
Proposition~\ref{A:cyclotomic-normalizer} would
then put $P$ in cusp inertia, and its full
centralizer would put all of $A$ there.
This contradicts the assumption. Consequently
\[
             i:P\hookrightarrow T_\ell J,
       \qquad J=\operatorname{Jac}(\overline Y),
                                                        \tag{A46}
\]
is injective and equivariant for the actual
quotient action of $G_L$ on $P$.

Pass to a further finite extension so that
$\overline Y$ has a rational point $y_0$.
The Abel--Jacobi morphism
$Y\to J$, $y\mapsto[y-y_0]$, gives an actual
arithmetic group homomorphism
\[
                  E_Y\longrightarrow
                  \left(\prod_qT_qJ\right)\rtimes G_L.
\]
Faltings semisimplicity and clearing denominators
give an equivariant $r:T_\ell J\to P$ with
$r i=\ell^k\operatorname{id}_P$ for some $k\ge0$.
After projecting to $T_\ell J$ and applying $r$,
the preceding arithmetic map restricts to
\[
                  \varphi:N\longrightarrow P\rtimes G_L.
                                                        \tag{A47}
\]
It kills $C$, since $C$ is pro-prime-to-$\ell$
and its geometric image is pro-$\ell$. It
therefore induces an injective homomorphism
\[
             \overline\varphi:N/C\hookrightarrow P\rtimes G_L:
                                                        \tag{A48}
\]
any element of its kernel lies in $D/C=P$,
where its map is multiplication by $\ell^k$.
This multiplication is injective.

The image $Q$ of (A48) projects onto $G_L$
and meets the geometric $P$ in $\ell^kP$.
Modulo $\ell^kP$ it is the graph of a
continuous finite-valued cocycle
$b:G_L\to P/\ell^kP$. On the open subgroup
where $b$ vanishes, $Q$ contains $(0,g)$
for every $g$: subtract the geometric
coordinate in $\ell^kP$ from any lift.
Taking an open normal subgroup inside this
kernel and replacing $L$ accordingly gives
an actual continuous section of $N/C$,
by the inverse of the homeomorphism (A48)
onto its compact image.

Let $B$ be the inverse image of this section
in $N$. Then $B\cap\pi=C$, $B\to G_L$
is onto, $B\cap P=1$, and $N=PB$.
As $P$ is normal, multiplication identifies
$N$ with $P\rtimes B$, proving (A44).
A section of $B\to G_L$ would be a section
of $N_E(A)\to G_L$, contradicting
Proposition~\ref{A:split-normalizer}; the same
argument applies after every finite extension.
\end{proof}

\begin{proposition}[Refuted: a prime-to-prime kernel alone does not imply virtual splitting]
\label{A:prime-complement-counterexample}
For every prime $\ell$ there is a profinite extension
\[
 1\longrightarrow\mathbf Z_q(1)\longrightarrow
 \mathcal B\longrightarrow G_{\mathbf Q}
       \longrightarrow1,
       \qquad q\ne\ell,
                                                        \tag{A49}
\]
whose kernel is pro-prime-to-$\ell$ and which has
no continuous section after any finite extension
of $\mathbf Q$. Thus the last extension in (A44)
cannot be split merely from the prime support
of its kernel. No embedding of (A49) into an
arithmetic curve group is asserted.
\end{proposition}
\begin{proof}
Choose an odd prime $q\ne\ell$ and two distinct
finite places $v,w$ of $\mathbf Q$.
For each $n\ge1$, the Brauer--Hasse--Noether
sequence gives a unique Brauer class with local
invariants
\[
 \operatorname{inv}_v(\beta_n)=q^{-n},\qquad
 \operatorname{inv}_w(\beta_n)=-q^{-n},\qquad
 \operatorname{inv}_u(\beta_n)=0\quad(u\ne v,w).
                                                        \tag{A50}
\]
The finite-coefficient Kummer sequence and Hilbert
90 identify this class with
$\beta_n\in H^2(G_{\mathbf Q},\mu_{q^n})$.
The coefficient transition
$\mu_{q^{n+1}}\to\mu_{q^n}$ induces multiplication
by $q$ on the Brauer group. Hence
$q\beta_{n+1}=\beta_n$ under these identifications.
The exact global sequence and the local
restriction formula used here are
\cite[IV, 2.14(e); III, (29)]{A_MilneCFT}.

We construct compatible actual continuous
cocycles, so that no finite-to-continuous
existence assertion is left implicit.
Suppose a normalized continuous two-cocycle
$z_n$ representing $\beta_n$ has been chosen.
Choose any such representative $z'_{n+1}$
of $\beta_{n+1}$. Its reduction differs from
$z_n$ by $\delta u_n$ for a continuous
one-cochain $u_n$ with values in $\mu_{q^n}$.
This finite-valued locally constant cochain
has a continuous lift $\widetilde u_n$ to
$\mu_{q^{n+1}}$, obtained by choosing a lift
of each of its finitely many values, and choosing
zero as the lift of zero. Since the discrepancy
is normalized, $u_n(1)=0$, so normalization
is preserved.
Replace $z'_{n+1}$ by
$z_{n+1}=z'_{n+1}-\delta\widetilde u_n$.
These representatives are now compatible.
Their inverse limit is a normalized continuous
two-cocycle
\[
 z:G_{\mathbf Q}^2\longrightarrow\mathbf Z_q(1).
\]
Continuity is exactly continuity of every
finite coordinate in the inverse-limit
topology. This construction lifts the specified
compatible nonzero obstruction classes; it does
not assert that vanishing finite obstruction
classes force vanishing of a continuous one.

Give $\mathbf Z_q(1)\times G_{\mathbf Q}$ the
product topology and multiplication
\[
 (a,g)(b,h)=\bigl(a+g b+z(g,h),gh\bigr).
                                                        \tag{A51}
\]
The cocycle identity proves associativity;
normalization gives the identity, and the
inverse is obtained by solving the displayed
formula. This is a compact totally disconnected
topological group, hence a profinite group,
and it gives (A49).

For a finite extension $L/\mathbf Q$ and a
place $v'\mid v$, local restriction multiplies
the invariant in (A50) by the positive integer
$[L_{v'}:\mathbf Q_v]$. For sufficiently large
$n$ the result
\[
                 [L_{v'}:\mathbf Q_v]/q^n
                         \quad\text{in }\mathbf Q/\mathbf Z
\]
is nonzero. Thus $\beta_n|_{G_L}\ne0$.
A section of (A49) over $G_L$ would make
$z|_{G_L}$ a continuous coboundary, and every
one of its finite reductions would then be
a coboundary. This contradiction proves
failure of splitting over every finite $L$.
\end{proof}

\subsection{The corresponding restrictions for proper hyperbolic curves}

\begin{lemma}[Known with independent proof: finitely generated pro-prime subgroups of a full surface group]
\label{A:surface-proprime-free}
Let $\Pi$ be the full profinite fundamental group
of a proper hyperbolic curve of genus $g\ge2$
over an algebraically closed field of characteristic
zero. Every finitely generated closed pro-$\ell$
subgroup $P\subset\Pi$ is free pro-$\ell$.
In particular, every finitely generated pro-$\ell$
subgroup having a nontrivial procyclic normal
subgroup is procyclic.
\end{lemma}
\begin{proof}
We give a direct surface-cohomology proof; the
broader pro-prime freeness theorem for profinite
completions of torsion-free hyperbolic virtually
special groups is \cite[Theorem F]{A_WiltonZalesskii2017}.
For a closed subgroup $P$ and finite trivial
coefficients, continuous cohomology satisfies
\[
 H^i(P,\mathbf F_\ell)
       =\varinjlim_{P\subset U\subset\Pi\ \mathrm{open}}
                          H^i(U,\mathbf F_\ell).          \tag{A52}
\]
Here is the continuity argument needed for this
specific use. A finite-valued continuous cochain
on the closed profinite subset $P^i$ extends
to a continuous cochain on $\Pi^i$, by extending
its finite clopen partition. Its cocycle
equations hold on $P^{i+1}$; compactness and
local constancy make them hold on $U^{i+1}$
for an open subgroup $U=P V$, with $V$ open
normal in $\Pi$. The same extension and
shrinking argument for a cochain whose
coboundary is prescribed proves injectivity
of the direct-limit comparison. This proves
(A52), for both $i=1$ and $i=2$.

For such an open $U$, let $g_U$ be the genus
of its actual finite covering curve. The
finite-coefficient group-to-curve comparison
can be checked directly here. In the tower of
all finite \'etale covers, every $H^1$ class
is killed by its finite torsor, and every
$H^2$ class is killed by a further cover of
degree $\ell$. Such a cover exists from
the nonzero abelianization of every proper
hyperbolic covering curve. Higher curve
cohomology vanishes. The universal-cover
Cartan--Leray comparison therefore identifies
the continuous group cohomology with the
curve cohomology, compatibly with pullback.
The
orientation class and the covering trace give
\[
 \dim H^1(U,\mathbf F_\ell)=2g_U,
 \quad H^2(U,\mathbf F_\ell)=\mathbf F_\ell,
 \quad g_U-1=[\Pi:U](g-1).                                \tag{A53}
\]
For $V\subset U$, restriction on $H^2$ is
multiplication by $[U:V]$ under these orientation
identifications. These are the actual
cohomological maps of the finite \'etale
coverings, or equivalently of their complex
surface coverings.

Suppose first that the integers
$v_\ell([\Pi:U])$ for open $U\supset P$
are bounded. Choose $U_0$ attaining their
maximum. Every smaller open $V\subset U_0$
containing $P$ has index prime to $\ell$ in
$U_0$, and the same holds for every further
transition. Each restriction on $H^1$ is
injective: a nonzero homomorphism
$V\to\mathbf F_\ell$ has kernel of index
$\ell$, so cannot vanish on a subgroup of
prime-to-$\ell$ index.

The subgroup $P$ cannot be open in $\Pi$,
because every open surface group has
abelianization $\widehat{\mathbf Z}^{\,2g_U}$
and thus has nontrivial quotients at primes
different from $\ell$. Therefore the indices
$[\Pi:V]$ in this cofinal family are unbounded.
By (A53), their $H^1$ dimensions are unbounded,
and all these spaces inject into the direct
limit in (A52). This contradicts finite
generation of the pro-$\ell$ group $P$.

Thus the $\ell$-parts of the indices are
unbounded. Given any open $U\supset P$, choose
a further open $V\subset U$ containing $P$
with $\ell\mid[U:V]$; intersecting $U$ with
an overgroup of sufficiently large
$\ell$-index gives such a $V$.
Its restriction map on $H^2$ is zero by
(A53). Formula (A52) consequently gives
$H^2(P,\mathbf F_\ell)=0$, which is the
pro-$\ell$ freeness criterion.

For the last assertion, let $C\ne1$ be the
procyclic normal subgroup. For each element
$h$ the subgroup generated by $C$ and $h$
is solvable and free pro-$\ell$, hence
procyclic. Thus every $h$ centralizes $C$.
A noncyclic free pro-$\ell$ group has trivial
center, so the whole finitely generated group
is procyclic.
\end{proof}

\begin{proposition}[Proved: proper-curve normalizers obey the same character and splitting restrictions]
\label{A:proper-normalizer-restrictions}
Let $X/K$ be proper hyperbolic over a number
field, and let $A\ne1$ be a closed procyclic
subgroup of its full geometric group whose
arithmetic normalizer has open image $G_L$.
Then the conclusions about arithmetic character
descent, excluded primes, and CM fields in
Lemma~\ref{A:normalizer-character-descent} and
Propositions~\ref{A:full-support-normalizer}--\ref{A:cm-normalizer-field}
hold, with the boundary alternative absent.
Explicitly, $L$ contains a CM subfield, the
prime support of $A$ omits the positive-density
complex-conjugation Frobenius set described
there, and each supported $\ell$ has the
stated CM-field embedding into $\mathbf Q_\ell$.
In particular no such $A$ has cofinite prime
support, and there is none when $L$ contains
no CM subfield.

Its normalizer has no section after any finite
arithmetic base extension, using the supplied
local theorem as in Proposition~\ref{A:split-normalizer}.
For every supported $\ell$ it also has the
exact open-subgroup structure (A44), with
the same residual pro-prime-to-$\ell$
extension and no virtual section.
These conclusions do not exclude all of
the remaining CM and restricted-support cases.
\end{proposition}
\begin{proof}
We specify the replacements of the affine
inputs, so the assertions compare the same
full geometric groups and actual arithmetic
maps. For a pro-$\ell$ Sylow subgroup $S$
of the geometric normalizer and any $h\in S$,
the subgroup $\langle A_\ell,h\rangle$ is a
finitely generated pro-$\ell$ subgroup of the
full surface group. It is free by
Lemma~\ref{A:surface-proprime-free}, and
solvable because its quotient by $A_\ell$
is procyclic. It is therefore cyclic.
This proves precisely the triviality of the
pro-$\ell$ part of the geometric normalizer
character used in
Lemma~\ref{A:normalizer-character-descent}.
The remaining image is finite and prime to
$\ell$, and the same power descends to the
open arithmetic quotient.

Choose the same detector $H=V A_\ell$, with
$V$ characteristic open in the surface group.
Its abelianization is again torsion-free,
now $\widehat{\mathbf Z}^{\,2g_H}$, and the
cyclic quotient $H/V$ gives the identical
injective homology line. The cover is proper,
so this homology is already the Tate module
of its Jacobian. Thus its weights, its
Hodge--Tate weights, and the descended power
character are exactly those in the compact
case of Lemma~\ref{A:normalizer-character-descent}.
The Serre algebraic-type proof in
Propositions~\ref{A:conjugation-prime-character}--\ref{A:cm-normalizer-field}
then applies to these actual characters.
There is no toric summand and no cusp case.
This proves every stated field and prime-support
restriction.

An actual arithmetic section in the normalizer
would give, by the same detector and Faltings
retraction, the explicit arithmetic map
(A43) from a Jacobian with nonzero procyclic
geometric image. Proposition~\ref{F-local-dichotomy}
has no cusp alternative for the proper
target $X$, and such an image is not open.
It therefore forces the image to be trivial,
a contradiction. This proves failure of
every virtual normalizer section.

Finally the Sylow step of
Proposition~\ref{A:normalizer-prime-complement}
also survives. After the geometric character
has been killed, every finitely generated
subgroup of $S$ containing a generator of
the central $A_\ell$ is free pro-$\ell$ with
nontrivial center, hence cyclic. Every finite
quotient of $S$ is consequently cyclic,
so $S$ itself is procyclic. The exact open
normal shrinking, central Sylow complement,
and intersection formula (A45) now apply.
The detector is proper and its compact line
is nonzero automatically. Its Abel--Jacobi
map and the integral Faltings retraction
give (A47)--(A48), and the same finite
cocycle gives the quotient section and
(A44). The non-splitting of the remaining
extension follows from the preceding
proper-target contradiction.
\end{proof}

\subsection{Deck-twisted fibre isolation and the fourth-to-third derived projection}

\begin{proposition}[Proved: one derived layer forces preservation of every element conjugacy class]\label{A:one-layer-class-preservation}
Let $K$ be a number field, let
$U=\mathbf P^1_K\setminus\{0,1,\infty\}$, and put
$S_j=F/F^{(j)}$ for its geometric derived quotients.
Write $C_j^\partial(K)$ for the centralizer of the
arithmetic outer action in the outer automorphisms of
$S_j$ preserving the three peripheral conjugacy classes
as an unordered set. For every $m\geq2$ and
$\gamma\in C_{m+1}^\partial(K)$, there is a unique
geometric cusp permutation $\sigma\in S_3$ such that
the projection of $\sigma^{-1}\gamma$ to
$\operatorname{Out}(S_m)$ has a representative
$\alpha$ satisfying
\[
 \text{for every }z\in S_m,\qquad
             \alpha(z)\text{ is conjugate to }z
             \text{ in }S_m.                         \tag{A66.1}
\]
In particular $\alpha$ preserves every closed normal
subgroup and every conjugacy class of closed procyclic
subgroups. The assertion in (A66.1) allows its conjugator
to depend on $z$.
\end{proposition}
\begin{proof}
Proposition~\ref{A:third-derived-projection} determines
the metabelian image of $\gamma$ as a unique cusp
permutation: uniqueness follows from the faithful
action of $S_3$ on $F^{\mathrm{ab}}$.
Remove that permutation and conjugate a representative
$\phi$ on $S_{m+1}$ geometrically so that its induced
automorphism on $M=S_2$ is literally the identity.
The center-free arithmetic-extension construction of
Proposition~\ref{A:solvable-centers} gives an actual
automorphism $\beta$ of the arithmetic extension
$E_{m+1}$ over $G_K$. Use the tangential section at
$0$ to write
\[
 \beta(s(g))=d_gs(g),\qquad
 \phi\rho_g\phi^{-1}=\operatorname{Int}(d_g)\rho_g.
\]
The projection of $d_g$ to $M$ is central, because
$\phi_M=1$. Since $Z(M)=1$, it is trivial.
Consequently the induced arithmetic automorphism
on $E_2$, and hence on $E_1$, is exactly the
identity. This keeps the constants in every Kummer
cover fixed.

Choose a characteristic open normal subgroup of
$S_m$ and denote its quotient by $G$.
Such subgroups are cofinal because $S_m$ is finitely
generated. The automorphisms induced by $\phi$ and
the actual tangential arithmetic action both act
on $G$. After a finite extension $K_0/K$, that
arithmetic action is trivial, so the corresponding
arithmetic quotient is
$G\times G_{K_0}$. The induced $\beta$ has the form
\[
 (g,\tau)\longmapsto
             \bigl(\alpha_G(g)z(\tau),\tau\bigr),
       \qquad z:G_{K_0}\longrightarrow Z(G).
                                                        \tag{A66.2}
\]
Indeed $(1,\tau)$ centralizes $G$, so its image
does also; the homomorphism law makes $z$ a
continuous homomorphism into the finite center.
A further finite extension kills $z$.
Thus (A66.2) becomes $(\alpha_G,\mathrm{id})$.
No center-free assumption on the finite group
$G$ has been made.

Let $C\to U_{K_0}$ be the resulting geometrically
connected finite etale $G$-torsor, with its constant
deck group. Fix an element $g\in G$.
Hilbert irreducibility supplies $a\in U(K_0)$
whose fibre is a connected $G$-torsor.
Its splitting field is a $G$-extension of $K_0$.
Apply Chebotarev in a normal closure over
$\mathbb Q$, extending the chosen $g$ while
fixing $K_0$. This supplies a place of $K_0$
of residue degree one over a rational prime $q$,
at which the fibre Frobenius is conjugate to $g$.
Exclude the finitely many bad places for the
marked cover, its deck maps, and $a$, and those
dividing $|G|$. These are exactly the Hilbert and
degree-one Chebotarev steps used in
Proposition~\ref{A:third-derived-IA}.

Put $N=q-1$ and enlarge the constants to
$K_1=K_0(\mu_N)$. The chosen place has an
extension still with residue field $\mathbb F_q$,
since all $N$th roots of unity lie there.
Over this field consider
\[
 B:\qquad v_1^N=t/a,\qquad
             v_2^N=(1-t)/(1-a),\qquad
                   Z=C\times_U B.                    \tag{A66.3}
\]
These are finite etale covers of the same tripod.
They have smooth marked reduction at the chosen
place: the original marked cover has good
reduction there, and the newly added ramification
indices divide $q-1$. Normalization, with the
punctures removed, gives the indicated models.
The arithmetic automorphism $\beta$ preserves
the two equations of $B$, including their
constant Kummer twists, because it induces the
identity on $E_1$.
It acts on the finite fibre of $Z$ by applying
$\alpha_G$ to its $G$-coordinate and leaving the
$B$-coordinate unchanged.

We give the cohomological comparison for every
component, without assuming that $Z$ is
geometrically connected. Its geometric monodromy
is a subgroup of $G\times(\mathbb Z/N)^2$ and
has derived length at most $m$.
The subgroup $H\subset F$ corresponding to any
component therefore contains $F^{(m)}$.
It follows that
\[
 F^{(m+1)}=[F^{(m)},F^{(m)}]\subset H',
 \qquad
 (H/F^{(m+1)})^{\mathrm{ab}}=H^{\mathrm{ab}}.
                                                        \tag{A66.4}
\]
Thus $E_{m+1}$ retains the full first homology
of each component, together with the maps
between components induced by arithmetic and
deck transformations. More explicitly, components
are the orbits of $F$ on the finite fibre;
$H$ is the stabilizer of a point of an orbit,
and changing that point conjugates $H$.
The automorphism $\beta$ maps the orbit and
stabilizer diagrams to themselves. The canonical
identification in (A66.4) commutes with these
conjugation maps. It therefore gives an
isomorphism on the direct sum of the first
cohomology groups which commutes with arithmetic
and intertwines a deck transformation $h\in G$
with $\alpha_G(h)$.
The same finite orbit diagram identifies
$H^0$. Poincare duality on each smooth component,
with its canonical orientation, then identifies
$H_c^1$ and $H_c^2$ equivariantly as well.
For any auxiliary coefficient prime different
from $q$, good marked reduction transfers these
identifications to the special fibre.

Let $\mathcal F$ denote its Frobenius operator
in the fixed-point convention. The preceding
comparison gives equality of alternating traces
on compactly supported cohomology for
$h\mathcal F$ and $\alpha_G(h)\mathcal F$.
These alternating traces count their fixed
geometric points. To apply the usual
Lefschetz trace formula here, one may twist the
finite-field descent datum by the finite-order
deck map $h$: it is defined over $\mathbb F_q$
and commutes with Frobenius, so the twisted
descent is effective and its rational points
are precisely the fixed points of $h\mathcal F$.
This also handles the possible permutation of
components.

All such fixed points lie above $t\in U(\mathbb F_q)$.
The deck map $h$ is trivial on $B$.
Since every nonzero element of $\mathbb F_q$
has $N$th power one, (A66.3) has rational points
over exactly $t=\overline a$, with $N^2$ points
in that fibre. Label the regular $G$-fibre
of $C$ there by $G$ so that Frobenius is right
multiplication by $g$ and the deck map $h$ is
left multiplication by $h$.
Its fixed-point condition is $hkg=k$.
Consequently
\[
 \#\operatorname{Fix}(h\mathcal F\mid Z)
  =N^2
    \begin{cases}
      |C_G(g)|,&h\text{ is conjugate to }g^{-1},\\
      0,&\text{otherwise}.
    \end{cases}                                      \tag{A66.5}
\]
Take $h=g^{-1}$ and compare it with
$\alpha_G(h)$. Equality of traces implies that
$\alpha_G(g)$ is conjugate to $g$.
The argument applies to every $g\in G$.

Finally fix $z\in S_m$. For each characteristic
finite quotient $G$, the set of $u\in S_m$ such
that $u z u^{-1}$ and $\alpha(z)$ have the same
image in $G$ is nonempty and closed.
The sets have the finite intersection property
by passage to a common finer characteristic
quotient. Compactness gives one conjugator for
this $z$ in $S_m$ itself. This proves (A66.1).
Applying it to the elements of a closed normal
subgroup shows preservation of that subgroup;
applying it to a topological generator gives
preservation of each procyclic conjugacy class.
\end{proof}

\begin{proposition}[Proved: class-preserving rigidity at derived length three and the fourth-level projection]\label{A:third-derived-class-rigidity}
Every continuous automorphism of $S_2$ or $S_3$
which sends each element to a conjugate of itself
is inner. Consequently, for every number field
$K$,
\[
 \operatorname{im}\bigl(
 C_4^\partial(K)\longrightarrow
                  \operatorname{Out}(S_3)\bigr)=S_3,
                                                        \tag{A67.1}
\]
where the group on the right is the geometric
group of cusp permutations. The same equality
holds with $C_j^\partial(K)$ for any $j\geq4$.
This is an equality of projected commutants;
the kernel in the same derived level is not
computed by it.
\end{proposition}
\begin{proof}
We first prove the completed-module facts on
which the rigidity argument depends.
For a profinite group $P$ write
$\Lambda_P=\widehat{\mathbb Z}[[P]]$.
If $u\in P$ has infinite $p$-primary order,
then both left and right multiplication by
$u-1$ have zero kernel on $\mathbb F_p[[P]]$.
For example, an element killed by left
multiplication is invariant under $\langle u\rangle$.
At a finite quotient it has constant
coefficients on the $u$-orbits.
Refine that quotient so that each such orbit
has its length multiplied by a factor divisible
by $p$. Projection of each invariant orbit sum
is then zero in characteristic $p$.
Hence its image at every original finite
quotient is zero. The right-hand assertion
uses the right orbits in the same way.
The required refinements exist because the
closed cyclic subgroup has its full induced
profinite topology and infinite $p$-part.
Successive division by $p$ gives the same
injectivity on $\mathbb Z_p[[P]]$.
If $u$ has infinite $p$-part for every $p$,
the assertion holds on $\Lambda_P$.

Suppose two commuting elements $u,v$ generate
a closed subgroup isomorphic to
$\widehat{\mathbb Z}^2$.
The quotient
$\Lambda_P/(u-1)\Lambda_P$ is the completed
free coefficient module on the left coset
space $\langle u\rangle\backslash P$.
One sees this at every finite quotient;
compactness identifies the kernel with
the closed image of multiplication by $u-1$.
Left multiplication by $v$ is well defined
on this space and has free procyclic orbits.
In finite quotient approximations its orbit
lengths can be refined by arbitrarily large
powers of any prime: the induced topology
on $\langle u,v\rangle/\langle u\rangle
\simeq\widehat{\mathbb Z}$ supplies these
refinements. The preceding orbit-sum proof
therefore also proves
\[
 v-1:\Lambda_P/(u-1)\Lambda_P
       \longrightarrow\Lambda_P/(u-1)\Lambda_P
             \quad\hbox{is injective}.                \tag{A67.2}
\]

Here is a degree-one cohomology consequence
in the exact topology needed below.
Suppose that $P$ has a normal abelian subgroup
$D$ containing such $u,v$.
For a continuous cocycle
$f:P\to\Lambda_P$, put $a=f(u)$ and $b=f(v)$.
The commuting relation gives
$(v-1)a=(u-1)b$.
By (A67.2), $a=(u-1)t$ for some
$t\in\Lambda_P$. The injectivity of $u-1$
then gives $b=(v-1)t$.
Subtract this coboundary. The new cocycle
vanishes on $\langle u,v\rangle$.
Commuting with $u$ forces its value at every
$d\in D$ to be $u$-invariant, and hence zero.
Normality of $D$ then forces its value at
every $g\in P$ to be $D$-invariant, again zero.
This proves, for continuous cocycles with
these compact coefficients,
\[
                  H^1(P,\Lambda_P)=0.                 \tag{A67.3}
\]
In particular (A67.3) holds for
$P=\widehat{\mathbb Z}^2$, with $u,v$ its
basis.

For $M=S_2$, let
$c=[x,y]$ and $b=xcx^{-1}$.
Its derived module is the free rank-one
module
$M'=\widehat{\mathbb Z}[[M^{\mathrm{ab}}]]c$.
The vectors $c,b$ correspond to $1,X$.
They generate a copy of
$\widehat{\mathbb Z}^2$: independence of
$1,X$ already follows after setting
$X^2=1$ in the completed group algebra.
Thus $D=M'$ supplies the normal abelian
subgroup needed for (A67.3).

Let first an automorphism of $S_2$ be
class-preserving. It induces the identity on
the abelianization. Its action on the free
rank-one derived module sends $c$ to $a c$
for some $a\in M^{\mathrm{ab}}$, since the
image of $c$ is conjugate to $c$.
Remove an inner automorphism lifting $a$.
The resulting automorphism is the identity
on both the abelian quotient and the derived
module. Its difference from the identity
is a continuous cocycle from
$\widehat{\mathbb Z}^2$ to that regular
module. By (A67.3), it is a coboundary, so
the automorphism is inner.

Now put
\[
 S=S_3,\qquad M=S_2,\qquad
 L=F''/F''',\qquad
 \Lambda=\widehat{\mathbb Z}[[M]].
\]
The Fox--Magnus exact sequence used in
Proposition~\ref{A:last-layer-unrestricted}
identifies $L$ with
\[
 \ker\bigl(\partial:\Lambda^2\to\Lambda\bigr),
 \qquad
       \partial(a_1,a_2)=a_1(x-1)+a_2(y-1).
                                                        \tag{A67.4}
\]
Coefficients act on the left. The homomorphism
$\partial$ is left $M$-equivariant.
An $L$-valued cocycle, considered in
$\Lambda^2$, is a coboundary by (A67.3).
If its bounding vector is $w$, then
$(g-1)\partial w=0$ for all $g\in M$.
The invariants of $\Lambda$ are zero, by
injectivity of $c-1$. Therefore $\partial w=0$,
so $w\in L$. We have proved
\[
                         H^1(M,L)=0.                 \tag{A67.5}
\]

We next construct a specific element of $L$
with no left annihilator over $\Lambda$,
including after reduction at every coefficient
prime. Regard
$r=[c,xcx^{-1}]$ as an element of
$F''/F'''$.
The left Fox rule
$\delta(uv)=\delta(u)+\overline u\,\delta(v)$
gives
\[
 \delta_y(c)=x-c,\qquad
 \delta_y(r)=A(x-c),\qquad
 A=(c-1)x-(b-1).
                                                        \tag{A67.6}
\]
The second identity uses that $c,b$ commute
in $M$. For $d=c-1$ we have
$d^x=b-1=xdx^{-1}$, and a literal multiplication
gives
\[
                   A d=d\,d^x(x-1).                   \tag{A67.7}
\]
All of $c,b,x$ and $xc^{-1}$ have infinite
$p$-part for every prime $p$:
this follows from the free derived-module
description for $c,b$ and from the
abelianization for $x,xc^{-1}$.
Hence each factor on the right of (A67.7)
has zero left annihilator, even over
$\mathbb F_p[[M]]$.
If $zA=0$, multiply on the right by $d$
and cancel these injective right
multiplications successively; this gives
$z=0$. Since $x-c=(xc^{-1}-1)c$ also has
zero left annihilator, (A67.6) proves that
the second coordinate of $r$ has zero
left annihilator. Thus
$\Lambda\to L$, $z\mapsto zr$, is injective.
Only this specified sided regularity is used;
the completed group algebra is not declared
to be a domain.

Let $\alpha$ be a class-preserving automorphism
of $S$. Its induced automorphism on $M$ is
inner by the metabelian argument already proved.
After removing that inner automorphism, assume
$\alpha_M=1$.
Its action $T$ on $L$ is then $\Lambda$-linear.
For any $v\in L$ generating a free
$\Lambda$-submodule, class preservation says
$T(v)=g v$ for some $g\in M$.
Apply class preservation also to
$(1+2h)v$ for arbitrary $h\in M$.
Freeness of the cyclic module generated by
$v$ gives, for some $g_h\in M$,
\[
                    g+2hg=g_h+2g_hh.                  \tag{A67.8}
\]
Distinct elements of $M$ are linearly
independent in $\widehat{\mathbb Z}[[M]]$,
as finite quotients separate any finite
list of them. When $h\ne1$, the two terms
on each side of (A67.8) have distinct
coefficient vectors $1$ and $2$.
It follows that $g_h=g$ and $hg=gh$.
This also holds for $h=1$.
Thus $g\in Z(M)=1$, and $T(v)=v$.
This coefficient comparison takes place
over $\widehat{\mathbb Z}$, so it includes
the coefficient prime two.

To apply this conclusion to every element
of $L$, choose the central scalar
$\nu\in\widehat{\mathbb Z}$ whose
$p$-coordinate is $p$ for every prime $p$.
For any $l\in L$, the second coordinate
of $r+\nu l$, in the coefficient-$p$
completed group algebra, reduces modulo
$p$ to the second coordinate of $r$.
The latter has zero left annihilator by
(A67.6)--(A67.7).
Successive division by $p$ shows that
the second coordinate of $r+\nu l$ has
zero left annihilator over
$\mathbb Z_p[[M]]$.
Taking the product over all $p$ proves
that $r+\nu l$ generates a free
$\Lambda$-submodule.
The preceding paragraph fixes both $r$
and $r+\nu l$, so $\nu(T(l)-l)=0$.
Multiplication by $\nu$ is injective
on $\Lambda^2$, and hence on $L$:
each coefficient-$p$ factor is
$p$-torsion-free.
Therefore $T(l)=l$ for all $l\in L$.

An automorphism acting identically on both
$M$ and $L$ now gives the continuous cocycle
\[
 M\longrightarrow L,\qquad
   \overline s\longmapsto\alpha(s)s^{-1}.
                                                        \tag{A67.9}
\]
The expression is independent of the
lift $s$, since $L$ is abelian and
$\alpha|L=1$. It is continuous by
the quotient topology.
Equation (A67.5) makes (A67.9) a coboundary,
so $\alpha$ is inner by an element of $L$.
This proves the class-preserving rigidity
for $S_3$.

Finally apply
Proposition~\ref{A:one-layer-class-preservation}
with $m=3$. After the unique cusp permutation
is removed, the projected automorphism is
class-preserving and hence inner.
Every cusp permutation is geometric and
commutes with arithmetic, which supplies
the reverse inclusion in (A67.1).
For $j>4$, project first to level four.
\end{proof}
